% ==============================================================================
% CH11_QUANTIZATION_LORA.TEX — CHƯƠNG 11: LƯỢNG TỬ HÓA, TINH CHỈNH LORA & CĂN CHỈNH SỞ THÍCH
% ==============================================================================
% Định dạng: Sách chuyên khảo kinh điển (Classic Engineering Monograph)
% Chuẩn mực: Không dùng hộp đóng khung 4 phía, toán học 5 bước, bài tập đa tầng
% ==============================================================================

\chapter{Lượng Tử Hóa, Tinh Chỉnh LoRA \& Căn Chỉnh Sở Thích (Alignment)}
\label{ch:quantization-lora}

Một mô hình nền tảng hiện đại như Llama-3-70B sở hữu năng lực suy luận phi thường, nhưng mức tiêu thụ tài nguyên phần cứng của nó đặt ra rào cản chi phí khổng lồ cho các ứng dụng thực tế. Lưu trữ 70 tỷ trọng số ở định dạng chuẩn nửa độ chính xác (FP16/BF16) đòi hỏi tối thiểu $140\text{ GB}$ VRAM chỉ riêng cho tham số tĩnh, chưa kể bộ nhớ đệm kích hoạt (Activation Memory) và bộ nhớ đệm ngữ cảnh KV Cache. Trong giai đoạn huấn luyện hoặc tinh chỉnh toàn phần (Full Fine-Tuning), các trạng thái bổ trợ của thuật toán tối ưu AdamW cùng với gradients có thể đẩy yêu cầu bộ nhớ vượt ngưỡng $840\text{ GB}$ VRAM, buộc các doanh nghiệp phải đầu tư các cụm GPU đắt đỏ gồm nhiều máy chủ HGX H100.

Để vượt qua bức tường tài nguyên này, Kỹ sư AI hiện đại bắt buộc phải làm chủ ba trụ cột kỹ nghệ nén, thích ứng và điều hướng hành vi mô hình:
\begin{enumerate}
  \item \textbf{Lượng tử hóa sau huấn luyện (Post-Training Quantization -- PTQ)}: Nén biểu diễn số học của tham số và kích hoạt từ số thực dấu phẩy động 16-bit xuống số nguyên 8-bit hoặc 4-bit (INT8, INT4, NF4, GGUF), giúp giảm $50\% - 75\%$ dung lượng VRAM và tăng tốc độ suy luận mà hầu như không làm tổn hại năng lực suy luận.
  \item \textbf{Tinh chỉnh thích ứng hiệu quả tham số (PEFT / LoRA / QLoRA)}: Đóng băng toàn bộ trọng số gốc và chỉ cập nhật một cặp ma trận bổ trợ có hạng số chiều nội tại siêu nhỏ (Low Intrinsic Dimension), cho phép tinh chỉnh các mô hình khổng lồ trên phần cứng tiêu dùng phổ thông.
  \item \textbf{Căn chỉnh sở thích con người (Human Preference Alignment)}: Điều chỉnh hành vi đầu ra của mô hình sau khi tinh chỉnh lệnh (SFT) để đảm bảo tính hữu ích, trung thực và an toàn thông qua các phương pháp tiên tiến: từ học tăng cường từ phản hồi con người (RLHF với PPO) đến tối ưu hóa sở thích trực tiếp (DPO, ORPO, SimPO).
\end{enumerate}

---

\section{Nền Tảng Số Học Máy Tính: Từ IEEE 754 Đến FP8 \& INT8}
\label{sec:computer-arithmetic-foundations}

Để hiểu bản chất của lượng tử hóa, trước hết người kỹ sư phải nắm vững cách máy tính số biểu diễn các đại lượng vô hạn của toán học liên tục trong một số lượng hữu hạn các bit nhị phân.

\subsection{Tiêu Chuẩn Biểu Diễn Số Thực Dấu Phẩy Động IEEE 754}
Theo tiêu chuẩn IEEE 754, một số thực dấu phẩy động $x \in \mathbb{R}$ được mã hóa nhị phân qua ba trường thông tin:
\begin{equation}
x = (-1)^s \times 2^{e - \text{bias}} \times \left(1 + \sum_{i=1}^{m_{\text{bits}}} b_i 2^{-i}\right)
\end{equation}
Trong đó:
\begin{itemize}
  \item $s \in \{0, 1\}$: Bit dấu (Sign bit, xác định số âm hay dương).
  \item $e$: Trường số mũ (Exponent, xác định dải động của đại lượng). Để biểu diễn cả số mũ âm và dương, một giá trị bù cố định $\text{bias} = 2^{e_{\text{bits}} - 1} - 1$ được trừ đi từ giá trị nguyên không dấu của trường này.
  \item $m$: Trường định trị hoặc phân số (Mantissa / Fraction, xác định độ chính xác và độ phân giải giữa hai số liên tiếp). Số $1$ phía trước là bit ẩn (implicit leading bit) của số chuẩn hóa.
\end{itemize}

\begin{table}[ht]
\centering
\footnotesize
\setlength{\tabcolsep}{3.5pt}
\caption{So sánh cấu trúc nhị phân và đặc tính số học của các định dạng dấu phẩy động.}
\label{tab:floating-point-formats}
\begin{tabular}{lcccccc}
\toprule
\textbf{Định dạng} & \textbf{Bit} & \textbf{Dấu ($s$)} & \textbf{Mũ ($e$)} & \textbf{Định trị ($m$)} & \textbf{Dải động xấp xỉ} & \textbf{Độ phân giải} \\
\midrule
\textbf{FP32} (Đơn)   & 32 & 1 & 8  & 23 & $1.18 \times 10^{-38} \to 3.40 \times 10^{38}$ & Rất cao ($10^{-7}$) \\
\textbf{FP16} (Nửa)   & 16 & 1 & 5  & 10 & $6.10 \times 10^{-5} \to 6.55 \times 10^{4}$  & Trung bình ($10^{-3}$) \\
\textbf{BF16} (Bfloat16)& 16 & 1 & 8  & 7  & $1.18 \times 10^{-38} \to 3.40 \times 10^{38}$ & Thấp ($10^{-2}$) \\
\textbf{FP8 E4M3}     & 8  & 1 & 4  & 3  & $1.95 \times 10^{-3} \to 448$                  & Thấp \\
\textbf{FP8 E5M2}     & 8  & 1 & 5  & 2  & $1.52 \times 10^{-5} \to 57344$                & Rất thấp \\
\bottomrule
\end{tabular}
\end{table}

\paragraph{Tại sao BF16 (Brain Floating Point) là cứu tinh của huấn luyện mô hình lớn?}
Trong những năm đầu của học sâu, việc huấn luyện bằng FP16 gặp phải vấn đề nghiêm trọng: dải số mũ của FP16 chỉ có 5 bits, do đó giá trị dương nhỏ nhất có thể biểu diễn được là xấp xỉ $6.1 \times 10^{-5}$. Trong các mô hình ngôn ngữ khổng lồ với hàng chục tỷ tham số, gradients của các tầng đầu thường xuyên nhỏ hơn $10^{-5}$, dẫn đến hiện tượng tràn số dưới (Underflow) và làm các đạo hàm triệt tiêu về $0.0$. Để khắc phục, kỹ sư phải sử dụng kỹ thuật co giãn hàm mất mát (Loss Scaling) rất phức tạp và dễ bất ổn định.

Định dạng \textbf{BF16} do Google Brain đề xuất giải quyết triệt để vấn đề này bằng cách giữ nguyên 8 bit số mũ giống hệt FP32 và chỉ cắt giảm trường định trị xuống 7 bits. Nhờ đó, BF16 có dải động tương đương FP32 ($10^{\pm 38}$), giúp quá trình lan truyền ngược hoàn toàn miễn nhiễm với lỗi Underflow mà không cần bất kỳ cơ chế Loss Scaling nào!

\paragraph{Định dạng FP8 thế hệ mới (E4M3 và E5M2) \& Kiến trúc FP8 Serving}:
Trên các kiến trúc GPU tiên tiến (NVIDIA Hopper H100, Blackwell B200), định dạng FP8 (8-bit Floating Point) được chuẩn hóa theo đặc tả OpenFP8 với hai biến thể chuyên biệt:
\begin{itemize}
  \item \textbf{FP8 E4M3 (1 bit dấu, 4 bits số mũ, 3 bits định trị)}:
  Sở hữu độ phân giải cao hơn với $\text{bias} = 7$. Giá trị thực dương lớn nhất có thể biểu diễn được là:
  \begin{equation}
  x_{\max}^{\text{E4M3}} = 2^{15 - 7} \times \left(1 + \frac{6}{8}\right) = 2^8 \times 1.75 = 448
  \end{equation}
  Đặc điểm kỹ thuật độc đáo của E4M3 là \textbf{không có mã biểu diễn vô cực ($\pm \infty$)}, và chỉ dành duy nhất một cấu hình bit cho giá trị bất định ($\text{NaN} = S.1111.111_2$). Việc loại bỏ vô cực giúp E4M3 tận dụng tối đa dải bit khả dụng để tăng gấp đôi độ phân giải số học, biến nó thành lựa chọn hoàn hảo cho trọng số và kích hoạt trong lượt suy luận xuôi (\term{Inference Forward Pass}) và serving thực tế.
  \item \textbf{FP8 E5M2 (1 bit dấu, 5 bits số mũ, 2 bits định trị)}:
  Sở hữu dải động khổng lồ tương đương FP16 với $\text{bias} = 15$. Giá trị lớn nhất đạt tới:
  \begin{equation}
  x_{\max}^{\text{E5M2}} = 2^{30 - 15} \times \left(1 + \frac{3}{4}\right) = 2^{15} \times 1.75 = 57344
  \end{equation}
  E5M2 tuân thủ cấu trúc IEEE 754 với các biểu diễn chuẩn cho cả $\pm \infty$ và $\text{NaN}$, đóng vai trò sống còn trong lượt lan truyền ngược (\term{Backward Pass}) để bảo vệ các gradients có biên độ biến thiên hàng nghìn lần không bị tràn số.
\end{itemize}

Trong các khung serving LLM hiện đại như vLLM, TensorRT-LLM và SGLang, việc triển khai \textbf{FP8 Serving} giúp tăng gấp đôi thông lượng giải mã (Tokens/sec) và giảm $50\%$ dung lượng KV Cache thông qua hai chiến lược điều chỉnh hệ số co giãn:
\begin{enumerate}
  \item \textbf{Co Giãn Động (Dynamic Scaling)}:
  Hệ số co giãn được tính toán thời gian thực ngay khi tensor kích hoạt $\mathbf{X}$ vừa được sinh ra:
  \begin{equation}
  S_x = \frac{\max(|\mathbf{X}|)}{448.0}, \qquad \mathbf{X}_{\text{FP8}} = \text{clip}\left(\frac{\mathbf{X}}{S_x}, -448.0, 448.0\right)
  \end{equation}
  Phương pháp này bảo toàn $99.9\%$ độ chính xác của FP16, nhưng đòi hỏi một kernel thu nhỏ CUDA (\term{Reduction Kernel}) để tìm giá trị cực đại trước mỗi phép tính GEMM, tạo ra một rào cản đồng bộ hóa phần cứng nhẹ.
  \item \textbf{Co Giãn Trễ / Tĩnh (Delayed / Static Scaling)}:
  Để triệt tiêu hoàn toàn rào cản đồng bộ hóa thời gian chạy, hệ thống ghi nhận lịch sử giá trị cực đại của kích hoạt qua $H$ bước trước đó (hoặc đo đạc trước trên tập dữ liệu hiệu chuẩn offline) để xác định một hệ số $S_{\text{delayed}}$ cố định:
  \begin{equation}
  S_{\text{delayed}} = \frac{1}{448.0} \max_{k=1..H} \left( \max(|\mathbf{X}_{t-k}|) \right)
  \end{equation}
  Nhờ đó, phép chuyển đổi sang FP8 và phép nhân ma trận FP8 GEMM trên Tensor Cores được thực thi liên tục và chồng lấn hoàn toàn (\term{Kernel Fusion \& Overlapping}), đạt thông lượng phần cứng tối đa.
\end{enumerate}

---

\subsection{Toán Học Lượng Tử Hóa Tuyến Tính Số Nguyên (Integer Quantization)}

Lượng tử hóa tuyến tính (Uniform Quantization) thiết lập một hàm ánh xạ từ miền giá trị liên tục của số thực $x \in [\alpha, \beta] \subset \mathbb{R}$ sang một tập các số nguyên có bước nhảy đều nhau $q \in [q_{\min}, q_{\max}] \subset \mathbb{Z}$.

\begin{definition}[Lượng Tử Hóa Tuyến Tính Bất Đối Xứng -- Asymmetric Quantization]
Cho tensor số thực $\mathbf{x}$. Hai tham số lượng tử hóa được xác định:
\begin{itemize}
  \item \textbf{Hệ số co giãn (Scale factor $S$)}:
  \begin{equation}
  S = \frac{\max(\mathbf{x}) - \min(\mathbf{x})}{q_{\max} - q_{\min}}
  \end{equation}
  \item \textbf{Điểm không (Zero-point $Z$)}: Đại lượng số nguyên tương ứng chính xác với giá trị số thực $0.0$, bảo đảm giá trị 0 không bị méo mó sau khi nén:
  \begin{equation}
  Z = \text{round}\left( -\frac{\min(\mathbf{x})}{S} \right) + q_{\min}
  \end{equation}
\end{itemize}
Phép biến đổi xuôi (Quantization) và biến đổi ngược (Dequantization):
\begin{equation}
q = \text{clip}\left( \left\lfloor \frac{x}{S} \right\rceil + Z, q_{\min}, q_{\max} \right)
\end{equation}
\begin{equation}
\hat{x} = S \cdot (q - Z)
\end{equation}
trong đó $\lfloor \cdot \rceil$ là phép làm tròn đến số nguyên gần nhất.
\end{definition}

\begin{definition}[Lượng Tử Hóa Tuyến Tính Đối Xứng -- Symmetric Quantization]
Khi miền giá trị được giả định đối xứng qua gốc tọa độ $(\alpha = -\beta = -\max(|\mathbf{x}|))$, điểm không $Z$ triệt tiêu hoàn toàn về $0$:
\begin{equation}
S = \frac{\max(|\mathbf{x}|)}{q_{\max}}, \qquad Z = 0
\end{equation}
\begin{equation}
q = \text{clip}\left( \left\lfloor \frac{x}{S} \right\rceil, -q_{\max}, q_{\max} \right), \qquad \hat{x} = S \cdot q
\end{equation}
\end{definition}

\paragraph{Ý nghĩa phần cứng then chốt của Lượng tử hóa đối xứng}:
Xét phép nhân ma trận trọng số và véc-tơ kích hoạt: $\mathbf{y} = \mathbf{x} \mathbf{W}$. Khi giải lượng tử hóa:
\begin{equation}
\mathbf{y} \approx \hat{\mathbf{x}} \hat{\mathbf{W}} = (S_x (\mathbf{q}_x - Z_x)) (S_w (\mathbf{q}_w - Z_w)) = S_x S_w \left( \mathbf{q}_x \mathbf{q}_w - Z_w \mathbf{q}_x - Z_x \mathbf{q}_w + Z_x Z_w \right)
\end{equation}
Nếu sử dụng lượng tử hóa bất đối xứng, mỗi phép nhân ma trận phải thực hiện thêm ba phép tính bù trừ số hạng chéo ($Z_w \mathbf{q}_x$, $Z_x \mathbf{q}_w$, $Z_x Z_w$), làm giảm hiệu suất phần cứng. Ngược lại, trong lượng tử hóa đối xứng ($Z_x = Z_w = 0$), công thức rút gọn thành:
\begin{equation}
\mathbf{y} \approx (S_x S_w) \cdot (\mathbf{q}_x \mathbf{q}_w)
\end{equation}
Toàn bộ phép tính nhân ma trận $\mathbf{q}_x \mathbf{q}_w$ là phép nhân ma trận số nguyên thuần túy, có thể đẩy trực tiếp vào các nhân phần cứng chuyên dụng (INT8 Tensor Cores hoặc chỉ lệnh DP4A trên GPU NVIDIA), đạt thông lượng tính toán (TFLOPS) cao gấp đôi so với phép nhân số thực FP16!

---

\subsection{Mức Độ Phân Khối Lượng Tử Hóa (Quantization Granularity)}
Sai số làm tròn phụ thuộc rất lớn vào phạm vi áp dụng của hệ số co giãn $S$:
\begin{itemize}
  \item \textbf{Per-tensor Quantization}: Dùng một hệ số $S$ duy nhất cho toàn bộ ma trận hàng triệu phần tử. Tốn ít bộ nhớ nhất (chỉ thêm một số float cho cả tầng), nhưng nếu có một vài phần tử cực đại cá biệt, toàn bộ các phần tử còn lại sẽ bị nén về $0$.
  \item \textbf{Per-channel / Per-column Quantization}: Mỗi hàng hoặc mỗi cột của ma trận trọng số sở hữu một hệ số $S$ riêng biệt. Phương pháp này bảo toàn đặc tính phân phối của từng nơ-ron độc lập, là chuẩn mực cho trọng số trong các thư viện hiện đại.
  \item \textbf{Group-wise / Block-wise Quantization}: Chia mỗi hàng của ma trận thành các nhóm liên tiếp (ví dụ: kích thước nhóm $G = 64$ hoặc $128$ phần tử), mỗi nhóm có một hệ số Scale riêng. Kỹ thuật này giảm thiểu triệt để ảnh hưởng của các giá trị ngoại lai cục bộ, là nền tảng của GPTQ và AWQ 4-bit.
\end{itemize}

---

\section{Giải Phẫu Các Thuật Toán Lượng Tử Hóa Hiện Đại (PTQ)}
\label{sec:ptq-modern-algorithms}

Trong phương pháp lượng tử hóa sau huấn luyện (Post-Training Quantization -- PTQ), mô hình đã hoàn thành huấn luyện được nén mà không cần chạy lại quá trình tiền huấn luyện tốn kém.

\subsection{Căn Bệnh Ngoại Lai Kích Hoạt (Activation Outliers) \& Giới Hạn 6.7B}
Năm 2022, trong công trình nghiên cứu nổi tiếng \textit{LLM.int8()}, Tim Dettmers và các cộng sự đã phát hiện một hiện tượng kỳ lạ xảy ra khi kích thước mô hình Transformer vượt qua ngưỡng \textbf{$6.7$ tỷ tham số (6.7B)}:
\begin{itemize}
  \item Trong các tầng sâu của Transformer, xuất hiện một số lượng rất nhỏ (khoảng $0.1\%$ các kênh kích hoạt) có giá trị biên độ đột biến, lớn gấp \textbf{$20$ đến $100$ lần} so với phần còn lại của tensor!
  \item Mặc dù số lượng cực ít, những kênh ngoại lai (Outliers) này nắm giữ toàn bộ thông tin phối hợp ngữ nghĩa và tư duy logic của mô hình.
  \item Nếu áp dụng lượng tử hóa INT8 trên toàn bộ tensor, hệ số Scale $S$ bị kéo lên rất lớn để bao phủ các giá trị ngoại lai này. Hệ quả là toàn bộ $99.9\%$ các giá trị kích hoạt bình thường bị làm tròn về $0$, phá hủy hoàn toàn độ chính xác của mô hình!
\end{itemize}

\subsection{INT8 SmoothQuant: Di Dời Độ Đột Biến Bằng Đại Số Tuyến Tính}
Được Xiao et al. (ICML 2023) đề xuất, \textbf{SmoothQuant} giải quyết triệt để bài toán hóc búa: \textit{Trọng số (Weights) rất dễ lượng tử hóa vì phân phối Gauss mượt mà không có ngoại lai; trong khi Kích hoạt (Activations) cực kỳ khó lượng tử hóa vì có các kênh ngoại lai khổng lồ.}

Ý tưởng toán học tinh tế của SmoothQuant là sử dụng một ma trận đường chéo $\mathbf{s} = \text{diag}(s_1, s_2, \dots, s_d)$ với $s_j > 0$ để chia tỷ lệ các kênh kích hoạt đầu vào và đồng thời nhân tỷ lệ bù trừ tương ứng vào các hàng của ma trận trọng số:
\begin{equation}
\mathbf{Y} = \mathbf{X} \mathbf{W} = \left(\mathbf{X} \cdot \text{diag}(\mathbf{s})^{-1}\right) \cdot \left(\text{diag}(\mathbf{s}) \cdot \mathbf{W}\right) = \hat{\mathbf{X}} \hat{\mathbf{W}}
\end{equation}
Trong đó phép nhân $\hat{\mathbf{X}} \hat{\mathbf{W}}$ hoàn toàn tương đương toán học $100\%$ với phép nhân ban đầu $\mathbf{X} \mathbf{W}$!

Hệ số di dời tối ưu cho kênh thứ $j$ được xác định theo công thức cân bằng:
\begin{equation}
s_j = \frac{\max(|\mathbf{X}_j|)^\alpha}{\max(|\mathbf{W}_j|)^{1-\alpha}}
\end{equation}
Trong đó $\max(|\mathbf{X}_j|)$ là giá trị kích hoạt lớn nhất trên kênh $j$, $\max(|\mathbf{W}_j|)$ là giá trị trọng số lớn nhất kết nối với kênh đó, và $\alpha \in [0, 1]$ là siêu tham số kiểm soát mức độ di dời:
\begin{itemize}
  \item Khi $\alpha = 1$: Toàn bộ độ khó được chuyển sang trọng số, kích hoạt trở nên cực kỳ phẳng và dễ lượng tử hóa.
  \item Khi $\alpha = 0$: Trọng số được giữ nguyên độ mượt tối đa.
  \item Khi $\alpha = 0.5$: Điểm cân bằng hình học tối ưu, chia sẻ đều độ lệch biên độ giữa kích hoạt và trọng số, cho phép lượng tử hóa đồng thời cả activations và weights sang INT8 mà không làm mất độ chính xác.
\end{itemize}

\subsection{Toán Học Giải Tích Bậc Hai: Optimal Brain Surgeon (OBS) \& Thuật Toán GPTQ}
\label{sec:obs-gptq-math}

Để hiểu tại sao GPTQ có thể nén một mô hình ngôn ngữ 70 tỷ tham số xuống INT4 hoặc INT3 mà không làm suy giảm chỉ số Perplexity, chúng ta phải xuất phát từ nền tảng giải tích của thuật toán \term{Optimal Brain Surgeon (OBS)} do Babak Hassibi và David Stork đề xuất năm 1993.

\paragraph{1. Động cơ: Tại sao làm tròn độc lập (Round-To-Nearest - RTN) thất bại thảm hại ở 4-bit?}
Xét một tầng tuyến tính bất kỳ của Transformer: $\mathbf{Y} = \mathbf{W} \mathbf{X}$, trong đó $\mathbf{W} \in \mathbb{R}^{d_{\text{out}} \times d_{\text{in}}}$ là ma trận trọng số và $\mathbf{X} \in \mathbb{R}^{d_{\text{in}} \times N}$ là ma trận kích hoạt đầu vào thu thập từ $N$ mẫu dữ liệu hiệu chuẩn (Calibration tokens).
Khi nén trọng số từ FP16 xuống INT4, phương pháp đơn giản nhất là làm tròn độc lập từng phần tử (Round-To-Nearest -- RTN):
\begin{equation}
\hat{w}_{ij} = \text{round}\left( \frac{w_{ij}}{S} \right) \cdot S
\end{equation}
Ở độ phân giải 8-bit ($256$ mức rời rạc), khoảng cách giữa hai mức lượng tử hóa $\Delta$ đủ nhỏ để sai số làm tròn $e = w_{ij} - \hat{w}_{ij}$ triệt tiêu lẫn nhau qua phân phối chuẩn. Nhưng ở độ phân giải 4-bit ($16$ mức rời rạc), khoảng cách $\Delta$ trở nên khổng lồ. Hàng triệu sai số cục bộ độc lập cộng dồn lại theo chiều sâu hàng chục tầng Transformer, khiến đầu ra $\widehat{\mathbf{Y}} = \widehat{\mathbf{W}} \mathbf{X}$ bị trôi dạt hoàn toàn khỏi phân phối ban đầu, dẫn đến hiện tượng sụp đổ suy luận (Perplexity bùng nổ lên hàng nghìn).

\paragraph{2. Khai triển chuỗi Taylor bậc hai của hàm mục tiêu sai số tái tạo}:
Thay vì tối ưu hóa mù quáng trên từng trọng số đơn lẻ, mục tiêu thực sự của lượng tử hóa bảo toàn độ chính xác là cực tiểu hóa bình phương sai số tái tạo kích hoạt đầu ra:
\begin{equation}
\min_{\widehat{\mathbf{W}}} \mathcal{E} = \|\mathbf{W} \mathbf{X} - \widehat{\mathbf{W}} \mathbf{X}\|_F^2 = \sum_{r=1}^{d_{\text{out}}} \|\mathbf{w}_r \mathbf{X} - \widehat{\mathbf{w}}_r \mathbf{X}\|_2^2
\end{equation}
Hàm mục tiêu này khả ly theo từng hàng của ma trận trọng số. Do đó, bài toán quy về việc tối ưu hóa độc lập cho từng vector trọng số hàng $\mathbf{w} \in \mathbb{R}^{1 \times d_{\text{in}}}$. Đặt $\Delta \mathbf{w} = \widehat{\mathbf{w}} - \mathbf{w}$ là vector sai số lượng tử hóa. Khai triển Taylor của hàm sai số $\mathcal{E}(\mathbf{w} + \Delta \mathbf{w})$ quanh điểm tối ưu ban đầu $\mathbf{w}$:
\begin{equation}
\mathcal{E}(\mathbf{w} + \Delta \mathbf{w}) = \mathcal{E}(\mathbf{w}) + \mathbf{g}^T \Delta \mathbf{w} + \frac{1}{2} \Delta \mathbf{w}^T \mathbf{H} \Delta \mathbf{w} + \mathcal{O}(\|\Delta \mathbf{w}\|^3)
\end{equation}
Trong đó:
\begin{itemize}
  \item Vì mô hình đã hội tụ sau huấn luyện, vector gradient bậc nhất triệt tiêu về 0: $\mathbf{g} = \nabla_{\mathbf{w}} \mathcal{E} = \mathbf{0}$.
  \item Ma trận đạo hàm bậc hai (Hessian Matrix) $\mathbf{H} \in \mathbb{R}^{d_{\text{in}} \times d_{\text{in}}}$ được xác định tường minh bằng:
  \begin{equation}
  \mathbf{H} = \nabla^2_{\mathbf{w}} \mathcal{E} = 2 \mathbf{X} \mathbf{X}^T
  \end{equation}
  \item Bỏ qua các số hạng bậc ba và cao hơn, mức biến thiên sai số thuần túy do lượng tử hóa gây ra là:
  \begin{equation}
  \Delta \mathcal{E} \approx \frac{1}{2} \Delta \mathbf{w}^T \mathbf{H} \Delta \mathbf{w}
  \end{equation}
\end{itemize}

\paragraph{3. Thiết lập bài toán tối ưu có ràng buộc với Nhân tử Lagrange}:
Giả sử ta tiến hành lượng tử hóa trọng số tại vị trí cột thứ $q$ ($w_q \to \hat{w}_q$), tức là cố định sai số lượng tử hóa cục bộ tại vị trí $q$:
\begin{equation}
\mathbf{e}_q^T \Delta \mathbf{w} + (w_q - \hat{w}_q) = 0
\end{equation}
trong đó $\mathbf{e}_q$ là vector đơn vị có phần tử thứ $q$ bằng $1$ và các phần tử khác bằng $0$.
Để giảm thiểu tối đa sự gia tăng sai số toàn cục $\Delta \mathcal{E}$, chúng ta không giữ nguyên các trọng số chưa nén còn lại, mà chủ động điều chỉnh chúng để hấp thụ phần sai số của $w_q$. Bài toán tối ưu hóa có ràng buộc được phát biểu:
\begin{equation}
\min_{\Delta \mathbf{w}} \frac{1}{2} \Delta \mathbf{w}^T \mathbf{H} \Delta \mathbf{w} \quad \text{với điều kiện ràng buộc } \mathbf{e}_q^T \Delta \mathbf{w} + (w_q - \hat{w}_q) = 0
\end{equation}

Thiết lập hàm Lagrange:
\begin{equation}
\mathcal{L}(\Delta \mathbf{w}, \lambda) = \frac{1}{2} \Delta \mathbf{w}^T \mathbf{H} \Delta \mathbf{w} + \lambda \left( \mathbf{e}_q^T \Delta \mathbf{w} + w_q - \hat{w}_q \right)
\end{equation}
với $\lambda \in \mathbb{R}$ là nhân tử Lagrange. Lấy đạo hàm riêng theo $\Delta \mathbf{w}$ và cho bằng $0$:
\begin{equation}
\nabla_{\Delta \mathbf{w}} \mathcal{L} = \mathbf{H} \Delta \mathbf{w} + \lambda \mathbf{e}_q = \mathbf{0} \implies \Delta \mathbf{w} = -\lambda \mathbf{H}^{-1} \mathbf{e}_q
\end{equation}
Nhân cả hai vế với $\mathbf{e}_q^T$:
\begin{equation}
\mathbf{e}_q^T \Delta \mathbf{w} = -\lambda \mathbf{e}_q^T \mathbf{H}^{-1} \mathbf{e}_q = -\lambda [\mathbf{H}^{-1}]_{qq}
\end{equation}
Mặt khác, theo điều kiện ràng buộc: $\mathbf{e}_q^T \Delta \mathbf{w} = -(w_q - \hat{w}_q) = \hat{w}_q - w_q$. Suy ra:
\begin{equation}
\lambda = \frac{w_q - \hat{w}_q}{[\mathbf{H}^{-1}]_{qq}}
\end{equation}
Thay $\lambda$ ngược trở lại phương trình của $\Delta \mathbf{w}$, ta thu được công thức bù sai số kinh điển của Optimal Brain Surgeon:
\begin{equation}
\Delta \mathbf{w} = - \frac{w_q - \hat{w}_q}{[\mathbf{H}^{-1}]_{qq}} \cdot \mathbf{H}^{-1}_{:, q}
\end{equation}
trong đó $\mathbf{H}^{-1}_{:, q}$ là cột thứ $q$ của ma trận Hessian nghịch đảo, và $[\mathbf{H}^{-1}]_{qq}$ là phần tử nằm trên đường chéo chính tại hàng $q$, cột $q$.

Mức gia tăng sai số toàn cục tối thiểu khi lượng tử hóa trọng số thứ $q$ là:
\begin{equation}
\Delta \mathcal{E}_q = \frac{1}{2} \Delta \mathbf{w}^T \mathbf{H} \Delta \mathbf{w} = \frac{1}{2} \left(-\lambda \mathbf{e}_q^T \mathbf{H}^{-1}\right) \mathbf{H} \left(-\lambda \mathbf{H}^{-1} \mathbf{e}_q\right) = \frac{1}{2} \lambda^2 [\mathbf{H}^{-1}]_{qq} = \frac{1}{2} \frac{(w_q - \hat{w}_q)^2}{[\mathbf{H}^{-1}]_{qq}}
\end{equation}

\paragraph{4. Thuật toán GPTQ hiện đại: Đột phá vượt qua giới hạn tính toán}:
Mặc dù công thức OBS có vẻ đẹp giải tích tuyệt đối, việc áp dụng nó cho mô hình LLM hiện đại từng gặp phải rào cản tính toán không thể vượt qua:
\begin{itemize}
  \item Trong OBS cổ điển, sau khi lượng tử hóa trọng số thứ $q$, ma trận Hessian phải được cập nhật lại theo công thức Sherman-Morrison với độ phức tạp $\mathcal{O}(d^2)$ cho mỗi tham số, dẫn đến tổng độ phức tạp toàn tầng là $\mathcal{O}(d^4)$. Với ma trận ẩn $d = 8192$ của Llama-3-70B, $d^4 \approx 4.5 \times 10^{15}$ phép tính, đòi hỏi hàng tuần tính toán chỉ cho một tầng duy nhất!
  \item \textbf{Cải tiến 1: Khử ma trận qua Phân rã Cholesky (Cholesky Factorization)}:
  Frantar et al. (2022) nhận ra rằng nếu cố định thứ tự lượng tử hóa các cột trọng số từ $1$ đến $d_{\text{in}}$, ta có thể tính toán trước ma trận Hessian nghịch đảo $\mathbf{H}^{-1}$ một lần duy nhất, sau đó dùng phân rã Cholesky của $\mathbf{H}^{-1}$:
  \begin{equation}
  \mathbf{H}^{-1} = \mathbf{L} \mathbf{L}^T
  \end{equation}
  trong đó $\mathbf{L}$ là ma trận tam giác dưới. Nhờ đó, việc cập nhật bù sai số cho toàn bộ các cột còn lại có thể thực hiện theo quy tắc khử Gaussian trực tiếp trên các cột của ma trận nghịch đảo mà không cần đảo ngược lại ma trận!
  \item \textbf{Cải tiến 2: Điều hòa Damping Factor $\lambda \mathbf{I}$ chống suy biến}:
  Trong thực tế, ma trận kích hoạt $\mathbf{X}$ có thể chứa các kênh tương quan tuyến tính hoặc các nơ-ron chết (giá trị 0), khiến ma trận Hessian $\mathbf{H} = 2 \mathbf{X} \mathbf{X}^T$ có các trị riêng gần bằng $0$ (ill-conditioned). Để đảm bảo tính khả nghịch và ổn định số học, GPTQ bổ sung một hệ số điều hòa đường chéo:
  \begin{equation}
  \widehat{\mathbf{H}} = 2 \mathbf{X} \mathbf{X}^T + \lambda \mathbf{I}, \quad \text{với } \lambda = \alpha \cdot \frac{\text{Tr}(2 \mathbf{X} \mathbf{X}^T)}{d_{\text{in}}} \quad (\alpha \approx 0.01)
  \end{equation}
  \item \textbf{Cải tiến 3: Cập nhật trễ theo khối (Lazy Batch / Block Updates)}:
  Thay vì cập nhật toàn bộ ma trận trọng số $\mathbf{W}$ sau mỗi cột đơn lẻ (gây nghẽn băng thông bộ nhớ GPU Memory Bandwidth do các thao tác ghi liên tục), GPTQ gom $B = 128$ cột lại thành một khối. Quá trình lượng tử hóa và bù sai số được tích lũy cục bộ bên trong khối, sau đó một phép nhân ma trận lớn duy nhất (\code{GEMM}) được thực thi trên GPU Tensor Cores để cập nhật toàn bộ phần còn lại của ma trận trọng số. Kỹ thuật này giúp tăng tốc độ nén lên hơn 10 lần!
\end{itemize}

\begin{example}[Tính Toán Bằng Tay Từng Bước Thuật Toán OBS Cho Tensor Kích Hoạt $2 \times 3$]
Xét ma trận kích hoạt đầu vào $\mathbf{X} \in \mathbb{R}^{2 \times 3}$ và vector trọng số ban đầu $\mathbf{w} \in \mathbb{R}^{1 \times 2}$:
\begin{equation*}
\mathbf{X} = \begin{bmatrix} 1 & 2 & 1 \\ 0 & 1 & 2 \end{bmatrix}, \qquad \mathbf{w} = \begin{bmatrix} 1.5 & -2.0 \end{bmatrix}
\end{equation*}
Đầu ra kích hoạt chuẩn ban đầu:
\begin{equation*}
\mathbf{y} = \mathbf{w} \mathbf{X} = \begin{bmatrix} 1.5 & -2.0 \end{bmatrix} \begin{bmatrix} 1 & 2 & 1 \\ 0 & 1 & 2 \end{bmatrix} = \begin{bmatrix} 1.5 & 1.0 & -2.5 \end{bmatrix}
\end{equation*}

\textbf{Bước 1: Tính ma trận Hessian $\mathbf{H}$ và ma trận nghịch đảo $\mathbf{H}^{-1}$}:
\begin{equation*}
\mathbf{X} \mathbf{X}^T = \begin{bmatrix} 1 & 2 & 1 \\ 0 & 1 & 2 \end{bmatrix} \begin{bmatrix} 1 & 0 \\ 2 & 1 \\ 1 & 2 \end{bmatrix} = \begin{bmatrix} 1^2 + 2^2 + 1^2 & 1\cdot 0 + 2\cdot 1 + 1\cdot 2 \\ 0\cdot 1 + 1\cdot 2 + 2\cdot 1 & 0^2 + 1^2 + 2^2 \end{bmatrix} = \begin{bmatrix} 6 & 4 \\ 4 & 5 \end{bmatrix}
\end{equation*}
Ma trận Hessian:
\begin{equation*}
\mathbf{H} = 2 \mathbf{X} \mathbf{X}^T = \begin{bmatrix} 12 & 8 \\ 8 & 10 \end{bmatrix}
\end{equation*}
Định thức: $\det(\mathbf{H}) = 12 \cdot 10 - 8 \cdot 8 = 120 - 64 = 56$.
Ma trận Hessian nghịch đảo:
\begin{equation*}
\mathbf{H}^{-1} = \frac{1}{56} \begin{bmatrix} 10 & -8 \\ -8 & 12 \end{bmatrix} = \begin{bmatrix} \frac{5}{28} & -\frac{1}{7} \\ -\frac{1}{7} & \frac{3}{14} \end{bmatrix} \approx \begin{bmatrix} 0.17857 & -0.14286 \\ -0.14286 & 0.21429 \end{bmatrix}
\end{equation*}

\textbf{Bước 2: Lượng tử hóa trọng số thứ nhất ($q=1$)}:
Giả sử ta làm tròn $w_1 = 1.5$ sang số nguyên gần nhất $\hat{w}_1 = 2.0$.
Sai số làm tròn cục bộ: $w_1 - \hat{w}_1 = 1.5 - 2.0 = -0.5$.
Phần tử đường chéo $[\mathbf{H}^{-1}]_{11} = \frac{5}{28} \approx 0.17857$.
Áp dụng công thức cập nhật bù trừ OBS cho vector trọng số:
\begin{equation*}
\Delta \mathbf{w} = - \frac{w_1 - \hat{w}_1}{[\mathbf{H}^{-1}]_{11}} \cdot \mathbf{H}^{-1}_{:, 1} = - \frac{-0.5}{\frac{5}{28}} \begin{bmatrix} \frac{5}{28} \\ -\frac{1}{7} \end{bmatrix} = \frac{0.5 \times 28}{5} \begin{bmatrix} \frac{5}{28} \\ -\frac{4}{28} \end{bmatrix} = 2.8 \begin{bmatrix} \frac{5}{28} \\ -\frac{4}{28} \end{bmatrix} = \begin{bmatrix} 0.5 \\ -0.4 \end{bmatrix}
\end{equation*}
Như vậy:
\begin{itemize}
  \item Trọng số thứ nhất nhận giá trị đã lượng tử hóa: $w_1 + \Delta w_1 = 1.5 + 0.5 = 2.0 = \hat{w}_1$.
  \item Trọng số thứ hai được tự động bù trừ: $w_2' = w_2 + \Delta w_2 = -2.0 + (-0.4) = -2.4$.
\end{itemize}

\textbf{Bước 3: So sánh sai số kích hoạt đầu ra giữa OBS và RTN}:
\begin{itemize}
  \item \textit{Nếu dùng Round-To-Nearest (RTN)}: Cả hai trọng số bị làm tròn độc lập: $\widehat{\mathbf{w}}_{\text{RTN}} = \begin{bmatrix} 2.0 & -2.0 \end{bmatrix}$.
  Đầu ra tái tạo: $\widehat{\mathbf{y}}_{\text{RTN}} = \begin{bmatrix} 2.0 & -2.0 \end{bmatrix} \begin{bmatrix} 1 & 2 & 1 \\ 0 & 1 & 2 \end{bmatrix} = \begin{bmatrix} 2.0 & 2.0 & -2.0 \end{bmatrix}$.
  Bình phương sai số tái tạo kích hoạt:
  \begin{equation*}
  \|\mathbf{y} - \widehat{\mathbf{y}}_{\text{RTN}}\|_2^2 = (1.5 - 2.0)^2 + (1.0 - 2.0)^2 + (-2.5 - (-2.0))^2 = 0.25 + 1.0 + 0.25 = \mathbf{1.50}
  \end{equation*}
  \item \textit{Nếu dùng OBS}: Vector trọng số sau khi lượng tử hóa $w_1$ và bù trừ $w_2$ là $\mathbf{w}_{\text{OBS}} = \begin{bmatrix} 2.0 & -2.4 \end{bmatrix}$.
  Đầu ra tái tạo: $\widehat{\mathbf{y}}_{\text{OBS}} = \begin{bmatrix} 2.0 & -2.4 \end{bmatrix} \begin{bmatrix} 1 & 2 & 1 \\ 0 & 1 & 2 \end{bmatrix} = \begin{bmatrix} 2.0 & 1.6 & -2.8 \end{bmatrix}$.
  Bình phương sai số tái tạo kích hoạt:
  \begin{equation*}
  \|\mathbf{y} - \widehat{\mathbf{y}}_{\text{OBS}}\|_2^2 = (1.5 - 2.0)^2 + (1.0 - 1.6)^2 + (-2.5 - (-2.8))^2 = 0.25 + 0.36 + 0.09 = \mathbf{0.70}
  \end{equation*}
\end{itemize}
\textbf{Kết luận định lượng}: Cơ chế bù trừ bậc hai của OBS đã triệt tiêu hơn $53\%$ sai số tái tạo kích hoạt so với làm tròn độc lập RTN ngay từ bước đầu tiên! Khi tiếp tục lượng tử hóa $w_2$, OBS sẽ đảm bảo tổng sai số trên toàn mạng được triệt tiêu tối ưu nhất.
\end{example}

\subsection{AWQ: Activation-aware Weight Quantization \& Kỹ Thuật LayerNorm Folding}
\label{sec:awq-math-optimization}

Được Lin et al. (MIT HAN Lab, MLSys 2024) đề xuất, \textbf{AWQ (Activation-aware Weight Quantization)} giải quyết triệt để sự đánh đổi giữa hiệu năng và khả năng tổng quát hóa trong lượng tử hóa 4-bit mà không cần huấn luyện lại hay tính toán ma trận Hessian cồng kềnh như GPTQ.

\paragraph{1. Quan Sát Thực Nghiệm Cốt Lõi: Trọng Số Không Bình Đẳng}:
Bằng việc phân tích các mô hình ngôn ngữ lớn từ 7B đến 70B, Lin et al. nhận thấy:
\begin{itemize}
  \item Nếu ta lựa chọn ngẫu nhiên $1\%$ trọng số để giữ lại ở định dạng FP16 và lượng tử hóa $99\%$ còn lại sang INT4, mô hình vẫn sụp đổ về mặt chất lượng (Perplexity tăng vọt).
  \item Nếu ta chọn $1\%$ trọng số có \textbf{biên độ tuyệt đối lớn nhất} của ma trận $\mathbf{W}$ để giữ FP16, chất lượng chỉ cải thiện khiêm tốn.
  \item Tuy nhiên, nếu ta chọn $1\%$ trọng số kết nối trực tiếp với các kênh \textbf{kích hoạt có biên độ trung bình lớn nhất} ($\mathbf{s}_X = \frac{1}{N}\sum |\mathbf{X}|$), mô hình bảo toàn được hơn $95\%$ năng lực suy luận và lập trình!
\end{itemize}
Điều này dẫn đến chân lý nền tảng: \textit{Tầm quan trọng của một tham số không nằm ở độ lớn tĩnh của chính nó, mà được quyết định bởi độ lớn của tín hiệu kích hoạt chảy qua nó.}

\paragraph{2. Thiết Lập Bài Toán Tối Ưu Hóa Tái Tạo Giải Tích}:
Cho tầng tuyến tính $\mathbf{Y} = \mathbf{W} \mathbf{X}$ với $\mathbf{W} \in \mathbb{R}^{d_{\text{out}} \times d_{\text{in}}}$ và $\mathbf{X} \in \mathbb{R}^{d_{\text{in}} \times N}$.
Nếu lưu riêng $1\%$ trọng số ở dạng FP16 (Mixed-Precision), phần cứng GPU sẽ bị phân mảnh bộ nhớ và không thể tận dụng các nhân ma trận SIMD/Tensor Core hiệu quả. AWQ đưa ra ý tưởng tinh tế: \textbf{Bảo vệ các trọng số quan trọng bằng cách co giãn biên độ của chúng trước khi lượng tử hóa!}

Xét một vector co giãn trên từng kênh $\mathbf{s} \in \mathbb{R}^{d_{\text{in}}}$ với $s_j > 0$. Phép nhân ma trận được biến đổi đồng nhất:
\begin{equation}
\mathbf{Y} = \mathbf{W} \mathbf{X} = \left( \mathbf{W} \cdot \text{diag}(\mathbf{s}) \right) \cdot \left( \text{diag}(\mathbf{s})^{-1} \mathbf{X} \right) = \mathbf{W}' \mathbf{X}'
\end{equation}
Mục tiêu là tìm vector $\mathbf{s}^*$ sao cho sau khi lượng tử hóa ma trận đã co giãn $\mathbf{W}'$, sai số tái tạo kích hoạt đầu ra đạt cực tiểu:
\begin{equation}
\label{eq:awq-optimization-objective}
\min_{\mathbf{s}} \mathcal{L}(\mathbf{s}) = \left\| \mathbf{W} \mathbf{X} - Q\left(\mathbf{W} \cdot \text{diag}(\mathbf{s})\right) \cdot \text{diag}(\mathbf{s})^{-1} \mathbf{X} \right\|_F^2
\end{equation}
trong đó $Q(\cdot)$ là toán tử lượng tử hóa số nguyên 4-bit (INT4).

\paragraph{Tại sao việc nhân $\mathbf{s} > 1$ lại bảo vệ được trọng số?}
Xét kênh đầu vào thứ $j$ có kích hoạt rất lớn. Khi ta chọn $s_j > 1$, các phần tử trên cột thứ $j$ của ma trận trọng số $\mathbf{w}_{:, j}' = \mathbf{w}_{:, j} \cdot s_j$ được phóng đại biên độ.
Sai số làm tròn của toán tử lượng tử hóa đối xứng $Q(\cdot)$ luôn bị chặn trong khoảng nửa bước lượng tử: $|\epsilon| \le \frac{\Delta}{2}$.
Khi tính toán đầu ra, giá trị đã lượng tử hóa được chia lại cho $s_j$:
\begin{equation}
\widehat{\mathbf{w}}_{:, j} = \frac{Q(\mathbf{w}_{:, j} \cdot s_j)}{s_j} = \frac{\mathbf{w}_{:, j} \cdot s_j + \boldsymbol{\epsilon}}{s_j} = \mathbf{w}_{:, j} + \frac{\boldsymbol{\epsilon}}{s_j}
\end{equation}
Sai số trọng số hiệu dụng tác động lên kênh kích hoạt lớn bị \textbf{triệt tiêu đi $s_j$ lần}! Bằng cách đó, thành phần sai số lớn nhất trong phép nhân $\sum_j \Delta \mathbf{w}_{:, j} \mathbf{x}_j$ bị khống chế hoàn toàn.

\paragraph{3. Thuật Toán Tìm Kiếm Theo Lưới Tối Ưu (Grid Search)}:
Việc giải trực tiếp bài toán (\ref{eq:awq-optimization-objective}) bằng hạ gradient (SGD) rất dễ rơi vào cực tiểu cục bộ và phụ thuộc vào bộ tối ưu. AWQ đơn giản hóa bài toán bằng cách tham số hóa vector $\mathbf{s}$ theo dạng lũy thừa của độ lớn kích hoạt:
\begin{equation}
\mathbf{s} = \mathbf{s}_X^\alpha, \qquad \text{với } s_{X, j} = \frac{1}{N} \sum_{i=1}^N |X_{i, j}|
\end{equation}
trong đó $\alpha \in [0, 1]$ là siêu tham số kiểm soát mức độ bảo vệ.
Hệ thống tiến hành tìm kiếm lưới (Grid Search) trên đoạn $\alpha \in [0, 1]$ với bước nhảy $\Delta \alpha = 0.05$ (chỉ 20 điểm đánh giá) trên một tập dữ liệu hiệu chuẩn cực nhỏ (chỉ cần $32 - 128$ chuỗi câu, không cần nhãn). Giá trị $\alpha^*$ tối ưu thỏa mãn:
\begin{equation}
\alpha^* = \arg\min_{\alpha \in [0, 1]} \left\| \mathbf{W} \mathbf{X} - Q\left(\mathbf{W} \cdot \text{diag}(\mathbf{s}_X^\alpha)\right) \cdot \text{diag}(\mathbf{s}_X^\alpha)^{-1} \mathbf{X} \right\|_F^2
\end{equation}
Nhờ không tối ưu hóa tự do hàng triệu tham số như GPTQ, AWQ hoàn toàn miễn nhiễm với hiện tượng quá khớp (Overfitting) vào tập dữ liệu hiệu chuẩn, bảo toàn trọn vẹn năng lực suy luận toán học và lập trình trên các tập kiểm thử mở rộng.

\paragraph{4. Kỹ Thuật Gập Hệ Số Scale Vào LayerNorm (Zero-Latency LayerNorm Folding)}:
Một rào cản lớn khi áp dụng hệ số co giãn động là chi phí thời gian chạy: nếu mỗi token trước khi nhân ma trận $\mathbf{W}'$ phải trải qua một phép chia vector $\mathbf{x} \cdot \mathbf{s}^{-1}$, thông lượng suy luận sẽ bị suy giảm nghiêm trọng do nghẽn băng thông bộ nhớ.

AWQ giải quyết vấn đề này một cách tuyệt mỹ thông qua cấu trúc mạng Transformer: Mọi tầng tuyến tính $\mathbf{W}_{\text{in}}$ (như ma trận chiếu $\mathbf{W}_q, \mathbf{W}_k, \mathbf{W}_v$ trong Attention hoặc ma trận $\mathbf{W}_{\text{gate}}, \mathbf{W}_{\text{up}}$ trong MLP) đều nhận đầu vào trực tiếp từ một tầng chuẩn hóa (LayerNorm hoặc RMSNorm).

Xét công thức tính của RMSNorm:
\begin{equation}
\mathbf{x} = \text{RMSNorm}(\mathbf{z}) = \frac{\mathbf{z}}{\text{RMS}(\mathbf{z})} \odot \boldsymbol{\gamma}_{\text{LN}} + \boldsymbol{\beta}_{\text{LN}}
\end{equation}
trong đó $\boldsymbol{\gamma}_{\text{LN}}$ là trọng số co giãn học được của LayerNorm.
Thay vì nhân $\mathbf{s}^{-1}$ vào $\mathbf{x}$ tại thời gian chạy, AWQ \textbf{gập trực tiếp hệ số này vào trọng số tĩnh của LayerNorm}:
\begin{equation}
\boldsymbol{\gamma}_{\text{LN}}' = \boldsymbol{\gamma}_{\text{LN}} \oslash \mathbf{s}, \qquad \boldsymbol{\beta}_{\text{LN}}' = \boldsymbol{\beta}_{\text{LN}} \oslash \mathbf{s}
\end{equation}
trong đó $\oslash$ là phép chia theo từng phần tử.
Khi đó, đầu ra của tầng LayerNorm tự động là $\mathbf{x}' = \mathbf{x} \oslash \mathbf{s}$.
Tầng tuyến tính tiếp theo chỉ cần thực hiện phép nhân ma trận thuần túy:
\begin{equation}
\mathbf{Y} = \mathbf{W}' \mathbf{x}' = Q(\mathbf{W} \cdot \text{diag}(\mathbf{s})) \cdot (\mathbf{x} \oslash \mathbf{s}) \approx \mathbf{W} \mathbf{x}
\end{equation}
\textbf{Kết luận}: Chi phí trễ suy luận bổ sung của AWQ là \textbf{chính xác bằng 0 ns (Zero Runtime Overhead)}! Toàn bộ quá trình bảo vệ được hoàn tất bằng các phép biến đổi trọng số tĩnh trước khi triển khai mô hình lên máy chủ.

\subsection{Định Dạng GGUF \& Thư Viện llama.cpp}
Được Georgi Gerganov phát triển, định dạng nhị phân \textbf{GGUF} (thay thế GGML) là chuẩn mực lưu trữ mô hình tối ưu cho CPU và thiết bị biên:
\begin{itemize}
  \item \textbf{Cơ chế nạp qua Memory Mapping (\code{mmap})}: Toàn bộ trọng số được sắp xếp thẳng hàng nhị phân. Khi ứng dụng khởi động, hệ điều hành chỉ cần ánh xạ tệp từ ổ SSD trực tiếp vào không gian địa chỉ bộ nhớ mà không cần tốn chu kỳ CPU để parse hay sao chép dữ liệu.
  \item \textbf{Cơ chế k-quants (Q4\_K\_M, Q5\_K\_M, Q8\_0)}: GGUF cho phép lượng tử hóa không đồng nhất: Các tầng nhạy cảm (như Embedding và Attention Output) được giữ ở mức chính xác cao (INT6 hoặc INT8), trong khi các ma trận MLP lớn được nén xuống INT4 để tối ưu tỷ lệ nén/độ chính xác.
\end{itemize}

\subsection{Lượng Tử Hóa Cực Hạn 1.58-Bit: BitNet b1.58 \& Mạng Nơ-ron Bậc Ba (Ternary Weights)}
\label{sec:bitnet-158b}

Trong các thuật toán lượng tử hóa truyền thống (INT8, INT4, NF4), mục tiêu là nén một mô hình đã huấn luyện ở dạng FP16 xuống một biểu diễn số nguyên có độ phân giải thấp hơn, chấp nhận đánh đổi một lượng sai số nhất định. Tuy nhiên, kiến trúc xử lý của phần cứng bên dưới vẫn là các bộ nhân số học (Hardware Multipliers) trong các lõi Tensor Cores. Năm 2024, một bước nhảy vọt mang tính phá vỡ giới hạn xuất hiện từ công trình \textit{The Era of 1-bit LLMs: All Large Language Models are in 1.58 Bits} của nhóm nghiên cứu Microsoft Research (Wang et al.): \textbf{BitNet b1.58}.

Thay vì dừng lại ở phép nhân ma trận dấu phẩy động, BitNet b1.58 định nghĩa lại hoàn toàn đại số nơ-ron: \textbf{Mọi tham số trọng số trong mô hình ngôn ngữ đều thuộc tập giá trị bậc ba (Ternary Set) $\{-1, 0, +1\}$, triệt tiêu hoàn toàn phép nhân số thực trong tầng tuyến tính và thay thế $100\%$ bằng các phép cộng và trừ số nguyên!}

\paragraph{1. Nền Tảng Lý Thuyết Thông Tin: Tại Sao Lại Là 1.58 Bit?}
Trong hệ thống nhị phân, một bit có thể biểu diễn $2^1 = 2$ trạng thái ($\{0, 1\}$ hoặc $\{-1, +1\}$). Với mô hình nhị phân 1-bit thuần túy (1-bit Binary Network như BitNet ban đầu), trọng số chỉ nhận giá trị $-1$ hoặc $+1$. Mô hình nhị phân 1-bit gặp phải suy giảm chất lượng biểu diễn nghiêm trọng vì thiếu vắng khả năng lọc tính năng: một trọng số bắt buộc phải kích hoạt hoặc ức chế, không thể ``tắt'' hoàn toàn một tín hiệu không quan trọng.

BitNet b1.58 bổ sung trạng thái thứ ba: $0$ (giá trị triệt tiêu / Feature Filtering). Theo công thức Entropi thông tin của Shannon, số bit tối thiểu để mã hóa một biến ngẫu nhiên rời rạc có 3 trạng thái phân phối đều là:
\begin{equation}
\label{eq:bitnet-shannon-capacity}
C = \log_2(3) = \frac{\ln 3}{\ln 2} \approx 1.5849625 \text{ bits/tham số}
\end{equation}
Chính vì vậy, kiến trúc này được định danh chính xác là \textbf{BitNet 1.58-bit (b1.58)}.

\paragraph{Khả năng đóng gói nhị phân (Bit-packing) tối ưu trong thực tế}:
Mặc dù máy tính số không hỗ trợ kiểu dữ liệu cơ sở $1.58$-bit, chúng ta có thể đóng gói nhiều trọng số bậc ba vào một từ nhị phân chuẩn mà hầu như không lãng phí dung lượng:
\begin{equation}
3^5 = 243 \le 256 = 2^8
\end{equation}
Một byte duy nhất ($8\text{ bits}$) có thể chứa trọn vẹn $5$ trọng số bậc ba độc lập ($243$ trạng thái trong tổng số $256$ trạng thái khả dĩ), đạt hiệu suất lưu trữ thực tế:
\begin{equation}
\frac{8\text{ bits}}{5\text{ tham số}} = 1.60\text{ bits/tham số}
\end{equation}
Mức hao phí lưu trữ lý thuyết chỉ là $\frac{1.60 - 1.585}{1.585} \approx 0.95\%$!

\paragraph{2. Công Thức Lượng Tử Hóa Absmean Bậc Ba}:
Cho ma trận trọng số $\mathbf{W} \in \mathbb{R}^{d_{\text{out}} \times d_{\text{in}}}$. Quá trình lượng tử hóa bậc ba trong BitNet b1.58 sử dụng hàm co giãn trung bình biên độ tuyệt đối (\term{Absmean Scaling}):
\begin{equation}
\gamma = \frac{1}{d_{\text{out}} \times d_{\text{in}}} \sum_{i=1}^{d_{\text{out}}} \sum_{j=1}^{d_{\text{in}}} |W_{i,j}|
\end{equation}
Ma trận trọng số bậc ba $\widetilde{\mathbf{W}} \in \{-1, 0, +1\}^{d_{\text{out}} \times d_{\text{in}}}$ được xác định bởi hàm làm tròn và cắt ngưỡng:
\begin{equation}
\label{eq:bitnet-weight-quant}
\widetilde{W}_{i,j} = \text{RoundClip}\left( \frac{W_{i,j}}{\gamma + \epsilon}, -1, 1 \right) = \max\left(-1, \min\left(1, \left\lfloor \frac{W_{i,j}}{\gamma + \epsilon} \right\rceil\right)\right)
\end{equation}
Trong đó $\epsilon > 0$ là hằng số bảo vệ chống chia cho 0. Khi cần giải mã để phân tích sai số: $\widehat{\mathbf{W}} = \gamma \times \widetilde{\mathbf{W}}$.

\paragraph{3. Lượng Tử Hóa Kích Hoạt 8-bit Đối Xứng}:
Để đảm bảo khả năng tính toán thuần số nguyên, véc-tơ kích hoạt đầu vào $\mathbf{x} \in \mathbb{R}^{1 \times d_{\text{in}}}$ được lượng tử hóa sang định dạng số nguyên 8-bit (INT8) theo từng token:
\begin{equation}
s_x = \frac{\max(|\mathbf{x}|)}{Q_b}, \qquad \text{với } Q_b = 2^{b-1} - 1 = 127
\end{equation}
\begin{equation}
\mathbf{x}_q = \text{clip}\left( \left\lfloor \frac{\mathbf{x}}{s_x} \right\rceil, -Q_b, Q_b \right) \in [-127, 127]^{1 \times d_{\text{in}}}
\end{equation}

\paragraph{4. Cách Mạng Phần Cứng: Triệt Tiêu Phép Nhân Bằng Phép Cộng Số Nguyên}:
Xét phép nhân ma trận - vector cơ bản trong tầng FFN hoặc Attention của Transformer: $\mathbf{y} = \mathbf{x} \mathbf{W}^T$.
Trong kiến trúc FP16/BF16 thông thường, để tính một phần tử đầu ra $y_j$:
\begin{equation}
y_j = \sum_{k=1}^{d_{\text{in}}} x_k W_{j,k} \implies \text{Đòi hỏi } d_{\text{in}} \text{ phép nhân FP16 và } d_{\text{in}}-1 \text{ phép cộng FP16.}
\end{equation}
Trong BitNet b1.58, vì $\widetilde{W}_{j,k} \in \{-1, 0, +1\}$, phép tính được phân tách thành:
\begin{equation}
\label{eq:bitnet-pure-additions}
y_{q, j} = \sum_{k=1}^{d_{\text{in}}} x_{q, k} \widetilde{W}_{j,k} = \sum_{k: \widetilde{W}_{j,k} = +1} x_{q, k} - \sum_{k: \widetilde{W}_{j,k} = -1} x_{q, k}
\end{equation}
Các trọng số có $\widetilde{W}_{j,k} = 0$ hoàn toàn bị bỏ qua!
Giá trị đầu ra số thực được phục hồi bằng một phép nhân vô hướng duy nhất trên toàn bộ vector:
\begin{equation}
y_j = y_{q, j} \times (s_x \cdot \gamma)
\end{equation}

\begin{table}[ht]
\centering
\footnotesize
\caption{So sánh mức tiêu thụ năng lượng và diện tích mạch bán dẫn giữa phép toán FP16 và cộng số nguyên INT8 trên tiến trình bán dẫn 7nm (Dữ liệu Horst et al., Horowitz 2014, Wang et al. 2024).}
\label{tab:bitnet-energy-comparison}
\begin{tabular}{lccc}
\toprule
\textbf{Phép toán số học} & \textbf{Năng lượng tiêu thụ (pJ)} & \textbf{Diện tích mạch ($\mu\text{m}^2$)} & \textbf{Mức tiết kiệm năng lượng} \\
\midrule
FP16 Multiplier (Nhân dấu phẩy động) & $1.10\text{ pJ}$ & $1640\,\mu\text{m}^2$ & Mốc cơ sở ($1\times$) \\
FP16 Adder (Cộng dấu phẩy động)      & $0.40\text{ pJ}$ & $1360\,\mu\text{m}^2$ & Tiết kiệm $2.75\times$ \\
INT32 Multiplier (Nhân số nguyên 32) & $3.10\text{ pJ}$ & $3200\,\mu\text{m}^2$ & Tốn hơn $2.8\times$ \\
\textbf{INT8 Adder (Cộng số nguyên 8)}& $\mathbf{0.03\text{ pJ}}$ & $\mathbf{36\,\mu\text{m}^2}$ & \textbf{Tiết kiệm hơn $36.7\times$!} \\
\bottomrule
\end{tabular}
\end{table}

\paragraph{Ý nghĩa kỹ nghệ và ứng dụng thực tế trên thiết bị biên}:
\begin{itemize}
  \item \textbf{Phá vỡ bức tường băng thông bộ nhớ (Memory Wall)}: Trong giai đoạn Auto-regressive Generation, tác vụ bị giới hạn bởi tốc độ đọc trọng số từ VRAM (Memory-bound). Nén trọng số từ 16-bit xuống 1.58-bit giúp cắt giảm $90\%$ lưu lượng truy xuất bộ nhớ, cho phép mô hình chạy với tốc độ suy luận nhanh gấp $2\times - 4\times$ trên cùng một phần cứng.
  \item \textbf{Kiến trúc Apple Silicon \& MLX}: Trên các chip M-series của Apple, các nhân CPU NEON và GPU Metal có thể thực hiện vector hóa các chỉ lệnh cộng trừ có dấu song song trên thanh ghi SIMD 128-bit/512-bit cực kỳ nhanh chóng mà không làm nóng máy hay tiêu hao pin.
  \item \textbf{Tiềm năng chip ASIC chuyên dụng không bộ nhân (Multiplier-less Silicon)}: Một chip gia tốc AI thiết kế riêng cho BitNet b1.58 không cần trang bị các khối phần cứng nhân ma trận cồng kềnh (Multiplier Array), tiết kiệm đến $95\%$ diện tích bán dẫn, cho phép tích hợp hàng chục nghìn bộ cộng tích lũy (Accumulator Trees) trên cùng một diện tích die silicon, mở đường cho việc nhúng LLM hàng chục tỷ tham số vào điện thoại thông minh, kính thực tế ảo và thiết bị IoT công nghiệp!
\end{itemize}

---

\section{Kỹ Nghệ Tinh Chỉnh Hiệu Quả Tham Số: LoRA, QLoRA \& DoRA}
\label{sec:peft-lora-qlora}

Khi muốn dạy mô hình học thêm một miền kiến thức mới hoặc định hình phong cách phản hồi chuyên biệt, phương pháp truyền thống là tinh chỉnh toàn phần (\textbf{Full Fine-Tuning} -- cập nhật toàn bộ tham số).

\subsection{Khủng Hoảng Bộ Nhớ Của Full Fine-Tuning}
Hãy làm một bài toán định cỡ bộ nhớ phần cứng để thấy sự bất khả thi của phương pháp này đối với mô hình 70B:
\begin{itemize}
  \item Trọng số mô hình tĩnh (FP16): $70 \times 10^9 \times 2\text{ Bytes} = 140\text{ GB}$.
  \item Gradients đạo hàm (FP16): $140\text{ GB}$.
  \item Trạng thái Optimizer AdamW (First moment $m_t$ và Second moment $v_t$ dạng FP32): $70 \times 10^9 \times (4 + 4)\text{ Bytes} = 560\text{ GB}$!
\end{itemize}
Tổng dung lượng VRAM tối thiểu để huấn luyện Full Fine-Tuning:
\begin{equation}
M_{\text{train}} = 140 + 140 + 560 = 840\text{ GB VRAM!}
\end{equation}
Chi phí này đòi hỏi một cụm ít nhất 12 GPU A100/H100 80GB kết nối NVLink đắt đỏ, vượt xa khả năng tiếp cận của phần lớn doanh nghiệp.

\subsection{Bản Chất Toán Học Của LoRA (Low-Rank Adaptation)}
Được Edward Hu và nhóm nghiên cứu Microsoft đề xuất năm 2021, \textbf{LoRA} dựa trên một giả thuyết giải tích nền tảng:

\begin{theorem}[Giả Thuyết Số Chiều Nội Tại Thấp -- Low Intrinsic Dimension]
Mặc dù các ma trận trọng số $\mathbf{W}_0 \in \mathbb{R}^{d \times k}$ của mô hình Transformer có số chiều biểu diễn khổng lồ, sự biến thiên trọng số $\Delta \mathbf{W}$ cần thiết để thích ứng với một tác vụ chuyên biệt thực chất chỉ nằm trên một không gian con có \textbf{hạng ma trận (Rank) rất nhỏ}:
\begin{equation}
\text{rank}(\Delta \mathbf{W}) = r \ll \min(d, k)
\end{equation}
\end{theorem}

Thay vì cập nhật trực tiếp ma trận kích thước lớn $\Delta \mathbf{W} \in \mathbb{R}^{d \times k}$, LoRA phân rã ma trận này thành tích của hai ma trận hạng thấp:
\begin{equation}
\mathbf{W} = \mathbf{W}_0 + \Delta \mathbf{W} = \mathbf{W}_0 + \frac{\alpha}{r} (\mathbf{A} \mathbf{B})
\end{equation}
Trong đó:
\begin{itemize}
  \item $\mathbf{W}_0 \in \mathbb{R}^{d \times k}$: Ma trận trọng số ban đầu, được \textbf{đóng băng hoàn toàn (Frozen)}, không tham gia tính đạo hàm hay cập nhật optimizer.
  \item $\mathbf{A} \in \mathbb{R}^{d \times r}$: Ma trận hạ chiều, được khởi tạo bằng phân phối Gauss chuẩn: $\mathcal{N}(0, \sigma^2)$.
  \item $\mathbf{B} \in \mathbb{R}^{r \times k}$: Ma trận tăng chiều, được \textbf{khởi tạo toàn bộ bằng số 0} ($\mathbf{B} = \mathbf{0}$).
  \item $r$: Siêu tham số Hạng (Rank), thường chọn $r \in \{4, 8, 16, 32, 64\}$.
  \item $\alpha$: Hệ số co giãn cố định (LoRA Alpha).
\end{itemize}

\begin{note}
Tại sao ma trận $\mathbf{B}$ bắt buộc phải được khởi tạo bằng $0$?
Bởi vì tại bước huấn luyện đầu tiên ($t=0$), ta có:
\begin{equation}
\Delta \mathbf{W} = \mathbf{A} \mathbf{B} = \mathbf{A} \times \mathbf{0} = \mathbf{0} \implies \mathbf{W} = \mathbf{W}_0
\end{equation}
Điều này đảm bảo rằng tại thời điểm bắt đầu tinh chỉnh, hành vi của mô hình hoàn toàn giống hệt mô hình gốc nguyên bản $100\%$, không bị nhiễu loạn ngẫu nhiên làm hỏng tri thức sẵn có!
\end{note}

\paragraph{Cơ chế hợp nhất trọng số không trễ suy luận (Zero-Latency Weight Merging)}:
Sau khi quá trình huấn luyện hoàn tất, chúng ta có thể cộng trực tiếp phần chênh lệch $\Delta \mathbf{W}$ vào ma trận trọng số gốc:
\begin{equation}
\mathbf{W}_{\text{merged}} = \mathbf{W}_0 + \frac{\alpha}{r} (\mathbf{A} \mathbf{B})
\end{equation}
Sau phép cộng ma trận tĩnh một lần duy nhất này, nhánh LoRA được giải phóng. Mô hình suy luận với cấu trúc mạng nguyên bản, đạt tốc độ suy luận đỉnh cao mà không tốn thêm bất kỳ phép tính rẽ nhánh nào!

\subsection{Đột Phá QLoRA (Quantized LoRA) \& Dẫn Xuất Phân Vị NF4}
\label{sec:qlora-nf4-math}

Năm 2023, Tim Dettmers và cộng sự công bố công trình đột phá \textbf{QLoRA (Quantized Low-Rank Adaptation)}, mở ra kỷ nguyên cho phép một cá nhân có thể tinh chỉnh mô hình ngôn ngữ 70 tỷ tham số trên một máy tính cá nhân trang bị 2 GPU tiêu dùng (RTX 3090/4090 24GB). Thành tựu này được xây dựng trên ba trụ cột toán học và kỹ nghệ hệ thống:

\paragraph{1. Dẫn Xuất Toán Học Của Kiểu Dữ Liệu Phân Vị NF4 (NormalFloat-4)}:
Trong các thuật toán lượng tử hóa số nguyên truyền thống (INT4), các mốc lượng tử hóa được phân bố đều nhau trên đoạn $[-a, a]$. Tuy nhiên, theo các nghiên cứu lý thuyết thông tin, phân phối trọng số của mô hình ngôn ngữ sau tiền huấn luyện tuân theo phân phối chuẩn Gauss $\mathcal{N}(0, \sigma^2)$ với mật độ tập trung cực cao quanh điểm 0 và thưa dần về hai phía đuôi. Lượng tử hóa đều INT4 lãng phí rất nhiều mức rời rạc ở vùng đuôi thưa và thiếu hụt trầm trọng độ phân giải ở vùng đỉnh chuông có mật độ xác suất dày đặc.

Để đạt \textbf{Entropi thông tin tối đa (Maximum Entropy Quantization)}, kiểu dữ liệu \textbf{NormalFloat-4 (NF4)} được xây dựng dựa trên nguyên lý \term{Equal-Quantile Information (Lượng thông tin phân vị đều)}: Chia miền giá trị của biến ngẫu nhiên Gauss chuẩn hóa $Z \sim \mathcal{N}(0, 1)$ thành $2^k = 16$ khoảng rời rạc sao cho xác suất tích lũy trong mỗi khoảng là hoàn toàn bằng nhau:
\begin{equation}
P(q_i \le Z \le q_{i+1}) = \frac{1}{16} = 0.0625, \qquad \forall i \in \{1, 2, \dots, 16\}
\end{equation}

Cho $\Phi(z) = \frac{1}{\sqrt{2\pi}} \int_{-\infty}^z e^{-t^2/2} dt$ là hàm phân phối tích lũy (CDF) của phân phối chuẩn tắc. Các điểm phân vị tối ưu được xác định tường minh qua hàm phân vị nghịch đảo (\term{Quantile / Probit Function}):
\begin{equation}
\label{eq:nf4-quantile-formula}
q_i = \Phi^{-1}\left( \frac{i - 0.5}{16} \right), \qquad i \in \{1, 2, \dots, 16\}
\end{equation}
Để đảm bảo giá trị số thực $0.0$ được biểu diễn chính xác tuyệt đối mà không có sai số làm tròn (bảo toàn tính thưa và giá trị ngắt kích hoạt), bảng mã 16 mốc của NF4 được chuẩn hóa bất đối xứng thành hai nhánh âm và dương:
\begin{itemize}
  \item 7 mức phân vị cho nửa miền âm $[-1, 0)$ sao cho $\Phi(q_i) = \frac{i}{16}$.
  \item 1 mức chính xác tại $0.0$.
  \item 8 mức phân vị cho nửa miền dương $(0, 1]$ sao cho $\Phi(q_i) = \frac{i - 0.5}{16}$.
\end{itemize}
Sau đó, toàn bộ 16 điểm mốc được chia tỷ lệ để dải giá trị nằm trọn vẹn trong đoạn đối xứng $[-1, 1]$. Bảng mã NF4 chính thức trở thành kiểu dữ liệu nén tối ưu thông tin lý thuyết tuyệt đối cho trọng số mạng nơ-ron sâu!

\paragraph{2. Lượng Tử Hóa Kép (Double Quantization -- DQ)}:
Trong lượng tử hóa khối (Block-wise Quantization), để giảm thiểu sai số ngoại lai, ma trận trọng số được chia thành các khối nhỏ gồm $B_1 = 64$ phần tử. Mỗi khối đòi hỏi một hệ số co giãn riêng $c_1 \in \mathbb{R}$.
Nếu lưu $c_1$ ở định dạng số thực chuẩn FP32 ($32\text{ bits}$):
\begin{equation}
\text{Dung lượng hệ số co giãn đơn} = \frac{32\text{ bits}}{64\text{ tham số}} = 0.50\text{ bits/tham số}
\end{equation}
Mức hao phí $0.50\text{ bit}$ chiếm tới $12.5\%$ tổng dung lượng lưu trữ của trọng số 4-bit!

\textbf{Lượng tử hóa kép (Double Quantization)} giải quyết bài toán này bằng cách gom các hệ số co giãn $c_1$ thành các khối lớn hơn với kích thước $B_2 = 256$, sau đó lượng tử hóa chính các hệ số $c_1$ từ FP32 xuống số nguyên 8-bit (FP8) với hệ số co giãn cấp hai $c_2$ (FP32).
Dung lượng bộ nhớ thực tế cho mỗi tham số sau khi áp dụng Double Quantization:
\begin{align}
\text{Dung lượng scale DQ} &= \frac{8\text{ bits}}{B_1} + \frac{32\text{ bits}}{B_1 \times B_2} = \frac{8}{64} + \frac{32}{64 \times 256} \notag \\
&= 0.125 + \frac{32}{16384} \approx 0.125 + 0.001953 = \mathbf{0.126953\text{ bits/tham số}}
\end{align}
Mức tiết kiệm bộ nhớ trên từng tham số là:
\begin{equation}
\Delta b = 0.50 - 0.126953 = \mathbf{0.373047\text{ bits/tham số}}
\end{equation}
Với mô hình Llama-3-70B ($70 \times 10^9$ tham số), lượng tử hóa kép giúp cắt giảm trực tiếp:
\begin{equation}
\Delta M = \frac{70 \times 10^9 \times 0.373047\text{ bits}}{8\text{ bits/byte} \times 1024^3\text{ bytes/GiB}} \approx \mathbf{3.04\text{ GiB VRAM!}}
\end{equation}
Mức tiết kiệm $3.04\text{ GiB}$ này mang tính chất sống còn, quyết định xem một mô hình 70B có thể chạy vừa khít trong giới hạn $48\text{ GB}$ VRAM của 2 card đồ họa RTX 3090/4090 hay không!

\paragraph{3. Cơ Chế Paged Optimizers Chống Sập Bộ Nhớ (OOM)}:
Trong quá trình tinh chỉnh, khi gặp các mẫu dữ liệu có độ dài chuỗi (Sequence Length) đột biến, bộ nhớ đệm kích hoạt (Activation Memory) tăng vọt cục bộ, dễ dẫn đến hiện tượng sập bộ nhớ (\term{CUDA Out-Of-Memory}).
QLoRA tích hợp tính năng \textbf{Paged Optimizers} dựa trên công nghệ bộ nhớ ảo phân trang thống nhất của NVIDIA (CUDA Unified Memory Page-locked Paging):
Khi VRAM chạm ngưỡng tới hạn ($>95\%$), hệ thống tự động hoán chuyển (Page-out) các trạng thái của bộ tối ưu AdamW (First Moments $\mathbf{m}_t$ và Second Moments $\mathbf{v}_t$) sang bộ nhớ RAM hệ thống thông qua bus PCIe tốc độ cao, giải phóng không gian cho tính toán kích hoạt. Sau khi hoàn thành lan truyền ngược, các trang dữ liệu được kéo ngược trở lại (Page-in) GPU để cập nhật trọng số. Kỹ thuật này triệt tiêu hoàn toàn nguy cơ sập luồng trong các đợt huấn luyện dài ngày!

---

\subsection{DoRA: Tinh Chỉnh Phân Rã Độ Lớn \& Hướng Trọng Số (Weight-Decomposed Low-Rank Adaptation)}
\label{sec:dora-math-decomposition}

Mặc dù LoRA đã trở thành tiêu chuẩn công nghiệp cho tinh chỉnh hiệu quả tham số, các công trình nghiên cứu sâu về động học tối ưu nhận thấy LoRA vẫn có khoảng cách hiệu năng nhất định so với tinh chỉnh toàn phần (Full Fine-Tuning) trên các tác vụ đòi hỏi tư duy logic cao như toán học và lập trình (GSM8K, MATH, HumanEval).
Năm 2024, Shih-Yang Liu và nhóm nghiên cứu NVIDIA/UT Austin công bố \textbf{DoRA (Weight-Decomposed Low-Rank Adaptation)}, giải mã chính xác nguyên nhân gốc rễ của khoảng cách này và đề xuất giải pháp đột phá.

\paragraph{1. Phân Tích Động Học Trọng Số: Độ Lớn (Magnitude) \& Hướng (Direction)}:
Bất kỳ ma trận trọng số nào $\mathbf{W} \in \mathbb{R}^{d \times k}$ đều có thể phân rã thành vector độ lớn và ma trận hướng chuẩn hóa theo từng cột (Column-wise L2 Norm):
\begin{equation}
\mathbf{W} = \mathbf{m} \odot \frac{\mathbf{V}}{\|\mathbf{V}\|_c}
\end{equation}
Trong đó:
\begin{itemize}
  \item $\mathbf{m} \in \mathbb{R}^{1 \times k}$ là vector chứa độ lớn Euclidean của từng cột: $m_j = \|\mathbf{w}_{:, j}\|_2 = \sqrt{\sum_{i=1}^d W_{i, j}^2}$.
  \item $\mathbf{V} \in \mathbb{R}^{d \times k}$ là ma trận hướng. Ma trận $\frac{\mathbf{V}}{\|\mathbf{V}\|_c}$ chứa các cột có chuẩn đơn vị ($\|\mathbf{v}_{:, j} / \|\mathbf{v}_{:, j}\|_2\|_2 = 1$).
  \item Ký hiệu $\odot$ đại diện cho phép nhân phát thanh theo từng cột.
\end{itemize}

\begin{figure}[ht]
\centering
\begin{tikzpicture}[scale=0.9, every node/.style={scale=0.9}]
  % Axes
  \draw[->, thick, color=gray!60] (-0.5, 0) -- (5.5, 0) node[right] {\footnotesize $w_1$};
  \draw[->, thick, color=gray!60] (0, -0.5) -- (0, 4.5) node[above] {\footnotesize $w_2$};
  
  % Base weight W0
  \draw[->, ultra thick, color=primarycolor!90] (0, 0) -- (3.0, 2.0) node[midway, above left] {$\mathbf{w}_0$};
  \draw[dashed, color=primarycolor!50] (0,0) circle (3.605);
  \node[color=primarycolor!90, font=\scriptsize] at (3.2, 1.6) {$\|\mathbf{w}_0\|_2 = m_0$};
  
  % LoRA update
  \draw[->, thick, color=dangercolor!80, dashed] (3.0, 2.0) -- (4.2, 3.2) node[midway, right] {\scriptsize $\Delta \mathbf{w} = \frac{\alpha}{r}\mathbf{A}\mathbf{B}$};
  \draw[->, ultra thick, color=dangercolor!90] (0, 0) -- (4.2, 3.2) node[above right] {\small $\mathbf{w}_{\text{LoRA}}$};
  
  % DoRA update: Decoupled
  \draw[->, thick, color=secondarycolor!80, dashed] (3.0, 2.0) -- (2.0, 3.0) node[midway, left] {\scriptsize Xoay hướng $\mathbf{V}$};
  \draw[->, ultra thick, color=secondarycolor!90] (0, 0) -- (2.6, 3.9) node[above] {\small $\mathbf{w}_{\text{DoRA}} = m \frac{\mathbf{V}}{\|\mathbf{V}\|}$};
  \node[draw, fill=secondarycolor!10, rounded corners=2pt, font=\scriptsize, align=center] at (4.5, 0.8) {DoRA tách biệt hoàn toàn\\Độ lớn $m$ và Hướng $\mathbf{V}$};
\end{tikzpicture}
\caption{Minh họa hình học sự khác biệt giữa LoRA và DoRA: LoRA buộc độ lớn và hướng thay đổi tỷ lệ thuận; DoRA cho phép xoay góc biểu diễn độc lập với việc điều chỉnh độ dài véc-tơ.}
\label{fig:dora-geometric-intuition}
\end{figure}

Phân tích thống kê về mức độ biến thiên giữa mô hình gốc $\mathbf{W}_0$ và mô hình sau tinh chỉnh $\mathbf{W}$:
\begin{itemize}
  \item \textbf{Trong Full Fine-Tuning (FT)}: Mức thay đổi độ lớn $\Delta m$ và mức thay đổi góc hướng $\Delta \theta$ có hệ số tương quan rất yếu hoặc hơi âm. Mô hình có thể xoay góc biểu diễn rất mạnh mà hầu như giữ nguyên độ lớn của nơ-ron, hoặc ngược lại.
  \item \textbf{Trong LoRA truyền thống}: Vì $\mathbf{W} = \mathbf{W}_0 + \frac{\alpha}{r} \mathbf{A}\mathbf{B}$, mọi cập nhật lên ma trận $\mathbf{A}, \mathbf{B}$ đều đồng thời tác động lên cả hướng và độ lớn theo một tỷ lệ cố định. Mức tương quan giữa $\Delta m$ và $\Delta \theta$ luôn luôn dương và rất chặt chẽ, tạo ra hiện tượng ``ghép đôi cứng nhắc'' (Coupled Dynamics), khiến LoRA khó lòng mô phỏng các biến đổi phức tạp của Full Fine-Tuning!
\end{itemize}

\paragraph{2. Công Thức Toán Học Của DoRA}:
DoRA giải quyết triệt để vấn đề này bằng cách \textbf{chỉ áp dụng LoRA để điều chỉnh ma trận hướng $\mathbf{V}$}, trong khi độ lớn $\mathbf{m}$ được tối ưu hóa độc lập như một vector tham số riêng:
\begin{equation}
\label{eq:dora-forward-formulation}
\mathbf{W} = \mathbf{m} \odot \frac{\mathbf{V}}{\|\mathbf{V}\|_c} = \mathbf{m} \odot \frac{\mathbf{W}_0 + \frac{\alpha}{r} \mathbf{A} \mathbf{B}}{\|\mathbf{W}_0 + \frac{\alpha}{r} \mathbf{A} \mathbf{B}\|_c}
\end{equation}
Trong đó:
\begin{itemize}
  \item $\mathbf{W}_0 \in \mathbb{R}^{d \times k}$: Trọng số cơ sở được đóng băng (Frozen).
  \item $\mathbf{m} \in \mathbb{R}^{1 \times k}$: Vector độ lớn khả huấn (Trainable Magnitude), được khởi tạo chính xác bằng chuẩn cột ban đầu: $\mathbf{m} = \|\mathbf{W}_0\|_c$.
  \item $\mathbf{A} \in \mathbb{R}^{d \times r}$ và $\mathbf{B} \in \mathbb{R}^{r \times k}$: Hai ma trận adapter hạng thấp, với $\mathbf{B} = \mathbf{0}$ tại $t=0$.
\end{itemize}
Tại bước huấn luyện đầu tiên ($t=0$):
\begin{equation}
\mathbf{W} = \|\mathbf{W}_0\|_c \odot \frac{\mathbf{W}_0 + \mathbf{0}}{\|\mathbf{W}_0 + \mathbf{0}\|_c} = \|\mathbf{W}_0\|_c \odot \frac{\mathbf{W}_0}{\|\mathbf{W}_0\|_c} = \mathbf{W}_0
\end{equation}
Đặc tính bảo toàn hành vi ban đầu tại bước khởi đầu được giữ nguyên vẹn $100\%$!

\paragraph{3. Dẫn Xuất Giải Tích Gradient: Cơ Chế Chiếu Trực Giao (Orthogonal Projection)}:
Hãy phân tích gradient của hàm mất mát $\mathcal{L}$ đối với vector độ lớn $\mathbf{m}$ và ma trận hướng $\mathbf{V}$ để hiểu tại sao DoRA lại ổn định và vượt trội:

\textbf{Gradient theo độ lớn $\mathbf{m}$}:
\begin{equation}
\nabla_{\mathbf{m}} \mathcal{L} = \sum_{i=1}^d \left( \nabla_{\mathbf{W}} \mathcal{L} \right)_{i, :} \odot \left( \frac{\mathbf{V}}{\|\mathbf{V}\|_c} \right)_{i, :}
\end{equation}
Gradient của $m_j$ chính là tích vô hướng giữa gradient trọng số và vector hướng đơn vị của cột đó, đo lường sự cần thiết phải tăng hoặc giảm cường độ phát xung của nơ-ron.

\textbf{Gradient theo ma trận hướng $\mathbf{V}$}:
Áp dụng quy tắc đạo hàm hàm hợp ma trận:
\begin{equation}
\nabla_{\mathbf{V}} \mathcal{L} = \frac{\mathbf{m}}{\|\mathbf{V}\|_c} \odot \left[ \nabla_{\mathbf{W}} \mathcal{L} - \frac{\mathbf{V}}{\|\mathbf{V}\|_c^2} \odot \sum_{i=1}^d \left( \mathbf{V} \odot \nabla_{\mathbf{W}} \mathcal{L} \right) \right]
\end{equation}
Viết dưới dạng toán tử ma trận cho từng cột vector $\mathbf{v}$ và đạo hàm $\mathbf{g} = \nabla_{\mathbf{w}} \mathcal{L}$:
\begin{equation}
\label{eq:dora-orthogonal-projection}
\nabla_{\mathbf{v}} \mathcal{L} = \frac{m}{\|\mathbf{v}\|_2} \left( \mathbf{I} - \frac{\mathbf{v} \mathbf{v}^T}{\|\mathbf{v}\|_2^2} \right) \mathbf{g}
\end{equation}
Ma trận $\mathbf{P}_{\mathbf{v}} = \left( \mathbf{I} - \frac{\mathbf{v} \mathbf{v}^T}{\|\mathbf{v}\|_2^2} \right)$ chính là \textbf{Ma trận chiếu trực giao (Orthogonal Projection Matrix)} lên không gian phần bù trực giao của vector $\mathbf{v}$!
Tính chất sống còn của toán tử này:
\begin{equation}
\mathbf{v}^T \left( \mathbf{I} - \frac{\mathbf{v} \mathbf{v}^T}{\|\mathbf{v}\|_2^2} \right) = \mathbf{v}^T - \frac{\mathbf{v}^T \mathbf{v}}{\|\mathbf{v}\|_2^2} \mathbf{v}^T = \mathbf{v}^T - \mathbf{v}^T = \mathbf{0}^T
\end{equation}
\textbf{Hệ quả toán học then chốt}: Gradient cập nhật cho nhánh hướng $\mathbf{V}$ luôn luôn vuông góc $100\%$ với vector hướng hiện tại!
Điều này đồng nghĩa với việc nhánh LoRA $(\mathbf{A}, \mathbf{B})$ trong DoRA được giải phóng hoàn toàn khỏi nhiệm vụ điều chỉnh độ lớn: nó \textbf{chỉ thuần túy xoay góc của không gian biểu diễn}. Việc thay đổi độ lớn được bàn giao toàn quyền cho vector tham số $\mathbf{m}$ một chiều siêu nhẹ.

\paragraph{4. Hợp Nhất Trọng Số Zero-Overhead Khi Suy Luận}:
Giống như LoRA, sau khi quá trình huấn luyện kết thúc, toàn bộ các tham số của DoRA được gộp tĩnh thành một ma trận duy nhất trước khi triển khai:
\begin{equation}
\mathbf{W}_{\text{deployed}} = \mathbf{m} \odot \frac{\mathbf{W}_0 + \frac{\alpha}{r} \mathbf{A} \mathbf{B}}{\|\mathbf{W}_0 + \frac{\alpha}{r} \mathbf{A} \mathbf{B}\|_c} \in \mathbb{R}^{d \times k}
\end{equation}
Mô hình serving chỉ cần tải ma trận $\mathbf{W}_{\text{deployed}}$, đạt tốc độ suy luận nguyên bản với chi phí trễ bằng 0!

---

\section{Căn Chỉnh Sở Thích Con Người: Từ RLHF Đến DPO, ORPO \& SimPO}
\label{sec:human-alignment}

Sau khi một mô hình được tiền huấn luyện (Pre-training) và tinh chỉnh có giám sát theo lệnh (Supervised Fine-Tuning -- SFT), nó đã có khả năng nói năng lưu loát và tuân thủ định dạng. Tuy nhiên, SFT đơn thuần mắc phải các hội chứng nguy hiểm:
\begin{itemize}
  \item \textbf{Nịnh người dùng (Sycophancy)}: Mô hình đồng tình với các tiền đề sai của người hỏi ngay cả khi chúng trái ngược với chân lý khoa học.
  \item \textbf{Ảo giác tự tin (Confident Hallucination)}: Bịa đặt các sự kiện lịch sử hoặc hàm thư viện không có thật với giọng điệu quả quyết.
  \item \textbf{Không an toàn (Safety Vulnerability)}: Dễ bị lợi dụng để sinh mã độc hoặc hướng dẫn hành vi nguy hại.
\end{itemize}
Quá trình căn chỉnh sở thích con người (Human Alignment) giải quyết bài toán này bằng cách dạy mô hình ưu tiên các câu trả lời hữu ích, trung thực và vô hại (\textit{Helpful, Honest, Harmless -- 3H}).

\subsection{Học Tăng Cường Từ Phản Hồi Con Người (RLHF với PPO)}
Khung làm việc RLHF truyền thống (Christiano et al., Ouyang et al. InstructGPT) tiến hành qua ba giai đoạn:

\paragraph{Giai đoạn 1: Thu thập dữ liệu sở thích cặp $(x, y_w, y_l)$}:
Với một câu hỏi $x$, mô hình sinh ra hai câu trả lời ứng viên $y_1, y_2$. Con người (hoặc mô hình giám định AI) đánh giá và chỉ định câu trả lời được ưa chuộng hơn ($y_w$, winner) và câu trả lời kém hơn ($y_l$, loser).

\paragraph{Giai đoạn 2: Huấn luyện Mô hình Điểm Thưởng (Reward Model)}:
Dựa trên mô hình xác suất lựa chọn Bradley-Terry (1952), xác suất câu trả lời $y_w$ được ưa chuộng hơn $y_l$ với đầu vào $x$ được mô hình hóa qua hàm sigmoid:
\begin{equation}
P(y_w \succ y_l \mid x) = \sigma\left(r_\psi(x, y_w) - r_\psi(x, y_l)\right) = \frac{1}{1 + e^{-(r_\psi(x, y_w) - r_\psi(x, y_l))}}
\end{equation}
Hàm mất mát để huấn luyện mô hình điểm thưởng $r_\psi$ bằng phương pháp cực đại hóa hợp lý (Maximum Likelihood):
\begin{equation}
\mathcal{L}_R(\psi) = -\mathbb{E}_{(x, y_w, y_l) \sim \mathcal{D}} \left[ \log \sigma\left(r_\psi(x, y_w) - r_\psi(x, y_l)\right) \right]
\end{equation}

\paragraph{Giai đoạn 3: Tối ưu hóa chính sách với PPO và phạt Phân kỳ KL}:
Mục tiêu là huấn luyện chính sách mô hình ngôn ngữ $\pi_\theta$ để tối đa hóa điểm thưởng dự đoán, đồng thời áp đặt một hình phạt dựa trên độ phân kỳ Kullback-Leibler ($D_{\text{KL}}$) so với mô hình tham chiếu ban đầu $\pi_{\text{ref}}$ (chính là mô hình SFT đã đóng băng):
\begin{equation}
\max_{\theta} \mathbb{E}_{x \sim \mathcal{D}, y \sim \pi_\theta} \left[ r_\psi(x, y) - \beta D_{\text{KL}}\left(\pi_\theta(y \mid x) \,\|\, \pi_{\text{ref}}(y \mid x)\right) \right]
\end{equation}
Trong đó $\beta$ là hệ số phạt KL. Hình phạt này đóng vai trò sống còn: nếu không có nó, mô hình chính sách $\pi_\theta$ sẽ nhanh chóng tìm ra các lỗ hổng của mô hình điểm thưởng để sinh ra các đoạn văn bản vô nghĩa nhưng đạt điểm thưởng cực cao (hiện tượng \textit{Reward Hacking} hoặc \textit{Overoptimization}).

\paragraph{Gánh nặng kiến trúc của RLHF}:
Để chạy RLHF với thuật toán PPO, hệ thống cần duy trì đồng thời \textbf{4 mô hình lớn} trên GPU:
\begin{enumerate}
  \item Mô hình Chính sách đang huấn luyện (Actor: $\pi_\theta$).
  \item Mô hình Đánh giá giá trị trạng thái (Critic: $V_\phi$).
  \item Mô hình Tham chiếu cố định (Reference Model: $\pi_{\text{ref}}$).
  \item Mô hình Điểm thưởng cố định (Reward Model: $r_\psi$).
\end{enumerate}
Sự phối hợp phức tạp giữa 4 mô hình này khiến quá trình huấn luyện RLHF cực kỳ tốn VRAM, kém ổn định và khó hội tụ.

---

\subsection{DPO (Direct Preference Optimization): Chứng Minh Giải Tích Đột Phá}
Năm 2023, Rafailov et al. (Stanford) công bố công trình mang tính cách mạng \textbf{DPO}, chứng minh rằng chúng ta có thể tối ưu hóa chính sách trực tiếp từ dữ liệu sở thích con người \textbf{mà không cần huấn luyện Reward Model riêng biệt và không cần vòng lặp Reinforcement Learning phức tạp!}

Dưới đây là chứng minh giải tích toàn diện từ bài toán tối ưu ràng buộc:

\begin{theorem}[Nghiệm Tối Ưu Giải Tích Của Bài Toán RLHF Khép Kín]
Xét bài toán tối ưu hóa có phạt KL cho một câu hỏi $x$:
\begin{equation}
\max_{\pi} \mathbb{E}_{y \sim \pi}\left[ r(x, y) \right] - \beta D_{\text{KL}}\left(\pi(y \mid x) \,\|\, \pi_{\text{ref}}(y \mid x)\right)
\end{equation}
với ràng buộc phân phối xác suất hợp lệ $\sum_y \pi(y \mid x) = 1$.
\end{theorem}

\begin{proof}
Khai triển định nghĩa của độ phân kỳ KL:
\begin{equation}
D_{\text{KL}}(\pi \,\|\, \pi_{\text{ref}}) = \sum_y \pi(y \mid x) \log \frac{\pi(y \mid x)}{\pi_{\text{ref}}(y \mid x)}
\end{equation}
Biểu thức hàm mục tiêu trở thành:
\begin{align*}
\mathcal{F}(\pi) &= \sum_y \pi(y \mid x) r(x, y) - \beta \sum_y \pi(y \mid x) \log \frac{\pi(y \mid x)}{\pi_{\text{ref}}(y \mid x)} \\
&= -\beta \sum_y \pi(y \mid x) \left[ \log \frac{\pi(y \mid x)}{\pi_{\text{ref}}(y \mid x)} - \frac{1}{\beta} r(x, y) \right] \\
&= -\beta \sum_y \pi(y \mid x) \log \left( \frac{\pi(y \mid x)}{\pi_{\text{ref}}(y \mid x) \exp\left(\frac{1}{\beta} r(x, y)\right)} \right)
\end{align*}
Đặt hàm phân hoạch chuẩn hóa:
\begin{equation}
Z(x) = \sum_y \pi_{\text{ref}}(y \mid x) \exp\left(\frac{1}{\beta} r(x, y)\right)
\end{equation}
Và định nghĩa phân phối xác suất tối ưu chuẩn hóa:
\begin{equation}
\pi^*(y \mid x) = \frac{1}{Z(x)} \pi_{\text{ref}}(y \mid x) \exp\left(\frac{1}{\beta} r(x, y)\right)
\end{equation}
Thay $Z(x) \pi^*(y \mid x)$ vào biểu thức ta có:
\begin{align*}
\mathcal{F}(\pi) &= -\beta \sum_y \pi(y \mid x) \log \left( \frac{\pi(y \mid x)}{Z(x) \pi^*(y \mid x)} \right) \\
&= -\beta \sum_y \pi(y \mid x) \log \frac{\pi(y \mid x)}{\pi^*(y \mid x)} + \beta \log Z(x) \sum_y \pi(y \mid x) \\
&= -\beta D_{\text{KL}}\left(\pi(y \mid x) \,\|\, \pi^*(y \mid x)\right) + \beta \log Z(x)
\end{align*}
Vì $\beta > 0$ và độ phân kỳ Kullback-Leibler $D_{\text{KL}}(\pi \,\|\, \pi^*) \ge 0$, giá trị cực đại của $\mathcal{F}(\pi)$ đạt được khi và chỉ khi:
\begin{equation}
D_{\text{KL}}\left(\pi(y \mid x) \,\|\, \pi^*(y \mid x)\right) = 0 \iff \pi(y \mid x) = \pi^*(y \mid x)
\end{equation}
Do đó, nghiệm tối ưu giải tích chính xác duy nhất là:
\begin{equation}
\label{eq:dpo-optimal-policy}
\pi^*(y \mid x) = \frac{1}{Z(x)} \pi_{\text{ref}}(y \mid x) \exp\left(\frac{1}{\beta} r(x, y)\right)
\end{equation}
\end{proof}

\paragraph{Bước biến đổi thiên tài: Rút hàm thưởng tiềm ẩn (Implicit Reward)}:
Từ công thức nghiệm tối ưu (\ref{eq:dpo-optimal-policy}), lấy logarit tự nhiên hai vế:
\begin{equation}
\log \pi^*(y \mid x) = \log \pi_{\text{ref}}(y \mid x) + \frac{1}{\beta} r(x, y) - \log Z(x)
\end{equation}
Chuyển vế để biểu diễn hàm thưởng $r(x, y)$ trực tiếp qua phân phối xác suất:
\begin{equation}
\label{eq:implicit-reward}
r(x, y) = \beta \log \frac{\pi^*(y \mid x)}{\pi_{\text{ref}}(y \mid x)} + \beta \log Z(x)
\end{equation}
Bây giờ, thay biểu thức hàm thưởng tiềm ẩn (\ref{eq:implicit-reward}) vào mô hình xác suất lựa chọn Bradley-Terry cho cặp sở thích $(y_w, y_l)$:
\begin{align*}
r(x, y_w) - r(x, y_l) &= \left[\beta \log \frac{\pi^*(y_w \mid x)}{\pi_{\text{ref}}(y_w \mid x)} + \beta \log Z(x)\right] - \left[\beta \log \frac{\pi^*(y_l \mid x)}{\pi_{\text{ref}}(y_l \mid x)} + \beta \log Z(x)\right] \\
&= \beta \log \frac{\pi^*(y_w \mid x)}{\pi_{\text{ref}}(y_w \mid x)} - \beta \log \frac{\pi^*(y_l \mid x)}{\pi_{\text{ref}}(y_l \mid x)}
\end{align*}
\textbf{Hàm phân hoạch phức tạp $Z(x)$ đã bị triệt tiêu hoàn toàn!}

Thay trực tiếp hiệu số này vào hàm mất mát cực đại hóa hợp lý, ta thu được hàm mất mát tối thượng của \textbf{DPO}:
\begin{equation}
\label{eq:dpo-loss}
\mathcal{L}_{\text{DPO}}(\theta; \pi_{\text{ref}}) = -\mathbb{E}_{(x, y_w, y_l) \sim \mathcal{D}} \left[ \log \sigma\left( \beta \log \frac{\pi_\theta(y_w \mid x)}{\pi_{\text{ref}}(y_w \mid x)} - \beta \log \frac{\pi_\theta(y_l \mid x)}{\pi_{\text{ref}}(y_l \mid x)} \right) \right]
\end{equation}

\paragraph{Giải phẫu cơ chế Gradient của DPO}:
Lấy đạo hàm của $\mathcal{L}_{\text{DPO}}$ theo trọng số $\theta$:
\begin{equation}
\nabla_\theta \mathcal{L}_{\text{DPO}} = -\beta \cdot \sigma\left(\hat{r}_\theta(x, y_l) - \hat{r}_\theta(x, y_w)\right) \cdot \left[ \nabla_\theta \log \pi_\theta(y_w \mid x) - \nabla_\theta \log \pi_\theta(y_l \mid x) \right]
\end{equation}
Trong đó $\hat{r}_\theta(x, y) = \beta \log \frac{\pi_\theta(y \mid x)}{\pi_{\text{ref}}(y \mid x)}$ là hàm thưởng tiềm ẩn. Công thức này cho thấy cơ chế tự điều tiết tuyệt đẹp:
\begin{itemize}
  \item Thừa số thứ hai $\left[ \nabla_\theta \log \pi_\theta(y_w \mid x) - \nabla_\theta \log \pi_\theta(y_l \mid x) \right]$: Tăng xác suất sinh ra câu trả lời thắng cuộc $y_w$ và giảm xác suất sinh ra câu trả lời thất bại $y_l$.
  \item Hệ số trọng số động $\sigma\left(\hat{r}_\theta(x, y_l) - \hat{r}_\theta(x, y_w)\right)$: Khi mô hình hiện tại đánh giá sai (tức $\hat{r}_\theta(y_l) > \hat{r}_\theta(y_w)$), hệ số này tiến gần về $1.0$, thúc đẩy cập nhật gradient cực mạnh. Khi mô hình đã phân biệt đúng với khoảng cách lớn, hệ số này tiệm cận $0.0$, ngăn chặn hiện tượng quá khớp!
\end{itemize}

---

\subsection{ORPO \& SimPO: Căn Chỉnh Tinh Gọn Không Cần Reference Model}
Mặc dù DPO loại bỏ được Reward Model và PPO, nó vẫn đòi hỏi duy trì một bản sao đóng băng của mô hình tham chiếu $\pi_{\text{ref}}$ chạy song song trên GPU VRAM để tính log-ratios. Hai kỹ thuật mới nhất đã loại bỏ nốt rào cản này:

\paragraph{1. ORPO (Odds Ratio Preference Optimization -- Hong et al., 2024)}:
ORPO định nghĩa tỷ số Odds của việc sinh chuỗi token $y$ dưới điều kiện ngữ cảnh $x$:
\begin{equation}
\text{odds}_\theta(y \mid x) = \frac{P_\theta(y \mid x)}{1 - P_\theta(y \mid x)}
\end{equation}
Hàm mất mát của ORPO kết hợp giữa mục tiêu SFT truyền thống và hàm phạt tỷ số Odds mà hoàn toàn không cần Reference Model:
\begin{equation}
\mathcal{L}_{\text{ORPO}} = \mathcal{L}_{\text{SFT}} + \lambda \cdot \mathcal{L}_{\text{OR}}
\end{equation}
\begin{equation}
\mathcal{L}_{\text{OR}} = -\mathbb{E}_{(x, y_w, y_l)} \left[ \log \sigma\left( \log \frac{\text{odds}_\theta(y_w \mid x)}{\text{odds}_\theta(y_l \mid x)} \right) \right]
\end{equation}
Nhờ đó, ORPO hoàn thành cả hai mục tiêu thích ứng lệnh và căn chỉnh sở thích trong đúng một lượt huấn luyện duy nhất!

\paragraph{2. SimPO (Simple Preference Optimization -- Meng et al., 2024)}:
SimPO tối giản hóa hơn nữa bằng cách chuẩn hóa log-xác suất theo chiều dài chuỗi token $|y|$ và áp đặt một lề mục tiêu $\gamma > 0$:
\begin{equation}
\mathcal{L}_{\text{SimPO}} = -\mathbb{E}_{(x, y_w, y_l)} \left[ \log \sigma\left( \frac{\beta}{|y_w|} \log \pi_\theta(y_w \mid x) - \frac{\beta}{|y_l|} \log \pi_\theta(y_l \mid x) - \gamma \right) \right]
\end{equation}
Việc chuẩn hóa chiều dài ngăn ngừa hiện tượng thiên vị các câu trả lời dài dòng (Verbosity Bias) của mô hình.

---

\subsection{Thuật Toán GRPO (Group Relative Policy Optimization) \& Kỷ Nguyên Reasoning}
\label{sec:grpo-deepseek-r1}

Đầu năm 2025, thế giới trí tuệ nhân tạo chứng kiến một cơn địa chấn công nghệ với sự ra đời của \textbf{DeepSeek-R1} và mô hình tiền thân \textbf{DeepSeek-R1-Zero} (DeepSeek-AI, 2025). Lần đầu tiên, một mô hình mã nguồn mở đạt năng lực suy luận toán học (AIME, MATH-500) và lập trình thuật toán (Codeforces, SWE-bench) sánh ngang với OpenAI o1 độc quyền, mở ra kỷ nguyên của \term{Reasoning Models} (mô hình tư duy Chuỗi suy luận mở rộng -- Extended Chain-of-Thought / System 2 Thinking).

Động cơ toán học trung tâm tạo nên kỳ tích này không phải là DPO hay PPO truyền thống, mà là thuật toán \textbf{Group Relative Policy Optimization (GRPO)} được kế thừa và phát triển từ DeepSeekMath (Shao et al., 2024).

\paragraph{1. Nút Thắt Trầm Kha Của PPO Trong Huấn Luyện Suy Luận Dài}:
Trong khung làm việc PPO kinh điển (Schulman et al., 2017) được dùng trong RLHF:
\begin{itemize}
  \item Để ước lượng độ lợi thế (Advantage $A_t$) theo phương pháp GAE (Generalized Advantage Estimation), hệ thống bắt buộc phải duy trì một mạng Đánh giá giá trị (Value Network / Critic $V_\phi(s)$) có số lượng tham số tương đương với mạng Chính sách (Policy / Actor $\pi_\theta$).
  \item Đối với các mô hình kích thước khổng lồ như DeepSeek 671B MoE hoặc Llama-3-70B, việc nạp thêm mạng Critic $V_\phi$ làm \textbf{tăng gấp đôi chi phí bộ nhớ VRAM} cho huấn luyện RL.
  \item Nguy hiểm hơn, trong các bài toán suy luận phức tạp, chuỗi suy luận có thể kéo dài hàng chục nghìn tokens ($10{,}000 - 30{,}000$ tokens) trước khi đi đến kết quả cuối cùng đúng hoặc sai ($r \in \{0, 1\}$). Mạng Critic $V_\phi$ cực kỳ khó học cách gán điểm giá trị chính xác cho từng token trung gian (bài toán Credit Assignment), dẫn đến phương sai ước lượng bùng nổ và làm sụp đổ quá trình hội tụ!
\end{itemize}

\paragraph{2. Bản Chất Toán Học Của GRPO: Triệt Tiêu Mạng Critic Bằng Chuẩn Hóa Nhóm}:
Ý tưởng cốt lõi của GRPO là: \textbf{Loại bỏ hoàn toàn mạng Critic ($V_\phi \equiv \emptyset$)}. Thay vì dùng mạng nơ-ron để đoán giá trị kỳ vọng nền của một câu hỏi, GRPO sinh ra một nhóm (Group) gồm $G$ câu trả lời ứng viên khác nhau cho cùng một câu hỏi, sau đó chuẩn hóa điểm thưởng ngay trong nhóm đó để xác định độ lợi thế tương đối:

\begin{definition}[Độ Lợi Thế Tương Đối Nhóm -- Group Relative Advantage]
Với mỗi câu hỏi (prompt) $q \sim \mathcal{D}$, mô hình chính sách cũ $\pi_{\theta_{\text{old}}}$ thực hiện lấy mẫu một nhóm gồm $G$ chuỗi câu trả lời độc lập:
\begin{equation}
\{o_1, o_2, \dots, o_G\} \sim \pi_{\theta_{\text{old}}}(O \mid q)
\end{equation}
Mỗi câu trả lời $o_i$ được đánh giá bởi hàm phần thưởng $\mathcal{R}(q, o_i)$ để cho ra một điểm thưởng số thực vô hướng $r_i$. Điểm thưởng trung bình và độ lệch chuẩn của cả nhóm được tính toán:
\begin{equation}
\bar{r} = \frac{1}{G} \sum_{i=1}^G r_i, \qquad \sigma_r = \sqrt{\frac{1}{G} \sum_{i=1}^G (r_i - \bar{r})^2}
\end{equation}
Độ lợi thế tương đối chuẩn hóa của câu trả lời thứ $i$ trong nhóm là:
\begin{equation}
\label{eq:grpo-advantage}
\widetilde{A}_i = \frac{r_i - \bar{r}}{\sigma_r + \epsilon}
\end{equation}
Trong đó $\epsilon > 0$ là hằng số chống chia cho 0. Nếu tất cả các câu trả lời trong nhóm đạt điểm thưởng giống hệt nhau ($\sigma_r = 0$), thì $\widetilde{A}_i = 0$ cho mọi $i$.
\end{definition}

\paragraph{Ý nghĩa trực giác sâu sắc của Độ lợi thế nhóm}:
\begin{itemize}
  \item \textit{Tự động khử độ khó của câu hỏi}: Đối với một câu hỏi toán học cực khó mà ngay cả các chuyên gia cũng dễ mắc lỗi, nếu 7 câu trả lời đạt điểm $0.0$ và chỉ 1 câu trả lời đạt điểm $1.0$, câu trả lời duy nhất đó sẽ nhận độ lợi thế dương cực lớn ($\widetilde{A}_i > 0$), thúc đẩy mô hình củng cố mạnh mẽ lộ trình tư duy này. Ngược lại, đối với câu hỏi quá dễ mà cả $G$ câu trả lời đều đúng ($r_i = 1.0$), độ lợi thế triệt tiêu về $0.0$, ngăn ngừa việc cập nhật gradient dư thừa lãng phí tài nguyên!
  \item \textit{Tiết kiệm $50\%$ VRAM}: Không cần bất kỳ tham số nào của mạng Critic $V_\phi$. Toàn bộ VRAM được dành trọn vẹn để nạp Actor $\pi_\theta$ và sinh chuỗi token ngữ cảnh siêu dài.
\end{itemize}

\paragraph{3. Hàm Mục Tiêu Clipped Surrogate Của GRPO}:
Hàm mục tiêu cực đại hóa của GRPO được xác định bởi:
\begin{equation}
\label{eq:grpo-loss}
\begin{split}
\mathcal{J}_{\text{GRPO}}(\theta) = \mathbb{E}_{q \sim \mathcal{D}, \{o_i\}_{i=1}^G \sim \pi_{\theta_{\text{old}}}} \Bigg[ \frac{1}{G} \sum_{i=1}^G \frac{1}{|o_i|} \sum_{t=1}^{|o_i|} & \min\Big( \rho_{i,t}(\theta) \widetilde{A}_i, \, \text{clip}\left(\rho_{i,t}(\theta), 1 - \epsilon_{\text{clip}}, 1 + \epsilon_{\text{clip}}\right) \widetilde{A}_i \Big) \\
&- \beta \mathbb{D}_{\text{KL}}(\pi_\theta \,\|\, \pi_{\text{ref}}) \Bigg]
\end{split}
\end{equation}
Trong đó:
\begin{itemize}
  \item $\rho_{i,t}(\theta) = \frac{\pi_\theta(o_{i,t} \mid q, o_{i,<t})}{\pi_{\theta_{\text{old}}}(o_{i,t} \mid q, o_{i,<t})}$: Tỷ số xác suất của token thứ $t$.
  \item $\epsilon_{\text{clip}}$: Ngưỡng cắt PPO (thường đặt $\epsilon_{\text{clip}} = 0.2$), ngăn chính sách mới đi quá xa khỏi chính sách cũ trong một bước tối ưu.
  \item $\frac{1}{|o_i|}$: Hệ số chuẩn hóa theo độ dài câu trả lời, bảo đảm các chuỗi suy luận dài không áp đảo hàm mất mát một cách nhân tạo.
  \item $\mathbb{D}_{\text{KL}}(\pi_\theta \,\|\, \pi_{\text{ref}})$: Hình phạt phân kỳ KL để giữ cho mô hình không bị sụp đổ phân phối ngôn ngữ tự nhiên. Để ước lượng không chệch và bảo đảm giá trị luôn không âm, công thức Schulman được áp dụng trên từng token:
  \begin{equation}
  \mathbb{D}_{\text{KL}}(\pi_\theta \,\|\, \pi_{\text{ref}}) \approx \frac{\pi_{\text{ref}}(o_{i,t})}{\pi_\theta(o_{i,t})} - \log \frac{\pi_{\text{ref}}(o_{i,t})}{\pi_\theta(o_{i,t})} - 1
  \end{equation}
\end{itemize}

\paragraph{4. Cơ Chế Thưởng Dựa Trên Quy Tắc (Rule-Based Verifiers) \& Miễn Nhiễm Reward Hacking}:
Trong các tác vụ đàm thoại chung, việc chấm điểm phải phụ thuộc vào Neural Reward Model (rất dễ bị khai thác lỗ hổng -- Reward Hacking theo Định luật Goodhart). Nhưng trong kỷ nguyên Reasoning, DeepSeek-R1 áp dụng nguyên tắc mang tính cách mạng: \textbf{Chỉ dùng các bộ kiểm tra tất định theo quy tắc (Rule-Based Verifiers) trên các miền chân lý khách quan (Toán học, Mã nguồn, Khoa học)}:
\begin{equation}
r_i = r_{\text{accuracy}} + \lambda_{\text{format}} r_{\text{format}}
\end{equation}
\begin{enumerate}
  \item \textbf{Phần thưởng tính chính xác ($r_{\text{accuracy}} \in \{0, 1\}$)}:
  \begin{itemize}
    \item \textit{Đối với bài toán toán học}: Mô hình được yêu cầu đưa đáp số cuối cùng vào thẻ chuẩn \verb|\boxed{...}|. Bộ kiểm tra tự động trích xuất chuỗi và dùng thư viện biểu thức đại số tượng trưng (như SymPy) để đối soát với đáp án chuẩn (Ground Truth). Nếu hoàn toàn trùng khớp $\implies r_{\text{acc}} = 1.0$, ngược lại $r_{\text{acc}} = 0.0$.
    \item \textit{Đối với bài toán lập trình}: Trích xuất đoạn mã trong khối \verb|```python| và ném vào môi trường Sandbox chạy qua bộ kiểm thử tự động (Unit Test Suite). Nếu vượt qua $100\%$ các ca kiểm thử $\implies r_{\text{acc}} = 1.0$.
  \end{itemize}
  \item \textbf{Phần thưởng định dạng ($r_{\text{format}} \in [0, 1]$)}:
  Yêu cầu mô hình bắt buộc phải đặt toàn bộ dòng suy tư nội tâm vào trong cặp thẻ đóng mở \verb|<think>| và \verb|</think>|, tiếp theo là câu trả lời kết luận. Nếu vi phạm cấu trúc thẻ $\implies r_{\text{format}} = 0.0$.
\end{enumerate}
Vì đáp án toán học và kết quả biên dịch mã nguồn là chân lý nhị phân khách quan, mô hình \textbf{hoàn toàn không có cách nào đánh lừa bộ chấm thưởng} ngoài việc tự mình suy nghĩ thật kỹ để tìm ra lời giải chính xác!

\paragraph{5. Sự Xuất Hiện Của Tư Duy System 2 \& Khoảnh Khắc Bừng Sáng (``Aha! Moment'')}:
Khi huấn luyện DeepSeek-R1-Zero thuần túy bằng thuật toán GRPO mà \textbf{hoàn toàn không có dữ liệu SFT mồi của con người}, một hiện tượng kỳ diệu tự sinh (Emergent Behavior) đã bộc lộ:
\begin{itemize}
  \item \textbf{Tự động kéo dài thời gian suy nghĩ (Test-Time Compute)}: Số lượng token suy luận trong thẻ \verb|<think>| tự động tăng vọt từ vài trăm tokens lên hàng chục nghìn tokens sau hàng nghìn bước RL.
  \item \textbf{Tự vấn và quay lui (Backtracking / Self-Correction)}: Khi đi vào ngõ cụt logic, mô hình tự phát sinh các cụm từ nội tâm: \textit{``Wait, let me rethink that assumption... Oh, that contradicts condition (2), let's try a different substitution...''}.
  \item \textbf{Khoảnh khắc bừng sáng (The Aha-Moment)}: Mô hình tự phát hiện ra cách tiếp cận bài toán từ một góc nhìn toán học thanh lịch hoàn toàn mới mà không cần bất kỳ sự can thiệp hay gợi ý nào từ kỹ sư huấn luyện!
\end{itemize}

---

\section{Ví Dụ Giải Mẫu Tính Toán Số Học Bằng Tay}
\label{sec:worked-arithmetic-examples}

\begin{example}[Tính Toán Lượng Tử Hóa Đối Xứng INT8 Bằng Tay]
Xét véc-tơ kích hoạt gồm 6 giá trị số thực dạng FP32:
\begin{equation*}
\mathbf{x} = [-10.0, -5.0, 0.0, 2.5, 7.5, 10.0]
\end{equation*}
Hãy thực hiện lượng tử hóa đối xứng sang định dạng số nguyên 8-bit có dấu ($q \in [-127, 127]$).
\end{example}

\paragraph{Lời giải chi tiết từng bước}:
\begin{itemize}
  \item[1)] Tính giá trị cực đại tuyệt đối và hệ số co giãn $S$:
  \begin{equation*}
  \max(|\mathbf{x}|) = 10.0 \implies S = \frac{10.0}{127} \approx 0.07874016
  \end{equation*}
  \item[2)] Lượng tử hóa từng phần tử: $q_i = \text{clip}\left(\text{round}(x_i / S), -127, 127\right)$:
  \begin{itemize}
    \item $q_1 = \text{round}(-10.0 / 0.07874016) = -127$
    \item $q_2 = \text{round}(-5.0 / 0.07874016) = -64$
    \item $q_3 = \text{round}(0.0 / 0.07874016) = 0$
    \item $q_4 = \text{round}(2.5 / 0.07874016) = 32$
    \item $q_5 = \text{round}(7.5 / 0.07874016) = 95$
    \item $q_6 = \text{round}(10.0 / 0.07874016) = 127$
  \end{itemize}
  Véc-tơ số nguyên thu được: $\mathbf{q} = [-127, -64, 0, 32, 95, 127]$.
  \item[3)] Giải lượng tử hóa $\hat{x}_i = q_i \times S$:
  \begin{equation*}
  \hat{\mathbf{x}} \approx [-10.0000, -5.0394, 0.0000, 2.5197, 7.4803, 10.0000]
  \end{equation*}
  \item[4)] Sai số bình phương trung bình:
  \begin{equation*}
  \text{MSE} = \frac{0.0 + (0.0394)^2 + 0.0 + (-0.0197)^2 + (0.0197)^2 + 0.0}{6} \approx \frac{0.00233}{6} \approx 0.000388
  \end{equation*}
\end{itemize}

---

\begin{example}[Tính Toán Vector Chuyển Đổi SmoothQuant Bằng Tay]
Xét ma trận kích hoạt $\mathbf{X} \in \mathbb{R}^{2 \times 2}$ và ma trận trọng số $\mathbf{W} \in \mathbb{R}^{2 \times 2}$:
\begin{equation*}
\mathbf{X} = \begin{pmatrix} 2.0 & 80.0 \\ 1.0 & 40.0 \end{pmatrix}, \qquad \mathbf{W} = \begin{pmatrix} 4.0 & 2.0 \\ 0.5 & 0.2 \end{pmatrix}
\end{equation*}
Nhận xét: Kênh thứ hai của $\mathbf{X}$ có giá trị ngoại lai rất lớn ($80.0$), trong khi trọng số tương ứng ở kênh đó lại rất nhỏ ($0.5$).
Hãy tính:
\begin{enumerate}
  \item Tích ma trận gốc $\mathbf{Y} = \mathbf{X} \mathbf{W}$.
  \item Véc-tơ hệ số di dời $\mathbf{s} = (s_1, s_2)$ với $\alpha = 0.5$.
  \item Ma trận kích hoạt mới $\hat{\mathbf{X}} = \mathbf{X} \cdot \text{diag}(\mathbf{s})^{-1}$ và ma trận trọng số mới $\hat{\mathbf{W}} = \text{diag}(\mathbf{s}) \cdot \mathbf{W}$.
  \item Kiểm chứng tính tương đương $\hat{\mathbf{X}} \hat{\mathbf{W}} = \mathbf{Y}$.
\end{enumerate}
\end{example}

\paragraph{Lời giải chi tiết}:
\begin{itemize}
  \item[1)] Tính tích ma trận gốc:
  \begin{equation*}
  \mathbf{Y} = \begin{pmatrix} 2.0 \times 4.0 + 80.0 \times 0.5 & 2.0 \times 2.0 + 80.0 \times 0.2 \\ 1.0 \times 4.0 + 40.0 \times 0.5 & 1.0 \times 2.0 + 40.0 \times 0.2 \end{pmatrix} = \begin{pmatrix} 48.0 & 20.0 \\ 24.0 & 10.0 \end{pmatrix}
  \end{equation*}

  \item[2)] Xác định biên độ cực đại từng kênh:
  \begin{itemize}
    \item Cột 1: $\max(|\mathbf{X}_1|) = 2.0$, hàng 1 của $\mathbf{W}$: $\max(|\mathbf{W}_1|) = 4.0$.
    \item Cột 2: $\max(|\mathbf{X}_2|) = 80.0$, hàng 2 của $\mathbf{W}$: $\max(|\mathbf{W}_2|) = 0.5$.
  \end{itemize}
  Với $\alpha = 0.5$, công thức $s_j = \sqrt{\frac{\max(|\mathbf{X}_j|)}{\max(|\mathbf{W}_j|)}}$:
  \begin{equation*}
  s_1 = \sqrt{\frac{2.0}{4.0}} = \sqrt{0.5} \approx 0.707107
  \end{equation*}
  \begin{equation*}
  s_2 = \sqrt{\frac{80.0}{0.5}} = \sqrt{160.0} \approx 12.64911
  \end{equation*}

  \item[3)] Biến đổi $\hat{\mathbf{X}}$ và $\hat{\mathbf{W}}$:
  \begin{equation*}
  \text{diag}(\mathbf{s})^{-1} = \begin{pmatrix} \frac{1}{0.707107} & 0 \\ 0 & \frac{1}{12.64911} \end{pmatrix} \approx \begin{pmatrix} 1.414214 & 0 \\ 0 & 0.079057 \end{pmatrix}
  \end{equation*}
  \begin{equation*}
  \hat{\mathbf{X}} = \mathbf{X} \cdot \text{diag}(\mathbf{s})^{-1} = \begin{pmatrix} 2.0 \times 1.414214 & 80.0 \times 0.079057 \\ 1.0 \times 1.414214 & 40.0 \times 0.079057 \end{pmatrix} \approx \begin{pmatrix} 2.8284 & 6.3246 \\ 1.4142 & 3.1623 \end{pmatrix}
  \end{equation*}
  Biên độ kênh 2 của kích hoạt đã giảm từ $80.0$ xuống còn $6.32$ (giảm hơn $12$ lần)!
  
  Ma trận trọng số mới:
  \begin{equation*}
  \hat{\mathbf{W}} = \text{diag}(\mathbf{s}) \cdot \mathbf{W} = \begin{pmatrix} 0.707107 \times 4.0 & 0.707107 \times 2.0 \\ 12.64911 \times 0.5 & 12.64911 \times 0.2 \end{pmatrix} \approx \begin{pmatrix} 2.8284 & 1.4142 \\ 6.3246 & 2.5298 \end{pmatrix}
  \end{equation*}

  \item[4)] Kiểm tra tích mới:
  \begin{align*}
  \hat{y}_{11} &= 2.8284 \times 2.8284 + 6.3246 \times 6.3246 = 8.00 + 40.00 = 48.0 \\
  \hat{y}_{12} &= 2.8284 \times 1.4142 + 6.3246 \times 2.5298 = 4.00 + 16.00 = 20.0 \\
  \hat{y}_{21} &= 1.4142 \times 2.8284 + 3.1623 \times 6.3246 = 4.00 + 20.00 = 24.0 \\
  \hat{y}_{22} &= 1.4142 \times 1.4142 + 3.1623 \times 2.5298 = 2.00 + 8.00 = 10.0
  \end{align*}
  Kết quả thu được $\hat{\mathbf{Y}} = \begin{pmatrix} 48.0 & 20.0 \\ 24.0 & 10.0 \end{pmatrix}$, khớp chính xác $100\%$ với ma trận gốc!
\end{itemize}

---

\begin{example}[Tính Toán Hàm Mất Mát DPO \& Implicit Reward Bằng Tay]
Xét một mẫu huấn luyện sở thích $(x, y_w, y_l)$ với hệ số phạt KL $\beta = 0.1$.
Xác suất mô hình hiện tại $\pi_\theta$ và mô hình tham chiếu $\pi_{\text{ref}}$ sinh ra hai câu trả lời là:
\begin{itemize}
  \item Câu trả lời thắng: $\pi_\theta(y_w \mid x) = 0.40$, $\pi_{\text{ref}}(y_w \mid x) = 0.20$.
  \item Câu trả lời thua: $\pi_\theta(y_l \mid x) = 0.10$, $\pi_{\text{ref}}(y_l \mid x) = 0.20$.
\end{itemize}
Hãy tính:
\begin{enumerate}
  \item Điểm thưởng tiềm ẩn $\hat{r}_\theta(y_w)$ và $\hat{r}_\theta(y_l)$.
  \item Lề chênh lệch điểm thưởng $\Delta r = \hat{r}_\theta(y_w) - \hat{r}_\theta(y_l)$.
  \item Giá trị hàm mất mát $\mathcal{L}_{\text{DPO}}$.
\end{enumerate}
\end{example}

\paragraph{Lời giải chi tiết}:
\begin{itemize}
  \item[1)] Tính log-ratios và điểm thưởng tiềm ẩn:
  \begin{equation*}
  \hat{r}_\theta(y_w) = \beta \log \frac{\pi_\theta(y_w \mid x)}{\pi_{\text{ref}}(y_w \mid x)} = 0.1 \times \log \left(\frac{0.40}{0.20}\right) = 0.1 \times \log(2.0) \approx 0.1 \times 0.69315 = 0.069315
  \end{equation*}
  \begin{equation*}
  \hat{r}_\theta(y_l) = \beta \log \frac{\pi_\theta(y_l \mid x)}{\pi_{\text{ref}}(y_l \mid x)} = 0.1 \times \log \left(\frac{0.10}{0.20}\right) = 0.1 \times \log(0.5) \approx 0.1 \times (-0.69315) = -0.069315
  \end{equation*}

  \item[2)] Lề chênh lệch điểm thưởng:
  \begin{equation*}
  \Delta r = \hat{r}_\theta(y_w) - \hat{r}_\theta(y_l) = 0.069315 - (-0.069315) = 0.13863
  \end{equation*}

  \item[3)] Giá trị hàm mất mát DPO:
  \begin{equation*}
  \mathcal{L}_{\text{DPO}} = -\log \sigma\left( \Delta r \right) = -\log \sigma\left( 0.13863 \right)
  \end{equation*}
  Ta có:
  \begin{equation*}
  \sigma(0.13863) = \frac{1}{1 + e^{-0.13863}} \approx \frac{1}{1 + 0.87055} \approx \frac{1}{1.87055} \approx 0.53460
  \end{equation*}
  Do đó:
  \begin{equation*}
  \mathcal{L}_{\text{DPO}} = -\log(0.53460) \approx 0.62623
  \end{equation*}
  \textit{Bình luận}: Khi $\beta = 0.1$, giá trị hàm mất mát là $0.62623$. Nếu không có hệ số co $\beta$ (tương đương $\beta = 1.0$), $\Delta r = \log(4) \approx 1.38629$, khi đó $\sigma(1.38629) = 0.80$ và mất mát giảm xuống $-\log(0.80) \approx 0.22314$. Hệ số $\beta$ đóng vai trò điều tiết mức độ tự tin và cường độ phạt độ lệch khỏi mô hình tham chiếu.
\end{itemize}

---

\begin{example}[Tính Toán Độ Lợi Thế Nhóm GRPO (DeepSeek-R1) Bằng Tay]
Xét quá trình huấn luyện mô hình suy luận toán học với thuật toán GRPO. Cho một câu hỏi toán học $q$. Mô hình chính sách cũ sinh ra một nhóm gồm $G = 4$ câu trả lời độc lập $\{o_1, o_2, o_3, o_4\}$.
Bộ chấm điểm quy tắc đánh giá:
\begin{itemize}
  \item $o_1$: Có khối thẻ \verb|<think>| lập luận chặt chẽ và kết luận đáp số đúng \verb|\boxed{42}| ($r_{\text{acc}} = 1.0$, $r_{\text{format}} = 1.0$).
  \item $o_2$: Đưa ra đáp số đúng \verb|\boxed{42}| nhưng không có thẻ suy nghĩ \verb|<think>| ($r_{\text{acc}} = 1.0$, $r_{\text{format}} = 0.0$).
  \item $o_3$: Có thẻ suy nghĩ \verb|<think>| nhưng tính toán sai đáp số \verb|\boxed{15}| ($r_{\text{acc}} = 0.0$, $r_{\text{format}} = 1.0$).
  \item $o_4$: Không có thẻ suy nghĩ và đáp số sai \verb|\boxed{99}| ($r_{\text{acc}} = 0.0$, $r_{\text{format}} = 0.0$).
\end{itemize}
Giả sử trọng số định dạng là $\lambda_{\text{format}} = 0.5$, tức hàm thưởng tổng là $r_i = r_{\text{acc}} + 0.5 \times r_{\text{format}}$.
Hãy tính:
\begin{enumerate}
  \item Điểm thưởng tuyệt đối $r_i$ cho từng câu trả lời trong nhóm.
  \item Điểm thưởng trung bình $\bar{r}$ và độ lệch chuẩn nhóm $\sigma_r$.
  \item Độ lợi thế tương đối chuẩn hóa $\widetilde{A}_i$ cho từng mẫu và nhận xét về hướng cập nhật gradient của mô hình.
\end{enumerate}
\end{example}

\paragraph{Lời giải chi tiết}:
\begin{itemize}
  \item[1)] Tính điểm thưởng từng câu trả lời:
  \begin{align*}
  r_1 &= 1.0 + 0.5 \times 1.0 = 1.5 \\
  r_2 &= 1.0 + 0.5 \times 0.0 = 1.0 \\
  r_3 &= 0.0 + 0.5 \times 1.0 = 0.5 \\
  r_4 &= 0.0 + 0.5 \times 0.0 = 0.0
  \end{align*}

  \item[2)] Tính điểm thưởng trung bình và độ lệch chuẩn của nhóm $G=4$:
  \begin{equation*}
  \bar{r} = \frac{1.5 + 1.0 + 0.5 + 0.0}{4} = \frac{3.0}{4} = 0.75
  \end{equation*}
  Độ lệch so với trung bình:
  \begin{align*}
  (r_1 - \bar{r})^2 &= (1.5 - 0.75)^2 = (0.75)^2 = 0.5625 \\
  (r_2 - \bar{r})^2 &= (1.0 - 0.75)^2 = (0.25)^2 = 0.0625 \\
  (r_3 - \bar{r})^2 &= (0.5 - 0.75)^2 = (-0.25)^2 = 0.0625 \\
  (r_4 - \bar{r})^2 &= (0.0 - 0.75)^2 = (-0.75)^2 = 0.5625
  \end{align*}
  Phương sai nhóm:
  \begin{equation*}
  \sigma_r^2 = \frac{0.5625 + 0.0625 + 0.0625 + 0.5625}{4} = \frac{1.25}{4} = 0.3125
  \end{equation*}
  Độ lệch chuẩn nhóm:
  \begin{equation*}
  \sigma_r = \sqrt{0.3125} \approx 0.559017
  \end{equation*}

  \item[3)] Tính độ lợi thế tương đối chuẩn hóa $\widetilde{A}_i = \frac{r_i - \bar{r}}{\sigma_r}$:
  \begin{align*}
  \widetilde{A}_1 &= \frac{1.5 - 0.75}{0.559017} = \frac{0.75}{0.559017} \approx \mathbf{+1.3416} \\
  \widetilde{A}_2 &= \frac{1.0 - 0.75}{0.559017} = \frac{0.25}{0.559017} \approx \mathbf{+0.4472} \\
  \widetilde{A}_3 &= \frac{0.5 - 0.75}{0.559017} = \frac{-0.25}{0.559017} \approx \mathbf{-0.4472} \\
  \widetilde{A}_4 &= \frac{0.0 - 0.75}{0.559017} = \frac{-0.75}{0.559017} \approx \mathbf{-1.3416}
  \end{align*}
  \textit{Bình luận ý nghĩa kỹ nghệ}:
  \begin{itemize}
    \item Tổng độ lợi thế nhóm triệt tiêu chính xác về 0: $\sum_{i=1}^4 \widetilde{A}_i = 1.3416 + 0.4472 - 0.4472 - 1.3416 = 0.0000$.
    \item Câu trả lời $o_1$ (vừa suy luận vừa đáp số chuẩn) nhận phần thưởng tối đa $\widetilde{A}_1 = +1.34$, kích thích chính sách tăng xác suất sinh ra cả thẻ suy nghĩ lẫn chuỗi logic.
    \item Câu trả lời $o_2$ (đoán mò đúng đáp số nhưng không suy nghĩ) chỉ nhận $\widetilde{A}_2 = +0.45$, thấp hơn đáng kể so với $o_1$.
    \item Câu trả lời $o_3$ (có cố gắng suy nghĩ nhưng sai) bị phạt nhẹ ($\widetilde{A}_3 = -0.45$), trong khi câu trả lời $o_4$ (vừa lười suy nghĩ vừa sai) bị phạt cực nặng ($\widetilde{A}_4 = -1.34$). Cơ chế phân bậc này ép mô hình bắt buộc phải kích hoạt chuỗi suy luận dài (System 2 thinking)!
  \end{itemize}
\end{itemize}

---

\begin{example}[Tính Toán Lượng Tử Hóa Bậc Ba BitNet b1.58 \& Phép Cộng Số Nguyên Bằng Tay]
Xét một ma trận trọng số thu nhỏ $\mathbf{W} \in \mathbb{R}^{2 \times 3}$ và véc-tơ kích hoạt đầu vào $\mathbf{x} \in \mathbb{R}^{1 \times 3}$:
\begin{equation*}
\mathbf{W} = \begin{pmatrix} 0.8 & -0.1 & 0.6 \\ -0.7 & 0.9 & 0.05 \end{pmatrix}, \qquad \mathbf{x} = [2.0, -4.0, 1.0]
\end{equation*}
Hãy thực hiện quy trình nén và tính toán theo chuẩn BitNet b1.58:
\begin{enumerate}
  \item Tính hệ số co giãn Absmean $\gamma$ và ma trận trọng số bậc ba $\widetilde{\mathbf{W}} \in \{-1, 0, +1\}^{2 \times 3}$.
  \item Lượng tử hóa véc-tơ kích hoạt $\mathbf{x}$ sang số nguyên 8-bit đối xứng $\mathbf{x}_q \in [-127, 127]^{1 \times 3}$.
  \item Thực hiện phép nhân ma trận $\mathbf{y}_q = \mathbf{x}_q \widetilde{\mathbf{W}}^T$ thuần túy bằng các phép cộng và trừ số nguyên (triệt tiêu hoàn toàn phép nhân phần cứng).
  \item Phục hồi giá trị đầu ra số thực $\hat{\mathbf{y}}$ và so sánh với giá trị tính toán FP32 chuẩn.
\end{enumerate}
\end{example}

\paragraph{Lời giải chi tiết}:
\begin{itemize}
  \item[1)] Tính hệ số Absmean $\gamma$:
  \begin{equation*}
  \gamma = \frac{\sum_{i,j} |W_{i,j}|}{2 \times 3} = \frac{0.8 + 0.1 + 0.6 + 0.7 + 0.9 + 0.05}{6} = \frac{3.15}{6} = 0.525
  \end{equation*}
  Lượng tử hóa bậc ba $\widetilde{W}_{i,j} = \text{RoundClip}\left(\frac{W_{i,j}}{\gamma}, -1, 1\right)$:
  \begin{itemize}
    \item Hàng 1:
      \begin{itemize}
        \item $\frac{0.8}{0.525} \approx 1.5238 \implies \text{round}(1.5238) = 2 \implies \text{clip}(2, -1, 1) = \mathbf{+1}$
        \item $\frac{-0.1}{0.525} \approx -0.1905 \implies \text{round}(-0.1905) = 0 \implies \text{clip}(0, -1, 1) = \mathbf{0}$
        \item $\frac{0.6}{0.525} \approx 1.1429 \implies \text{round}(1.1429) = 1 \implies \text{clip}(1, -1, 1) = \mathbf{+1}$
      \end{itemize}
    \item Hàng 2:
      \begin{itemize}
        \item $\frac{-0.7}{0.525} \approx -1.3333 \implies \text{round}(-1.3333) = -1 \implies \text{clip}(-1, -1, 1) = \mathbf{-1}$
        \item $\frac{0.9}{0.525} \approx 1.7143 \implies \text{round}(1.7143) = 2 \implies \text{clip}(2, -1, 1) = \mathbf{+1}$
        \item $\frac{0.05}{0.525} \approx 0.0952 \implies \text{round}(0.0952) = 0 \implies \text{clip}(0, -1, 1) = \mathbf{0}$
      \end{itemize}
  \end{itemize}
  Ma trận trọng số bậc ba thu được:
  \begin{equation*}
  \widetilde{\mathbf{W}} = \begin{pmatrix} 1 & 0 & 1 \\ -1 & 1 & 0 \end{pmatrix}
  \end{equation*}

  \item[2)] Lượng tử hóa véc-tơ kích hoạt $\mathbf{x} = [2.0, -4.0, 1.0]$ sang INT8:
  Biên độ cực đại: $\max(|\mathbf{x}|) = 4.0$. Hệ số tỷ lệ: $s_x = \frac{4.0}{127} \approx 0.031496$.
  \begin{align*}
  x_{q, 1} &= \text{round}\left(\frac{2.0}{0.031496}\right) = \text{round}(63.5) = 64 \\
  x_{q, 2} &= \text{round}\left(\frac{-4.0}{0.031496}\right) = \text{round}(-127.0) = -127 \\
  x_{q, 3} &= \text{round}\left(\frac{1.0}{0.031496}\right) = \text{round}(31.75) = 32
  \end{align*}
  Véc-tơ số nguyên INT8: $\mathbf{x}_q = [64, -127, 32]$.

  \item[3)] Phép nhân ma trận thuần cộng/trừ số nguyên:
  \begin{itemize}
    \item Phần tử đầu ra 1: Hàng 1 của $\widetilde{\mathbf{W}}$ là $[1, 0, 1]$. Vị trí 1 là $+1$, vị trí 2 là $0$ (bỏ qua), vị trí 3 là $+1$:
    \begin{equation*}
    y_{q, 1} = x_{q, 1} + x_{q, 3} = 64 + 32 = \mathbf{96}
    \end{equation*}
    \item Phần tử đầu ra 2: Hàng 2 của $\widetilde{\mathbf{W}}$ là $[-1, 1, 0]$. Vị trí 1 là $-1$ (trừ), vị trí 2 là $+1$ (cộng), vị trí 3 là $0$ (bỏ qua):
    \begin{equation*}
    y_{q, 2} = -x_{q, 1} + x_{q, 2} = -64 + (-127) = \mathbf{-191}
    \end{equation*}
  \end{itemize}
  Kết quả tích lũy số nguyên: $\mathbf{y}_q = [96, -191]$. Toàn bộ quá trình \textbf{hoàn toàn không sử dụng bất kỳ mạch nhân phần cứng nào}!

  \item[4)] Khôi phục giá trị số thực $\hat{\mathbf{y}} = \mathbf{y}_q \times (s_x \cdot \gamma)$:
  Tích hệ số co giãn: $s_x \cdot \gamma = 0.031496 \times 0.525 \approx 0.0165354$.
  \begin{align*}
  \hat{y}_1 &= 96 \times 0.0165354 \approx \mathbf{1.5874} \\
  \hat{y}_2 &= -191 \times 0.0165354 \approx \mathbf{-3.1583}
  \end{align*}
  So sánh với kết quả phép nhân FP32 chuẩn:
  \begin{align*}
  y_{\text{FP32}, 1} &= 2.0 \times 0.8 + (-4.0) \times (-0.1) + 1.0 \times 0.6 = 1.6 + 0.4 + 0.6 = \mathbf{2.6000} \\
  y_{\text{FP32}, 2} &= 2.0 \times (-0.7) + (-4.0) \times 0.9 + 1.0 \times 0.05 = -1.4 - 3.6 + 0.05 = \mathbf{-4.9500}
  \end{align*}
  Tỷ số tương quan hướng giữa hai vector đạt $0.999$, bảo toàn chuẩn xác cực tính và cấu trúc không gian đặc trưng của tín hiệu!
\end{itemize}

---

\section{Hiện Thực Hóa: Mã Nguồn Python Chuẩn Kỹ Nghệ}
\label{sec:python-production-modules}

Dưới đây là các module Python thuần chuẩn mực, tự chứa và không phụ thuộc vào thư viện ngoài, hiện thực hóa bộ lượng tử hóa đối xứng, bộ chuyển đổi SmoothQuant, lớp LoRA Linear đa khách hàng, và động cơ tính mất mát DPO:

\begin{lstlisting}[language=Python, caption={Bộ lượng tử hóa đối xứng INT8 và giải lượng tử hóa thuần NumPy.}]
from typing import Tuple
import numpy as np

def symmetric_quantize_int8(x: np.ndarray) -> Tuple[np.ndarray, float]:
    """
    Luong tu hoa doi xung tensor so thuc sang INT8 (-127 den 127).
    Cong thuc: Scale = max(|x|) / 127, q = clip(round(x / Scale), -127, 127).
    """
    max_val = float(np.max(np.abs(x)))
    if max_val == 0.0:
        return np.zeros_like(x, dtype=np.int8), 1.0

    scale = max_val / 127.0
    q = np.clip(np.round(x / scale), -127, 127).astype(np.int8)
    return q, scale

def dequantize_int8(q: np.ndarray, scale: float) -> np.ndarray:
    """Giai luong tu hoa INT8 ve float32: x_hat = q * scale."""
    return q.astype(np.float32) * scale
\end{lstlisting}

\begin{lstlisting}[language=Python, caption={Bộ mô phỏng chuyển đổi SmoothQuant cân bằng activations và weights.}]
class SmoothQuantSimulator:
    """
    Mo phong bien doi SmoothQuant di doi do dot bien giua X va W:
    Y = X @ W = (X * diag(s)^-1) @ (diag(s) * W) = X_hat @ W_hat
    """
    @staticmethod
    def calculate_scale_vector(x: np.ndarray, w: np.ndarray, alpha: float = 0.5) -> np.ndarray:
        """
        Tinh vector ti le s_j = (max|X_j|^alpha) / (max|W_j|^(1-alpha)).
        X co dang (batch_size, hidden_dim), W co dang (hidden_dim, out_dim).
        """
        act_max = np.max(np.abs(x), axis=0)  # (hidden_dim,)
        weight_max = np.max(np.abs(w), axis=1)  # (hidden_dim,)

        act_max = np.maximum(act_max, 1e-5)
        weight_max = np.maximum(weight_max, 1e-5)

        scale_vector = (act_max ** alpha) / (weight_max ** (1.0 - alpha))
        return scale_vector

    @staticmethod
    def transform(x: np.ndarray, w: np.ndarray, s: np.ndarray) -> Tuple[np.ndarray, np.ndarray]:
        """Thuc hien phep bien doi: X_hat = X / s va W_hat = W * s."""
        # s: (hidden_dim,) -> chia theo tung cot cua X va nhan vao tung hang cua W
        x_hat = x / s
        w_hat = w * s[:, np.newaxis]
        return x_hat, w_hat
\end{lstlisting}

\begin{lstlisting}[language=Python, caption={Module LoRA Linear với cơ chế Zero-Latency Weight Merging.}]
class SimpleLoRALayer:
    """
    Tang tuyen tinh gan adapter LoRA: y = x @ W0 + (alpha / r) * (x @ A @ B)
    """
    def __init__(self, d_in: int, d_out: int, r: int = 8, alpha: float = 16.0):
        self.d_in = d_in
        self.d_out = d_out
        self.r = r
        self.alpha = alpha
        self.scaling = alpha / r

        # Trong so goc dong bang (Frozen)
        self.w0 = np.random.standard_normal((d_in, d_out)).astype(np.float32)

        # Ma tran adapter: A khoi tao Gauss, B khoi tao Zero
        self.a = np.random.standard_normal((d_in, r)).astype(np.float32) * 0.01
        self.b = np.zeros((r, d_out), dtype=np.float32)

    def forward(self, x: np.ndarray) -> np.ndarray:
        """Lan truyen xuoi: nhanh goc cong voi nhanh LoRA."""
        base_out = np.matmul(x, self.w0)
        lora_out = self.scaling * np.matmul(np.matmul(x, self.a), self.b)
        return base_out + lora_out

    def merge_weights(self) -> np.ndarray:
        """Hop nhat vinh vien trong so LoRA vao trong so goc cho suy luan zero-overhead."""
        delta_w = self.scaling * np.matmul(self.a, self.b)
        return self.w0 + delta_w
\end{lstlisting}

\begin{lstlisting}[language=Python, caption={Động cơ tính toán hàm mất mát và phần thưởng tiềm ẩn DPO (Direct Preference Optimization).}]
class DPOLossEngine:
    """
    Dong co tinh toan ham mat mat DPO va implicit rewards:
    L_DPO = -E [ log(sigmoid(beta * (log_pi(yw) - log_ref(yw)) - beta * (log_pi(yl) - log_ref(yl)))) ]
    """
    def __init__(self, beta: float = 0.1):
        self.beta = beta

    def compute_loss(
        self,
        policy_chosen_logps: np.ndarray,
        policy_rejected_logps: np.ndarray,
        reference_chosen_logps: np.ndarray,
        reference_rejected_logps: np.ndarray
    ) -> Tuple[float, float, float, float]:
        """
        Dau vao la log-probabilities cua cau tra loi thang (chosen) va thua (rejected).
        Tra ve: (dpo_loss, chosen_reward, rejected_reward, accuracy).
        """
        # Tinh log-ratio cho chosen va rejected
        pi_logratios = policy_chosen_logps - policy_rejected_logps
        ref_logratios = reference_chosen_logps - reference_rejected_logps

        logits = pi_logratios - ref_logratios
        scaled_logits = self.beta * logits

        # Ham sigmoid on dinh so hoc
        losses = -self._log_sigmoid(scaled_logits)
        mean_loss = float(np.mean(losses))

        # Tinh implicit rewards
        chosen_rewards = self.beta * (policy_chosen_logps - reference_chosen_logps)
        rejected_rewards = self.beta * (policy_rejected_logps - reference_rejected_logps)

        reward_accuracy = float(np.mean(chosen_rewards > rejected_rewards))

        return (
            mean_loss,
            float(np.mean(chosen_rewards)),
            float(np.mean(rejected_rewards)),
            reward_accuracy
        )

    @staticmethod
    def _log_sigmoid(x: np.ndarray) -> np.ndarray:
        """Tinh log(sigmoid(x)) on dinh so hoc tranh tran so exp."""
        return -np.logaddexp(0.0, -x)
\end{lstlisting}

\begin{lstlisting}[language=Python, caption={Bộ ước lượng lợi thế nhóm GRPO và bộ kiểm tra quy tắc suy luận (DeepSeek-R1).}]
import re
from typing import Dict, List, Optional, Tuple
import numpy as np

class GRPOAdvantageEstimator:
    """
    Uoc luong do loi the nhom (Group Relative Advantage) trong GRPO (DeepSeek-R1).
    A_i = (r_i - mean(R)) / (std(R) + eps)
    Triet tieu hoan toan mang Critic V_phi, tiet kiem 50% VRAM.
    """
    @staticmethod
    def extract_boxed_answer(text: str) -> Optional[str]:
        """Trich xuat dap an trong the \\boxed{...}."""
        pattern = r"\\boxed\{([^}]+)\}"
        match = re.search(pattern, text)
        return match.group(1).strip() if match else None

    @classmethod
    def rule_based_verifier(
        cls, completion: str, ground_truth: str, format_weight: float = 0.5
    ) -> Dict[str, float]:
        """
        Bo kiem tra quy tac tat dinh:
        - format_reward: bat buoc co the <think>...</think> va the \\boxed{...}
        - accuracy_reward: kiem tra dap an trung khop voi ground_truth
        """
        has_think = bool(re.search(r"<think>.*?</think>", completion, re.DOTALL))
        extracted = cls.extract_boxed_answer(completion)
        has_boxed = extracted is not None

        format_reward = 1.0 if (has_think and has_boxed) else (0.5 if (has_think or has_boxed) else 0.0)
        acc_reward = 1.0 if (extracted is not None and extracted.lower() == ground_truth.strip().lower()) else 0.0
        total_reward = acc_reward + format_weight * format_reward
        return {
            "format_reward": format_reward,
            "accuracy_reward": acc_reward,
            "total_reward": total_reward,
        }

    @staticmethod
    def compute_group_advantages(rewards: List[float], eps: float = 1e-6) -> np.ndarray:
        """
        Chuan hoa diem thuong trong nhom G cau tra loi:
        A_i = (r_i - mean(R)) / (std(R) + eps)
        """
        r = np.array(rewards, dtype=np.float64)
        mean_r = np.mean(r)
        std_r = np.std(r)
        if std_r < 1e-8:
            return np.zeros_like(r)
        return (r - mean_r) / (std_r + eps)

    @staticmethod
    def compute_clipped_surrogate_loss(
        pi_logps: np.ndarray,
        old_logps: np.ndarray,
        advantages: np.ndarray,
        clip_eps: float = 0.2
    ) -> Tuple[float, float]:
        """
        Tinh ton that Clipped Surrogate PPO:
        loss = -min(ratio * A, clip(ratio, 1-eps, 1+eps) * A)
        """
        ratios = np.exp(pi_logps - old_logps)
        surr1 = ratios * advantages
        clipped = np.clip(ratios, 1.0 - clip_eps, 1.0 + clip_eps) * advantages
        loss = -float(np.mean(np.minimum(surr1, clipped)))
        clip_frac = float(np.mean(ratios != np.clip(ratios, 1.0 - clip_eps, 1.0 + clip_eps)))
        return loss, clip_frac
\end{lstlisting}

\begin{lstlisting}[language=Python, caption={Lớp tuyến tính BitNet b1.58 với phép nhân ma trận thuần cộng/trừ số nguyên.}]
class BitNetTernaryLinear:
    """
    Lop tuyen tinh BitNet b1.58 voi trong so bac ba {-1, 0, +1} va kich hoat INT8.
    Phep tinh ma tran Y = X @ W.T duoc thuc hien thuan tuy bang phep cong va tru so nguyen:
    y_q = sum_{W=+1} x_q - sum_{W=-1} x_q (triet tieu hoan toan phep nhan phan cung).
    """
    def __init__(self, d_in: int, d_out: int):
        self.d_in = d_in
        self.d_out = d_out
        self.w = np.random.standard_normal((d_out, d_in)).astype(np.float32)

    @staticmethod
    def quantize_weights(w: np.ndarray, eps: float = 1e-6) -> Tuple[np.ndarray, float]:
        """Luong tu hoa trong so sang {-1, 0, +1} bang Absmean scaling."""
        gamma = float(np.mean(np.abs(w)))
        scale = max(gamma, eps)
        w_ternary = np.clip(np.round(w / scale), -1, 1).astype(np.int8)
        return w_ternary, scale

    @staticmethod
    def quantize_activations(x: np.ndarray, bits: int = 8) -> Tuple[np.ndarray, float]:
        """Luong tu hoa kich hoat sang INT8 doi xung [-127, 127]."""
        q_max = (1 << (bits - 1)) - 1
        max_abs = float(np.max(np.abs(x)))
        scale = 1.0 if max_abs < 1e-8 else max_abs / q_max
        x_q = np.clip(np.round(x / scale), -q_max, q_max).astype(np.int8)
        return x_q, scale

    @staticmethod
    def forward_pure_addition(x_q: np.ndarray, w_ternary: np.ndarray) -> np.ndarray:
        """
        Phep tinh lan truyen xuoi thuan cong/tru so nguyen (Zero Multiplications):
        - Cac phan tu co W = +1: duoc cong vao bo tich luy
        - Cac phan tu co W = -1: duoc tru khoi bo tich luy
        - Cac phan tu co W = 0: bo qua
        """
        batch_size, d_in = x_q.shape
        d_out, _ = w_ternary.shape
        y_acc = np.zeros((batch_size, d_out), dtype=np.int32)

        for b in range(batch_size):
            row_x = x_q[b].astype(np.int32)
            for j in range(d_out):
                row_w = w_ternary[j]
                plus_sum = np.sum(row_x[row_w == 1])
                minus_sum = np.sum(row_x[row_w == -1])
                y_acc[b, j] = plus_sum - minus_sum

        return y_acc
\end{lstlisting}

\begin{lstlisting}[language=Python, caption={Bộ tối ưu hóa hệ số tỷ lệ AWQ và kỹ thuật gập vào LayerNorm (Lin et al., 2024).}]
class AWQScaleOptimizer:
    """
    Activation-aware Weight Quantization (AWQ) Scale Optimizer (Lin et al., 2023).
    Finds per-channel protection scales s = s_X^alpha that minimize:
    min_{alpha} || W @ X - Q(W * diag(s)) @ diag(s)^-1 @ X ||_F^2
    """
    def __init__(self, bits: int = 4, n_grid: int = 20):
        self.bits = bits
        self.n_grid = n_grid
        self.q_max = (1 << (bits - 1)) - 1  # 7 for 4-bit signed symmetric

    def quantize_symmetric(self, w: np.ndarray) -> np.ndarray:
        """Per-channel symmetric quantization along axis 1 (per row)."""
        scales = np.max(np.abs(w), axis=1, keepdims=True) / self.q_max
        scales = np.maximum(scales, 1e-8)
        q = np.clip(np.round(w / scales), -self.q_max, self.q_max)
        return q * scales

    def optimize_scale(
        self,
        w: np.ndarray,
        x: np.ndarray
    ) -> Tuple[np.ndarray, float, float]:
        """
        w: (d_out, d_in)
        x: (d_in, N) calibration activations
        Returns: (best_scale_s, best_alpha, best_loss)
        """
        act_mean = np.mean(np.abs(x), axis=1)
        act_mean = np.maximum(act_mean, 1e-6)

        # Baseline loss without AWQ (s = 1.0)
        y_true = np.matmul(w, x)
        w_q_base = self.quantize_symmetric(w)
        y_base = np.matmul(w_q_base, x)
        baseline_loss = float(np.mean((y_true - y_base) ** 2))

        best_loss = baseline_loss
        best_alpha = 0.0
        best_s = np.ones_like(act_mean)

        for alpha in np.linspace(0.0, 1.0, self.n_grid):
            s_cand = act_mean ** alpha
            geom_mean = np.exp(np.mean(np.log(s_cand)))
            s_cand = s_cand / max(geom_mean, 1e-6)

            # Transformed weights: W' = W * diag(s)
            w_scaled = w * s_cand[np.newaxis, :]
            w_scaled_q = self.quantize_symmetric(w_scaled)
            # Reconstruct effective weight: W_recon = W_scaled_q * diag(s)^-1
            w_recon = w_scaled_q / s_cand[np.newaxis, :]

            y_pred = np.matmul(w_recon, x)
            loss = float(np.mean((y_true - y_pred) ** 2))

            if loss < best_loss:
                best_loss = loss
                best_alpha = float(alpha)
                best_s = s_cand.copy()

        return best_s, best_alpha, best_loss

    @staticmethod
    def fold_scale_into_layernorm(ln_weight: np.ndarray, s: np.ndarray) -> np.ndarray:
        """
        Folds the AWQ scale vector s into preceding LayerNorm weights:
        W_LN_new = W_LN / s
        Ensuring zero runtime latency overhead at inference time!
        """
        return ln_weight / s
\end{lstlisting}

\begin{lstlisting}[language=Python, caption={Tầng tuyến tính DoRA phân rã độc lập Độ lớn và Hướng (Liu et al., ICML 2024).}]
class DoRALinearLayer:
    """
    Weight-Decomposed Low-Rank Adaptation (DoRA, Liu et al., ICML 2024).
    Decouples weight magnitude m and direction V:
    W = m * (W_0 + (alpha / r) * A @ B) / ||W_0 + (alpha / r) * A @ B||_c
    Where ||.||_c is the column-wise L2 norm over input dimensions.
    """
    def __init__(self, d_in: int, d_out: int, r: int = 8, alpha: float = 16.0):
        self.d_in = d_in
        self.d_out = d_out
        self.r = r
        self.alpha = alpha
        self.scaling = alpha / r

        # Base weights (frozen)
        self.w0 = np.random.standard_normal((d_in, d_out)).astype(np.float32)

        # Trainable magnitude vector m: initialized to column-wise L2 norm of W0
        self.m = np.linalg.norm(self.w0, axis=0, keepdims=True)

        # Trainable low-rank direction matrices
        self.a = np.random.standard_normal((d_in, r)).astype(np.float32) * 0.01
        self.b = np.zeros((r, d_out), dtype=np.float32)

    def get_decomposed_weight(self) -> np.ndarray:
        delta_w = self.scaling * np.matmul(self.a, self.b)
        v = self.w0 + delta_w
        v_norm = np.linalg.norm(v, axis=0, keepdims=True)
        v_norm = np.maximum(v_norm, 1e-8)
        return self.m * (v / v_norm)

    def forward(self, x: np.ndarray) -> np.ndarray:
        w_dora = self.get_decomposed_weight()
        return np.matmul(x, w_dora)

    def merge_weights(self) -> np.ndarray:
        """Merges DoRA direction and magnitude into static base weight."""
        return self.get_decomposed_weight()
\end{lstlisting}

\begin{lstlisting}[language=Python, caption={Bộ tính toán giải tích dung lượng và mức tiết kiệm của QLoRA Double Quantization.}]
class QLoRADoubleQuantCalculator:
    """
    Calculates exact memory footprint and savings of QLoRA Double Quantization (DQ).
    Weights: 4-bit (NF4). Block size B1 (default 64).
    First-level scales: FP32 (32 bits) without DQ -> 32 / B1 = 0.50 bits/param.
    With Double Quantization (DQ):
    First-level scales quantized to FP8 (8 bits) with block size B2 (default 256).
    Second-level scales in FP32 (32 bits).
    Memory per param for scales: 8 / B1 + 32 / (B1 * B2)
    Total saving: 32 / B1 - (8 / B1 + 32 / (B1 * B2)) = 0.373 bits/param.
    """
    def __init__(self, block_size_b1: int = 64, block_size_b2: int = 256):
        self.b1 = block_size_b1
        self.b2 = block_size_b2

    def calculate_bits_per_param(self) -> dict:
        weight_bits = 4.0
        single_quant_scale_bits = 32.0 / self.b1
        single_quant_total = weight_bits + single_quant_scale_bits

        dq_c1_scale_bits = 8.0 / self.b1
        dq_c2_scale_bits = 32.0 / (self.b1 * self.b2)
        dq_total_scale_bits = dq_c1_scale_bits + dq_c2_scale_bits
        dq_total = weight_bits + dq_total_scale_bits

        savings_per_param = single_quant_scale_bits - dq_total_scale_bits
        return {
            "weight_bits": weight_bits,
            "single_quant_scale_bits": single_quant_scale_bits,
            "single_quant_total_bits": single_quant_total,
            "double_quant_scale_bits": dq_total_scale_bits,
            "double_quant_total_bits": dq_total,
            "savings_bits_per_param": savings_per_param,
        }

    def calculate_model_vram_gb(
        self,
        num_parameters: int,
        include_lora: bool = True,
        lora_params: int = 100_000_000
    ) -> dict:
        bits_info = self.calculate_bits_per_param()
        bytes_per_gb = 1024 ** 3

        fp16_bytes = num_parameters * 2.0
        fp16_gb = fp16_bytes / bytes_per_gb

        sq_bytes = num_parameters * (bits_info["single_quant_total_bits"] / 8.0)
        sq_gb = sq_bytes / bytes_per_gb

        dq_bytes = num_parameters * (bits_info["double_quant_total_bits"] / 8.0)
        dq_gb = dq_bytes / bytes_per_gb

        savings_bytes = num_parameters * (bits_info["savings_bits_per_param"] / 8.0)
        savings_gb = savings_bytes / bytes_per_gb

        lora_gb = 0.0
        if include_lora:
            # LoRA weights (BF16, 2B) + Gradients (BF16, 2B) + AdamW (FP32, 8B) = 12B/param
            lora_bytes = lora_params * 12.0
            lora_gb = lora_bytes / bytes_per_gb

        return {
            "fp16_gb": fp16_gb,
            "single_quant_gb": sq_gb,
            "double_quant_gb": dq_gb,
            "vram_saved_gb": savings_gb,
            "lora_overhead_gb": lora_gb,
            "total_qlora_training_gb": dq_gb + lora_gb,
        }
\end{lstlisting}

\begin{simulationbox}[Mô Phỏng Trực Quan Hóa Thuật Toán Lượng Tử Hóa, LoRA \& Alignment]
Để củng cố trực giác vật lý và kiến trúc phần cứng của các đột phá công nghệ trong chương này, dưới đây là kịch bản và prompt chuẩn để yêu cầu AI sinh mã nguồn hoạt họa 60 FPS bằng thư viện Manim Community Edition (Manim CE):

\paragraph{Prompt 1: Hoạt Họa 60 FPS Cơ Chế Chuẩn Hóa Nhóm GRPO (DeepSeek-R1)}
\begin{quote}\small\itshape
``Hãy viết một đoạn mã Python hoàn chỉnh sử dụng thư viện Manim CE (\code{from manim import *}) để tạo một video giải thích trực quan về thuật toán GRPO (Group Relative Policy Optimization) từ DeepSeek-R1 với các yêu cầu sau:
\begin{enumerate}
  \item \textbf{Cảnh 1: Lấy Mẫu Nhóm (Group Rollout Sampling)}: Thể hiện một nút câu hỏi trung tâm (Prompt Node: $q$) phát ra đồng thời $G = 4$ nhánh suy luận chạy song song. Mỗi nhánh hiển thị cấu trúc thẻ \verb|<think>| đang mở rộng độ dài động.
  \item \textbf{Cảnh 2: Đánh Giá Quy Tắc Tất Định (Rule-Based Verifiers)}: Các nhánh đi qua hai khối kiểm định nhị phân: Khối kiểm tra định dạng (\verb|<think>| \& \verb|\boxed{}|) và Khối kiểm tra đáp án SymPy. Hiển thị điểm số $r_1=1.5$, $r_2=1.0$, $r_3=0.5$, $r_4=0.0$ với màu sắc tương ứng (Xanh lá / Vàng / Đỏ).
  \item \textbf{Cảnh 3: Chuẩn Hóa Nhóm (Group Normalization \& Baseline Shift)}: Xuất hiện một biểu đồ cột hiển thị mức điểm trung bình $\bar{r} = 0.75$ dưới dạng một đường nét đứt màu trắng. Sau đó, các cột biến đổi chuyển dịch vị trí: các nhánh trên đường cơ sở nâng lên thành $\widetilde{A} > 0$ (màu xanh dương), các nhánh dưới đường cơ sở hạ xuống thành $\widetilde{A} < 0$ (màu cam).
  \item \textbf{Cảnh 4: Cập Nhật Chính Sách}: Thể hiện các mũi tên gradient kéo xác suất của nhánh $o_1$ lên cao và dìm nhánh $o_4$ xuống.
  \item Sử dụng font Serif cổ điển, bố cục chia vùng rõ ràng và có thể render trực tiếp bằng lệnh: \code{manim -pql scene.py GRPOGroupAdvantageScene}.''
\end{enumerate}
\end{quote}

\paragraph{Prompt 2: Hoạt Họa 60 FPS Kiến Trúc Mạch Bán Dẫn Cộng Số Nguyên BitNet b1.58}
\begin{quote}\small\itshape
``Hãy viết một kịch bản Manim CE 60 FPS so sánh trực quan kiến trúc phần cứng vi mạch giữa phép nhân ma trận FP16 truyền thống và bộ cộng tích lũy số nguyên BitNet b1.58:
\begin{enumerate}
  \item \textbf{Cảnh 1: Màn Hình Phân Đôi (Split Screen Layout)}:
  \begin{itemize}
    \item Bên trái: Kiến trúc FP16 GEMM truyền thống. Thể hiện một mảng ma trận các bộ nhân số thực (Multiplier Array) cồng kềnh, các vạch nhiệt lượng màu đỏ tỏa ra và đồng hồ công suất tiêu thụ ghi nhận $1.10\text{ pJ/op}$.
    \item Bên phải: Kiến trúc BitNet b1.58. Trọng số chỉ mang ba trạng thái $\{-1, 0, +1\}$. Các bộ nhân bị mờ đi và biến mất hoàn toàn, thay bằng một cây bộ cộng nhị phân (Adder Tree) siêu nhỏ gọn, đồng hồ tiêu thụ chỉ $0.03\text{ pJ/op}$ (tiết kiệm $36\times$).
  \end{itemize}
  \item \textbf{Cảnh 2: Luồng Dữ Liệu Thực Thi (Dataflow in Action)}: Dòng token kích hoạt INT8 đi vào. Các phần tử gặp trọng số $+1$ được dẫn thẳng vào cổng cộng (Green wire), phần tử gặp trọng số $-1$ được dẫn vào cổng trừ (Red wire), phần tử gặp $0$ bị ngắt mạch (Disconnected wire), chứng minh không có bất kỳ phép nhân dấu phẩy động nào diễn ra.
  \item \textbf{Cảnh 3: Biểu Đồ Hiệu Suất Băng Thông}: Mô phỏng việc nén mô hình 70B từ $140\text{ GB}$ VRAM xuống còn $14\text{ GB}$, chạy mượt mà trên bộ nhớ thống nhất (Unified Memory) của Apple M-series hoặc GPU tiêu dùng.
  \item Đảm bảo mã nguồn chuẩn mực, sử dụng bảng màu Navy/Amber/Teal và render trực tiếp với lệnh: \code{manim -pql scene.py BitNetSiliconAdderScene}.''
\end{enumerate}
\end{quote}

\paragraph{Prompt 3: Hoạt Họa 60 FPS Tối Ưu Hóa AWQ \& Gập Hệ Số LayerNorm Folding}
\begin{quote}\small\itshape
``Hãy viết một đoạn mã Python sử dụng thư viện Manim CE để tạo một video giải thích trực quan về thuật toán AWQ (Activation-aware Weight Quantization) và cơ chế Zero-Latency LayerNorm Folding:
\begin{enumerate}
  \item \textbf{Cảnh 1: Phân Phối Kích Hoạt \& Kênh Ngoại Lai (Activation Outliers)}: Thể hiện một ma trận kích hoạt $\mathbf{X}$ chạy qua mạng. Zoom cận cảnh vào kênh thứ $0$ và $1$ với các cột kích hoạt cao vọt lên (màu đỏ phát sáng), lớn gấp $50\times$ so với các kênh thông thường (màu xám nhạt).
  \item \textbf{Cảnh 2: Cơ Chế Co Giãn Bảo Vệ (Protection Scaling)}: Hiển thị ma trận trọng số $\mathbf{W}$. Khi áp dụng hệ số $s > 1$, các cột trọng số kết nối với kênh ngoại lai được kéo dãn biên độ lên $s$ lần. Minh họa bằng lưới phân vị INT4: khi biên độ cột trọng số tăng cao, sai số làm tròn tương đối $\Delta w / w$ co nhỏ lại rõ rệt.
  \item \textbf{Cảnh 3: Gập Tỷ Lệ Vào LayerNorm (Zero-Latency Folding)}: Xuất hiện khối LayerNorm trước khối tuyến tính. Thể hiện vector tỷ lệ $\mathbf{s}^{-1}$ bay từ phép nhân ma trận ngược trở lại khối LayerNorm và hợp nhất trực tiếp vào vector trọng số $\boldsymbol{\gamma}_{\text{LN}}$, biến đổi thành $\boldsymbol{\gamma}_{\text{LN}}' = \boldsymbol{\gamma}_{\text{LN}} / \mathbf{s}$. Một đồng hồ đếm độ trễ bổ sung hiển thị con số $0\text{ ns}$ phát sáng màu xanh lá.
  \item Render trực tiếp bằng lệnh: \code{manim -pql scene.py AWQLayerNormFoldingScene}.''
\end{enumerate}
\end{quote}

\paragraph{Prompt 4: Hoạt Họa 60 FPS Phân Rã Độ Lớn \& Hướng DoRA So Với LoRA Truyền Thống}
\begin{quote}\small\itshape
``Hãy viết một kịch bản Manim CE 60 FPS minh họa hình học sự khác biệt giữa LoRA và DoRA trên không gian vector 2 chiều và mặt cầu chuẩn hóa:
\begin{enumerate}
  \item \textbf{Cảnh 1: Vector Trọng Số Gốc $\mathbf{w}_0$}: Vẽ một vector gốc màu xanh navy bắt đầu từ tâm tọa độ và một đường tròn nét đứt biểu diễn mặt cầu chuẩn hóa $\|\mathbf{w}\|_2 = m_0$.
  \item \textbf{Cảnh 2: Hạn Chế Ghép Đôi Của LoRA (Coupled Dynamics)}: Thể hiện cập nhật LoRA $\Delta \mathbf{w} = \frac{\alpha}{r}\mathbf{A}\mathbf{B}$. Vector kết quả $\mathbf{w}_{\text{LoRA}}$ vừa bị kéo dài độ lớn vừa bị đổi hướng theo một quỹ đạo dính chéo cứng nhắc. Đồ thị phân tán $\Delta m$ vs $\Delta \theta$ hiển thị một đường thẳng hồi quy dốc đứng (hệ số tương quan dương $r \approx 0.85$).
  \item \textbf{Cảnh 3: Phân Rã Độc Lập DoRA (Decoupled Freedom)}:
  \begin{itemize}
    \item Bước 1: Vector hướng $\mathbf{V}$ trượt linh hoạt trên mặt cầu đơn vị nhờ ma trận chiếu trực giao $\mathbf{P}_{\mathbf{v}} = \mathbf{I} - \frac{\mathbf{v}\mathbf{v}^T}{\|\mathbf{v}\|^2}$.
    \item Bước 2: Vector độ lớn $\mathbf{m}$ co dãn tự do độ dài vector mà không làm đổi góc hướng.
    \item Đồ thị $\Delta m$ vs $\Delta \theta$ phân tán tự do giống hệt đám mây điểm của Full Fine-Tuning.
  \end{itemize}
  \item Render trực tiếp với lệnh: \code{manim -pql scene.py DoRAOrthogonalDecompositionScene}.''
\end{enumerate}
\end{quote}
\end{simulationbox}

---

\section{Hệ Thống Bài Tập \& Câu Hỏi Tư Duy Đa Tầng}
\label{sec:ch11-exercises}

\subsection{Nhóm A: Câu Hỏi Tư Duy Bản Chất (Conceptual Reasoning)}
\begin{enumerate}
  \item \textbf{Câu A.1}: Tại sao việc giữ ma trận $\mathbf{B}$ của LoRA khởi tạo bằng $0.0$ tại thời điểm bắt đầu huấn luyện ($t=0$) lại là yêu cầu kỹ thuật bắt buộc thay vì khởi tạo ngẫu nhiên như mạng nơ-ron truyền thống?
  \item \textbf{Câu A.2}: Phân tích sự khác biệt về mặt bản chất giữa FP16 và BF16. Tại sao việc cắt giảm độ chính xác định trị (từ 10 bits xuống 7 bits) của BF16 lại được chấp nhận rộng rãi trong huấn luyện mô hình ngôn ngữ lớn?
  \item \textbf{Câu A.3}: Tại sao thuật toán AWQ lại bảo toàn được năng lực giải thuật và viết code của mô hình sau khi nén INT4 tốt hơn đáng kể so với phương pháp làm tròn độc lập (Round-To-Nearest -- RTN)?
  \item \textbf{Câu A.4}: Trong bài toán căn chỉnh DPO, điều gì sẽ xảy ra nếu ta đặt hệ số phạt $\beta \to 0$ hoặc $\beta \to \infty$?
  \item \textbf{Câu A.5}: So sánh ưu và nhược điểm kiến trúc giữa DPO và ORPO. Khi nào một Kỹ sư AI nên chọn ORPO thay vì DPO trong dự án thực tế?
  \item \textbf{Câu A.6 (Bản Chất Của Hessian Nghịch Đảo Trong Optimal Brain Surgeon)}: \\
  Trong công thức gia tăng sai số của OBS:
  \begin{equation*}
  \Delta \mathcal{E}_q = \frac{1}{2} \frac{(w_q - \hat{w}_q)^2}{[\mathbf{H}^{-1}]_{qq}}
  \end{equation*}
  Tại sao phần tử đường chéo $[\mathbf{H}^{-1}]_{qq}$ của ma trận nghịch đảo càng lớn thì mức gia tăng sai số khi lượng tử hóa trọng số $w_q$ lại càng nhỏ? Giải thích ý nghĩa hình học của elip đường đồng mức (Level Sets) của hàm mất mát trong không gian tham số.
  \item \textbf{Câu A.7 (So Sánh Độ Phức Tạp Thuật Toán Nén Trọng Số)}: \\
  So sánh độ phức tạp tính toán và bộ nhớ giữa ba chiến lược nén trọng số: Làm tròn độc lập (Round-To-Nearest -- RTN), Optimal Brain Surgeon (OBS) cổ điển cập nhật từng trọng số, và thuật toán GPTQ dùng phân rã Cholesky cập nhật theo khối (Lazy Block Updates). Tại sao OBS cổ điển không thể mở rộng cho các mô hình ngôn ngữ lớn như Llama-3-70B, trong khi GPTQ lại hoàn thành trong vòng vài giờ?

  \item \textbf{Câu A.8 (Cơ Chế Khử Critic Trong Thuật Toán GRPO)}: \\
  Tại sao thuật toán GRPO có thể triệt tiêu hoàn toàn mạng Critic $V_\phi$ mà vẫn đảm bảo tính ổn định hội tụ của thuật toán PPO? Phân tích cách việc chuẩn hóa điểm thưởng trong nhóm (Group Advantage) đóng vai trò như một bộ ước lượng đường cơ sở (Baseline Estimator) thích ứng động theo từng câu hỏi.

  \item \textbf{Câu A.9 (Ý Nghĩa Của Trọng Số Bậc Ba Trong BitNet b1.58)}: \\
  So sánh BitNet 1-bit nhị phân thuần túy ($\{-1, +1\}$) và BitNet b1.58 bậc ba ($\{-1, 0, +1\}$). Tại sao việc bổ sung trạng thái $0$ lại là bước ngoặt quyết định giúp mô hình đạt hiệu năng tương đương FP16 ở cùng quy mô tham số? Phân tích theo lý thuyết lọc tính năng (Feature Filtering) và tính thưa động (Dynamic Sparsity).

  \item \textbf{Câu A.10 (Bản Chất LayerNorm Folding Trong Thuật Toán AWQ)}: \\
  Tại sao trong thuật toán AWQ, việc nhân ma trận co giãn $\text{diag}(\mathbf{s})^{-1}$ vào kích hoạt $\mathbf{X}$ tại thời gian chạy (Runtime) lại gây ra hiện tượng nghẽn băng thông bộ nhớ (Memory Bandwidth Bottleneck)? Trình bày cơ chế toán học của kỹ thuật gập trực tiếp (LayerNorm Folding) vào trọng số của tầng chuẩn hóa đi trước và giải thích tại sao kỹ thuật này triệt tiêu hoàn toàn độ trễ bổ sung ($0\text{ ns}$).

  \item \textbf{Câu A.11 (Nguyên Lý Phân Rã Độ Lớn \& Hướng Của DoRA)}: \\
  So sánh cơ chế cập nhật tham số giữa LoRA truyền thống và DoRA. Bằng cách nào việc phân rã ma trận trọng số $\mathbf{W} = \mathbf{m} \odot \frac{\mathbf{V}}{\|\mathbf{V}\|_c}$ kết hợp với ma trận chiếu trực giao $\mathbf{P}_{\mathbf{v}} = \mathbf{I} - \frac{\mathbf{v}\mathbf{v}^T}{\|\mathbf{v}\|_2^2}$ giúp loại bỏ sự ghép đôi cứng nhắc (Coupled Dynamics) giữa độ lớn và góc hướng?
\end{enumerate}

\subsection{Nhóm B: Bài Toán Tính Toán Giải Tích Bằng Tay (Analytical Problems)}
\begin{enumerate}
  \item \textbf{Bài B.1 (Tính Định Cỡ VRAM Cho QLoRA 70B Trên GPU Tiêu Dùng)}: \\
  Xét mô hình Llama-3-70B có 70 tỷ tham số.
  Chúng ta tinh chỉnh mô hình bằng QLoRA:
  Trọng số gốc được lượng tử hóa sang NF4 ($0.5\text{ byte/tham số}$).
  Adapter LoRA có $100\text{ triệu}$ tham số ở định dạng BF16 ($2\text{ bytes/tham số}$).
  \begin{itemize}
    \item[a)] Tính dung lượng VRAM để lưu trọng số gốc dạng NF4.
    \item[b)] Tính dung lượng VRAM để lưu trọng số LoRA, gradients của LoRA, và trạng thái Optimizer AdamW của LoRA (giả sử trạng thái AdamW chiếm 8 bytes cho mỗi tham số LoRA).
    \item[c)] Cộng thêm $4\text{ GB}$ bộ đệm kích hoạt (Activation memory) cho batch size nhỏ, kiểm tra xem toàn bộ quá trình tinh chỉnh có thể chạy vừa trên 2 card đồ họa RTX 3090/4090 (mỗi card $24\text{ GB}$, tổng $48\text{ GB}$) hay không.
  \end{itemize}

  \item \textbf{Bài B.2 (Phương Sai Sai Số Lượng Tử Hóa Tuyến Tính)}: \\
  Giả sử biến ngẫu nhiên kích hoạt $X$ tuân theo phân phối đều liên tục $X \sim \mathcal{U}[-a, a]$.
  Chúng ta thực hiện lượng tử hóa đối xứng sang $K$ mức rời rạc với bước nhảy $\Delta = \frac{2a}{K}$.
  Chứng minh rằng sai số làm tròn $e = X - \hat{X}$ có kỳ vọng $E[e] = 0$ và phương sai sai số là:
  \begin{equation*}
  \text{Var}(e) = \frac{\Delta^2}{12}
  \end{equation*}

  \item \textbf{Bài B.3 (Chứng Minh Giải Tích Đạo Hàm Gradient Của DPO)}: \\
  Cho hàm mất mát DPO:
  \begin{equation*}
  \mathcal{L}_{\text{DPO}}(\theta) = -\log \sigma\left( \beta \log \frac{\pi_\theta(y_w \mid x)}{\pi_{\text{ref}}(y_w \mid x)} - \beta \log \frac{\pi_\theta(y_l \mid x)}{\pi_{\text{ref}}(y_l \mid x)} \right)
  \end{equation*}
  Bằng cách đặt $u = \beta \log \frac{\pi_\theta(y_w \mid x)}{\pi_{\text{ref}}(y_w \mid x)} - \beta \log \frac{\pi_\theta(y_l \mid x)}{\pi_{\text{ref}}(y_l \mid x)}$ và áp dụng tính chất của hàm sigmoid $\sigma'(u) = \sigma(u)(1 - \sigma(u)) = \sigma(u)\sigma(-u)$, hãy chứng minh rằng:
  \begin{equation*}
  \nabla_\theta \mathcal{L}_{\text{DPO}}(\theta) = -\beta \cdot \sigma\left( \hat{r}_\theta(x, y_l) - \hat{r}_\theta(x, y_w) \right) \cdot \left[ \nabla_\theta \log \pi_\theta(y_w \mid x) - \nabla_\theta \log \pi_\theta(y_l \mid x) \right]
  \end{equation*}

  \item \textbf{Bài B.4 (Tính Toán Thủ Công Bù Sai Số OBS Cho Tensor $2 \times 3$)}: \\
  Cho ma trận kích hoạt $\mathbf{X} = \begin{bmatrix} 1 & 2 & 1 \\ 0 & 1 & 2 \end{bmatrix}$ và vector trọng số ban đầu $\mathbf{w} = [1.5, -2.0]$.
  \begin{enumerate}
    \item Tính ma trận Hessian $\mathbf{H} = 2 \mathbf{X} \mathbf{X}^T$ và ma trận Hessian nghịch đảo $\mathbf{H}^{-1}$.
    \item Nếu lượng tử hóa trọng số thứ nhất từ $w_1 = 1.5$ thành $\hat{w}_1 = 2.0$, hãy áp dụng công thức OBS để tính vector bù sai số $\Delta \mathbf{w}$ cho trọng số thứ hai $w_2$.
    \item Tính giá trị mới của $w_2'$ và so sánh bình phương sai số kích hoạt đầu ra $\|\mathbf{y} - \hat{\mathbf{y}}\|_2^2$ giữa phương pháp OBS và phương pháp làm tròn độc lập RTN.
  \end{enumerate}

  \item \textbf{Bài B.5 (Chứng Minh Giải Tích Đầy Đủ Công Thức Cập Nhật OBS Bằng Nhân Tử Lagrange)}: \\
  Cho hàm sai số bậc hai $\mathcal{E}(\Delta \mathbf{w}) = \frac{1}{2} \Delta \mathbf{w}^T \mathbf{H} \Delta \mathbf{w}$ với $\mathbf{H}$ là ma trận đối xứng xác định dương. Xét bài toán tối thiểu hóa sai số khi cố định lượng tử hóa trọng số thứ $q$:
  \begin{equation*}
  \min_{\Delta \mathbf{w}} \frac{1}{2} \Delta \mathbf{w}^T \mathbf{H} \Delta \mathbf{w} \quad \text{s.t.} \quad \mathbf{e}_q^T \Delta \mathbf{w} + (w_q - \hat{w}_q) = 0
  \end{equation*}
  Bằng phương pháp nhân tử Lagrange, hãy thiết lập hàm Lagrange $\mathcal{L}(\Delta \mathbf{w}, \lambda)$ và giải tích chi tiết để chứng minh:
  \begin{equation*}
  \Delta \mathbf{w} = - \frac{w_q - \hat{w}_q}{[\mathbf{H}^{-1}]_{qq}} \cdot \mathbf{H}^{-1}_{:, q} \qquad \text{và} \qquad \Delta \mathcal{E}_q = \frac{1}{2} \frac{(w_q - \hat{w}_q)^2}{[\mathbf{H}^{-1}]_{qq}}
  \end{equation*}

  \item \textbf{Bài B.6 (Chứng Minh Độ Lợi Thế Nhóm Trong GRPO Luôn Có Kỳ Vọng Mẫu Bằng 0)}: \\
  Cho một nhóm gồm $G$ câu trả lời độc lập với điểm thưởng $r_1, r_2, \dots, r_G$. Độ lợi thế chuẩn hóa nhóm được xác định:
  \begin{equation*}
  \widetilde{A}_i = \frac{r_i - \bar{r}}{\sigma_r + \epsilon}, \qquad \text{với } \bar{r} = \frac{1}{G}\sum_{i=1}^G r_i, \quad \sigma_r = \sqrt{\frac{1}{G}\sum_{i=1}^G (r_i - \bar{r})^2}
  \end{equation*}
  \begin{enumerate}
    \item Chứng minh giải tích rằng tổng $\sum_{i=1}^G \widetilde{A}_i = 0$ và suy ra kỳ vọng mẫu của độ lợi thế nhóm luôn triệt tiêu về $0$.
    \item Xét trường hợp $\sigma_r \to 0$ (tất cả các câu trả lời đều cùng đạt điểm thưởng bằng nhau). Phân tích ý nghĩa học tăng cường tại sao độ lợi thế triệt tiêu về $0$ lại ngăn ngừa việc cập nhật gradient dư thừa khi câu hỏi quá dễ hoặc quá khó.
  \end{enumerate}

  \item \textbf{Bài B.7 (Tính Toán Dung Lượng Thông Tin Shannon \& Đóng Gói Bit-packing Cho BitNet b1.58)}: \\
  \begin{enumerate}
    \item Chứng minh rằng dung lượng thông tin lý thuyết của một biến ngẫu nhiên nhận 3 giá trị đẳng xác suất là $\log_2(3) \approx 1.58496\text{ bits}$.
    \item Chứng minh rằng một byte ($8\text{ bits}$) có thể chứa tối đa 5 tham số bậc ba độc lập. Tính dung lượng bộ nhớ thực tế (theo đơn vị GiB) cần thiết để lưu trữ toàn bộ ma trận trọng số của mô hình ngôn ngữ 70 tỷ tham số khi được đóng gói bằng phương pháp này, và so sánh với mô hình FP16 gốc.
  \end{enumerate}

  \item \textbf{Bài B.8 (Tính Toán Giải Tích Mức Tiết Kiệm VRAM Của QLoRA Double Quantization)}: \\
  Xét hai mô hình ngôn ngữ mã nguồn mở hàng đầu: Llama-3-70B ($70 \times 10^9$ tham số) và Llama-3.1-405B ($405 \times 10^9$ tham số).
  Trọng số 4-bit được chia thành các khối kích thước $B_1 = 64$.
  \begin{enumerate}
    \item Tính dung lượng hệ số co giãn cấp 1 (FP32) trên từng tham số khi không áp dụng Double Quantization.
    \item Khi áp dụng Double Quantization với block size cấp hai $B_2 = 256$, hệ số $c_1$ được lượng tử hóa sang FP8 ($8\text{ bits}$) và hệ số $c_2$ giữ dạng FP32 ($32\text{ bits}$). Tính dung lượng scale trung bình trên từng tham số và chứng minh rằng mức tiết kiệm là $\Delta b \approx 0.373\text{ bits/tham số}$.
    \item Tính chính xác dung lượng VRAM tiết kiệm được (theo đơn vị GiB) cho cả hai mô hình 70B và 405B.
  \end{enumerate}

  \item \textbf{Bài B.9 (Chứng Minh Giải Tích Tính Trực Giao Của Gradient Hướng DoRA)}: \\
  Cho một vector cột trọng số $\mathbf{w} = m \frac{\mathbf{v}}{\|\mathbf{v}\|_2}$ với $m = \|\mathbf{w}\|_2 > 0$ và $\mathbf{v} \neq \mathbf{0}$.
  Gọi $\mathbf{g} = \nabla_{\mathbf{w}} \mathcal{L}$ là gradient của hàm mất mát đối với vector trọng số hiệu dụng $\mathbf{w}$.
  \begin{enumerate}
    \item Chứng minh rằng gradient đối với vector hướng $\mathbf{v}$ có dạng:
    \begin{equation*}
    \nabla_{\mathbf{v}} \mathcal{L} = \frac{m}{\|\mathbf{v}\|_2} \left( \mathbf{I} - \frac{\mathbf{v}\mathbf{v}^T}{\|\mathbf{v}\|_2^2} \right) \mathbf{g}
    \end{equation*}
    \item Chứng minh rằng tích vô hướng $\mathbf{v}^T \nabla_{\mathbf{v}} \mathcal{L} = 0$, khẳng định rằng gradient cập nhật cho vector hướng luôn trực giao hoàn toàn với vector hướng hiện tại.
  \end{enumerate}
\end{enumerate}

\subsection{Nhóm C: Thực Hành Kỹ Nghệ \& Lập Trình (Implementation Labs)}
\begin{enumerate}
  \item \textbf{Bài C.1 (Bộ Quản Lý Đa Adapter LoRA Cho Nền Tảng SaaS)}: \\
  Xây dựng lớp Python \code{MultiLoRAManager} duy trì một mô hình nền tảng gốc duy nhất $\mathbf{W}_0$ cố định trên GPU VRAM và hỗ trợ nạp động nhiều bộ adapter LoRA cho từng khách hàng (\code{tenant_id}), tự động kích hoạt cặp ma trận $(\mathbf{A}, \mathbf{B})$ tương ứng mà không làm nhân bản mô hình gốc.

  \item \textbf{Bài C.2 (Bộ Chuyển Đổi SmoothQuant Tự Động)}: \\
  Viết hàm \code{apply_smoothquant_layer} nhận vào tensor kích hoạt $\mathbf{X}$ và ma trận trọng số $\mathbf{W}$, tính toán véc-tơ tỷ lệ $\mathbf{s}$, thực hiện biến đổi cân bằng và chứng minh rằng sai số cực đại tuyệt đối $\|\mathbf{X}\mathbf{W} - \hat{\mathbf{X}}\hat{\mathbf{W}}\|_\infty < 10^{-5}$.

  \item \textbf{Bài C.3 (Bộ Kiểm Định Sở Thích DPO)}: \\
  Hiện thực hóa bộ kiểm định tự động đo lường độ chính xác phân loại sở thích (Reward Accuracy) và tỷ lệ lề phần thưởng trung bình qua 100 mẫu huấn luyện.

  \item \textbf{Bài C.4 (Bộ Lượng Tử Hóa OBS/GPTQ Một Tầng Với Cholesky)}: \\
  Xây dựng lớp Python \code{OBSQuantizerEngine} nhận ma trận kích hoạt $\mathbf{X}$ và ma trận trọng số $\mathbf{W}$. Thực hiện:
  \begin{itemize}
    \item Tính ma trận Hessian $\mathbf{H} = 2 \mathbf{X} \mathbf{X}^T$ với hệ số điều hòa $\lambda \mathbf{I}$ ($\lambda = 0.01 \times \text{mean}(\text{diag}(\mathbf{H}))$).
    \item Tính nghịch đảo Hessian $\mathbf{H}^{-1}$.
    \item Tuần tự duyệt qua từng cột của $\mathbf{W}$, lượng tử hóa sang số nguyên đối xứng INT8/INT4 và bù trừ sai số giải tích lên các cột còn lại.
    \item So sánh tổng sai số bình phương tái tạo kích hoạt MSE với phương pháp làm tròn độc lập Round-to-Nearest (RTN) và kiểm định rằng MSE của OBS luôn nhỏ hơn RTN ít nhất $40\%$.
  \end{itemize}

  \item \textbf{Bài C.5 (Bộ Ước Lượng Lợi Thế Nhóm GRPO \& Kiểm Định Quy Tắc)}: \\
  Xây dựng lớp Python \code{GRPOAdvantageEstimator} hiện thực hóa các phương thức:
  \begin{itemize}
    \item Bộ kiểm định quy tắc \code{rule_based_verifier} (chấm điểm thẻ \verb|<think>| và \verb|\boxed{}|).
    \item Hàm chuẩn hóa độ lợi thế nhóm \code{compute_group_advantages}.
    \item Hàm tính tổn thất cắt xén PPO \code{compute_clipped_surrogate_loss}.
  \end{itemize}

  \item \textbf{Bài C.6 (Lớp Tuyến Tính BitNet b1.58 Thuần Phép Cộng Số Nguyên)}: \\
  Xây dựng lớp Python \code{BitNetTernaryLinear} thực hiện lượng tử hóa trọng số bậc ba theo Absmean, lượng tử hóa kích hoạt INT8, và phép lan truyền xuôi \code{forward_pure_addition} hoàn toàn không sử dụng phép nhân số thực, kiểm chứng tính tương đương chính xác với phép nhân ma trận số nguyên chuẩn.

  \item \textbf{Bài C.7 (Bộ Tối Ưu Hóa Tỷ Lệ AWQScaleOptimizer \& Gập Vào LayerNorm)}: \\
  Xây dựng lớp Python \code{AWQScaleOptimizer} thực hiện:
  \begin{itemize}
    \item Đo đạc biên độ trung bình kích hoạt theo từng kênh $s_{X, j} = \text{mean}(|X_{j, :}|)$.
    \item Quét lưới siêu tham số $\alpha \in [0, 1]$ với 20 điểm mẫu để tìm tỷ lệ $\mathbf{s} = \mathbf{s}_X^\alpha$ tối thiểu hóa sai số $\|\mathbf{W}\mathbf{X} - Q(\mathbf{W} \cdot \text{diag}(\mathbf{s})) \cdot \text{diag}(\mathbf{s})^{-1}\mathbf{X}\|_F^2$.
    \item Hiện thực phương thức \code{fold_scale_into_layernorm} gập trực tiếp tỷ lệ $\mathbf{s}$ vào trọng số LayerNorm đi trước.
  \end{itemize}

  \item \textbf{Bài C.8 (Tầng Tuyến Tính DoRALinearLayer Phân Rã Độ Lớn \& Hướng)}: \\
  Xây dựng lớp Python \code{DoRALinearLayer} mô phỏng tầng tuyến tính:
  \begin{itemize}
    \item Khởi tạo vector độ lớn $\mathbf{m} = \|\mathbf{W}_0\|_c$ và hai ma trận LoRA $\mathbf{A}, \mathbf{B}$ với $\mathbf{B} = \mathbf{0}$.
    \item Hiện thực hóa lan truyền xuôi theo công thức $\mathbf{W} = \mathbf{m} \odot \frac{\mathbf{W}_0 + \frac{\alpha}{r}\mathbf{A}\mathbf{B}}{\|\mathbf{W}_0 + \frac{\alpha}{r}\mathbf{A}\mathbf{B}\|_c}$.
    \item Hiện thực phương thức \code{merge_weights} hợp nhất vĩnh viễn trọng số phục vụ serving zero-latency.
  \end{itemize}

  \item \textbf{Bài C.9 (Bộ Tính Toán QLoRADoubleQuantCalculator Định Cỡ Bộ Nhớ)}: \\
  Xây dựng lớp Python \code{QLoRADoubleQuantCalculator} tính toán chính xác số bit trên tham số và phân tích chi tiết dung lượng VRAM tĩnh, bộ nhớ LoRA (trọng số, gradient, AdamW) và bộ đệm kích hoạt cho các mô hình $7\text{B}, 13\text{B}, 70\text{B}, 405\text{B}$.
\end{enumerate}

\subsection{Nhóm D: Phân Tích Thất Bại \& Debugging Thực Tế (Post-Mortem Analysis)}
\begin{enumerate}
  \item \textbf{Tình huống sự cố 1 (Suy thoái năng lực giải thuật sau khi lượng tử hóa)}: \\
  Đội ngũ AI tiến hành lượng tử hóa mô hình Qwen-72B từ FP16 xuống INT4 bằng phương pháp GPTQ. Sau khi nén, chỉ số Perplexity trên tập tin tức tiếng Anh chỉ tăng nhẹ $0.2\%$, nhưng tỷ lệ giải quyết bài toán lập trình trên tập HumanEval (Pass@1) giảm nghiêm trọng từ $75\%$ xuống $31\%$. Phân tích nguyên nhân gốc rễ và đề xuất 3 giải pháp khắc phục chuẩn công nghiệp.

  \item \textbf{Tình huống sự cố 2 (Hiện tượng sụp đổ phân phối xác suất trong DPO)}: \\
  Trong quá trình huấn luyện DPO cho mô hình 8B với tốc độ học $\eta = 5 \times 10^{-6}$ và $\beta = 0.5$, sau 500 bước huấn luyện, độ dài trung bình của câu trả lời sinh ra bất ngờ tăng vọt từ 200 tokens lên hơn 1800 tokens, chứa đầy các cụm từ lặp lại vô nghĩa. Phân tích nguyên nhân của hội chứng này và phương pháp phòng ngừa.

  \item \textbf{Tình huống sự cố 3 (Sụp đổ số học do ma trận Hessian suy biến trong GPTQ)}: \\
  Khi thực hiện lượng tử hóa một tầng MLP lớn của mô hình 70B bằng GPTQ không có hệ số điều hòa ($\lambda = 0$), hàm nghịch đảo ma trận \code{torch.linalg.cholesky} bất ngờ ném ngoại lệ \code{LinAlgError: The factorization could not be completed because the matrix is not positive-definite}. Sau khi bỏ qua lỗi và dùng giả nghịch đảo Moore-Penrose, các trọng số sau lượng tử hóa bùng nổ lên giá trị vô cực ($NaN / \pm\infty$). Phân tích nguyên nhân suy biến của ma trận kích hoạt và nguyên lý khắc phục bằng Damping Factor.
\end{enumerate}

---

\section{Hướng Dẫn Lời Giải \& Bình Luận Kỹ Thuật Chi Tiết}
\label{sec:ch11-solutions}

\subsection{Lời Giải Nhóm A: Câu Hỏi Tư Duy Bản Chất}
\begin{itemize}
  \item \textbf{Lời giải Câu A.1}:
  Nếu ma trận $\mathbf{B}$ được khởi tạo ngẫu nhiên, tích $\mathbf{A}\mathbf{B} \neq \mathbf{0}$, khiến $\mathbf{W} = \mathbf{W}_0 + \frac{\alpha}{r}\mathbf{A}\mathbf{B}$ bị sai lệch ngay tại bước huấn luyện $0$. Mô hình sẽ lập tức mất đi tri thức ngôn ngữ quý giá đã được tích lũy qua hàng tháng tiền huấn luyện, dẫn đến hiện tượng quên thảm họa (Catastrophic Forgetting). Bằng cách khởi tạo $\mathbf{B} = \mathbf{0}$, ta bảo toàn chính xác $100\%$ hành vi ban đầu của mô hình.

  \item \textbf{Lời giải Câu A.2}:
  Mạng nơ-ron học sâu là các hệ thống tối ưu hóa mang tính thống kê có độ dung sai cao với nhiễu số học nhỏ (noise-tolerant). Độ chính xác tuyệt đối ở chữ số thập phân thứ 4 hay thứ 5 (được cung cấp bởi 10 bit mantissa của FP16) hầu như không ảnh hưởng đến hướng đi tổng thể của véc-tơ gradient. Ngược lại, dải động (Dynamic Range) là yếu tố sống còn: một khi gradient bị underflow về 0, trọng số sẽ vĩnh viễn ngừng cập nhật. Do đó, đổi 3 bits mantissa để lấy thêm 3 bits exponent của BF16 là một cuộc đánh đổi hoàn hảo cho huấn luyện mô hình lớn.

  \item \textbf{Lời giải Câu A.3}:
  Phương pháp Round-To-Nearest (RTN) xem mọi tham số có tầm quan trọng bình đẳng. Trong các bài toán viết code và suy luận logic, các phép tính phụ thuộc sâu sắc vào các kênh kích hoạt ngoại lai (Activation Outliers). RTN làm tròn thô bạo các trọng số nối với các kênh này, làm lệch hướng biểu diễn ngữ nghĩa. AWQ nhận diện và phóng đại chính xác $1\%$ trọng số nhạy cảm nhất này lên trước khi làm tròn, bảo tồn trọn vẹn mạch tư duy logic của mô hình.

  \item \textbf{Lời giải Câu A.4}:
  Khi $\beta \to 0$, hình phạt phân kỳ KL bị vô hiệu hóa hoàn toàn, mô hình chính sách $\pi_\theta$ chỉ tập trung đẩy xác suất của $y_w$ lên $1$ và dìm $y_l$ về $0$, nhanh chóng dẫn đến hiện tượng quá khớp và phá hủy khả năng tạo sinh tự nhiên. Khi $\beta \to \infty$, trọng số phạt KL áp đảo hàm mục tiêu, buộc $\pi_\theta$ phải giữ nguyên trạng thái giống hệt mô hình tham chiếu $\pi_{\text{ref}}$, khiến mô hình hoàn toàn không học được bất kỳ sở thích mới nào từ con người.

  \item \textbf{Lời giải Câu A.5}:
  DPO có ưu điểm là nền tảng giải tích chặt chẽ và khả năng kiểm soát độ trôi dạt chính sách chính xác qua mô hình tham chiếu $\pi_{\text{ref}}$, nhưng nhược điểm là tốn gấp đôi VRAM vì phải nạp cả hai mô hình lớn song song. ORPO kết hợp SFT và căn chỉnh trong một pha duy nhất, không cần mô hình tham chiếu, giúp tiết kiệm $50\%$ VRAM và rút ngắn thời gian huấn luyện. Do đó, ORPO là lựa chọn tối ưu khi nguồn tài nguyên GPU bị giới hạn hoặc khi cần tinh chỉnh nhanh trên phần cứng biên.

  \item \textbf{Lời giải Câu A.6}:
  Theo công thức OBS, $\Delta \mathcal{E}_q = \frac{1}{2} \frac{(w_q - \hat{w}_q)^2}{[\mathbf{H}^{-1}]_{qq}}$. Về mặt hình học, hàm mất mát $\mathcal{E}$ là một paraboloid elip nhiều chiều trong không gian tham số.
  \begin{itemize}
    \item Ma trận Hessian $\mathbf{H}$ xác định độ cong: Trị riêng lớn ứng với hướng có độ cong dốc đứng, trị riêng nhỏ ứng với hướng có độ cong thoai thoải.
    \item Tuy nhiên, khi ta lượng tử hóa $w_q$, ta không chỉ dịch chuyển trên trục $q$ mà còn được phép dịch chuyển trên tất cả các trục còn lại $\Delta w_{j \neq q}$ để tìm điểm có sai số cực tiểu trên siêu phẳng cắt $\Delta w_q = \hat{w}_q - w_q$.
    \item Theo bổ đề ma trận khối, phần tử đường chéo $[\mathbf{H}^{-1}]_{qq}$ tỉ lệ nghịch với khoảng cách từ mặt phẳng cắt đến mặt elip. Cụ thể: $[\mathbf{H}^{-1}]_{qq} = \frac{1}{H_{qq} - \mathbf{h}_q^T \mathbf{H}_{-q, -q}^{-1} \mathbf{h}_q}$. Đại lượng trong mẫu số chính là ``độ cong hiệu dụng'' (Effective Curvature) của trục $q$ sau khi đã chiếu tối ưu lên không gian con của các trọng số còn lại.
    \item Khi $[\mathbf{H}^{-1}]_{qq}$ rất lớn, độ cong hiệu dụng rất nhỏ, nghĩa là các trọng số còn lại có khả năng hấp thụ và bù trừ hầu như toàn bộ sai số của $w_q$, dẫn đến mức tăng sai số toàn cục $\Delta \mathcal{E}_q$ rất nhỏ!
  \end{itemize}

  \item \textbf{Lời giải Câu A.7}:
  So sánh ba thuật toán nén trọng số:
  \begin{itemize}
    \item \textbf{Round-To-Nearest (RTN)}: Độ phức tạp tính toán $\mathcal{O}(d_{\text{in}} d_{\text{out}})$, không cần dữ liệu hiệu chuẩn hay bộ nhớ phụ. Tuy nhiên, sai số làm tròn tích lũy độc lập khiến mô hình 4-bit bị sụp đổ hoàn toàn về năng lực suy luận.
    \item \textbf{Optimal Brain Surgeon (OBS) Cổ Điển}: Mỗi lần lượng tử hóa 1 trọng số, phải cập nhật ma trận Hessian nghịch đảo theo công thức Sherman-Morrison tốn $\mathcal{O}(d_{\text{in}}^2)$, dẫn đến tổng thời gian cho một tầng là $\mathcal{O}(d_{\text{in}}^3 \cdot d_{\text{out}})$ hoặc $\mathcal{O}(d_{\text{in}}^4)$. Với mô hình Llama-3-70B ($d = 8192$), con số này đòi hỏi hàng triệu tỷ phép tính tuần tự không thể song song hóa, khiến OBS cổ điển bất khả thi trên thực tế.
    \item \textbf{Cholesky GPTQ}: Bằng cách cố định thứ tự nén và sử dụng phân rã Cholesky $\mathbf{H}^{-1} = \mathbf{L} \mathbf{L}^T$, ma trận nghịch đảo chỉ cần tính một lần duy nhất ($\mathcal{O}(d_{\text{in}}^3)$). Hơn thế nữa, bằng kỹ thuật Lazy Block Updates (gom $B = 128$ cột), toàn bộ các thao tác bù sai số được thực hiện qua các phép nhân ma trận GEMM trên GPU Tensor Cores, nén trọn vẹn mô hình 70B chỉ trong $2 - 4$ giờ!
  \end{itemize}

  \item \textbf{Lời giải Câu A.8}:
  Trong thuật toán PPO truyền thống, mạng Critic $V_\phi(s)$ được dùng để dự đoán kỳ vọng điểm thưởng nền (Baseline) nhằm trừ bớt khỏi điểm thưởng thực tế, thu được độ lợi thế $A = R - V(s)$ giúp giảm thiểu phương sai của gradient. Tuy nhiên, việc huấn luyện một mạng nơ-ron Critic lớn gây tốn gấp đôi VRAM và cực kỳ bất ổn định khi chuỗi suy luận kéo dài hàng chục nghìn tokens.
  GRPO triệt tiêu hoàn toàn mạng Critic bằng cách lấy mẫu một nhóm $G$ câu trả lời độc lập cho cùng một câu hỏi và dùng giá trị trung bình mẫu của nhóm $\bar{r} = \frac{1}{G}\sum r_i$ làm đường cơ sở. Về mặt lý thuyết xác suất, $\bar{r}$ chính là ước lượng Monte Carlo không chệch của giá trị kỳ vọng trạng thái $\mathbb{E}_{o \sim \pi}[r(q, o)] = V^\pi(q)$ mà hoàn toàn không tốn bất kỳ tham số hay trạng thái bộ tối ưu nào! Hơn nữa, việc chia cho độ lệch chuẩn nhóm $\sigma_r$ chuẩn hóa độ lớn gradient giữa các câu hỏi dễ và khó, giữ cho quá trình hội tụ cực kỳ ổn định.

  \item \textbf{Lời giải Câu A.9}:
  Trong mô hình nhị phân 1-bit thuần túy ($W \in \{-1, +1\}$), mọi liên kết nơ-ron bắt buộc phải mang tính kích hoạt (+1) hoặc ức chế (-1). Khi xử lý ngôn ngữ tự nhiên, một token cụ thể chỉ liên quan đến một tập hợp nhỏ các đặc trưng ngữ nghĩa. Việc ép buộc các liên kết không liên quan phải nhận giá trị $\pm 1$ sẽ bơm một lượng nhiễu khổng lồ vào các tầng sâu, phá hủy hoàn toàn độ chính xác của mô hình.
  Việc bổ sung trạng thái thứ ba ($0$) trong BitNet b1.58 mang lại hai năng lực then chốt:
  \begin{itemize}
    \item \textbf{Lọc tính năng tường minh (Explicit Feature Filtering)}: Mô hình có khả năng chủ động ngắt hoàn toàn các kết nối không quan trọng ($W_{ij} = 0$).
    \item \textbf{Tính thưa động (Dynamic Sparsity)}: Chỉ những kênh kích hoạt thực sự quan trọng mới được tích lũy vào nơ-ron đích, phản ánh chính xác cấu trúc thưa của mạng nơ-ron sinh học và bảo toàn trọn vẹn năng lực suy luận tương đương FP16.
  \end{itemize}

  \item \textbf{Lời giải Câu A.10}:
  Trong quá trình suy luận theo lô hoặc sinh từng token (Auto-regressive Generation), tầng giải mã Transformer bị giới hạn bởi băng thông bộ nhớ (Memory Bandwidth Bound). Nếu sau mỗi tầng chuẩn hóa, hệ thống phải thực hiện thêm một phép chia véc-tơ $\mathbf{x} \oslash \mathbf{s}$ trên kích hoạt trước khi đưa vào nhân GEMM, GPU bắt buộc phải thực thi thêm một kernel tính toán riêng biệt (Elementwise Kernel), gây tiêu tốn băng thông đọc/ghi bộ nhớ VRAM và tạo ra rào cản đồng bộ luồng (Kernel Launch Overhead).
  Kỹ thuật LayerNorm Folding nhận diện rằng phép chuẩn hóa $\text{RMSNorm}(\mathbf{z}) = \frac{\mathbf{z}}{\text{RMS}(\mathbf{z})} \odot \boldsymbol{\gamma}_{\text{LN}}$ kết thúc bằng một phép nhân vô hướng với vector trọng số tĩnh $\boldsymbol{\gamma}_{\text{LN}}$. Do đó:
  \begin{equation*}
  \frac{\mathbf{z}}{\text{RMS}(\mathbf{z})} \odot \boldsymbol{\gamma}_{\text{LN}} \oslash \mathbf{s} = \frac{\mathbf{z}}{\text{RMS}(\mathbf{z})} \odot \left( \boldsymbol{\gamma}_{\text{LN}} \oslash \mathbf{s} \right) = \frac{\mathbf{z}}{\text{RMS}(\mathbf{z})} \odot \boldsymbol{\gamma}_{\text{LN}}'
  \end{equation*}
  Bằng cách chia trước ma trận trọng số tĩnh $\boldsymbol{\gamma}_{\text{LN}}' = \boldsymbol{\gamma}_{\text{LN}} \oslash \mathbf{s}$ một lần duy nhất khi nạp mô hình vào bộ nhớ, đầu ra của tầng LayerNorm tự động là $\mathbf{x}'$ mà hoàn toàn không tốn thêm bất kỳ chỉ lệnh máy nào trong quá trình suy luận, triệt tiêu $100\%$ độ trễ thời gian chạy ($0\text{ ns}$).

  \item \textbf{Lời giải Câu A.11}:
  Trong LoRA truyền thống, $\mathbf{W} = \mathbf{W}_0 + \frac{\alpha}{r}\mathbf{A}\mathbf{B}$. Bất kỳ bước cập nhật gradient nào lên $\mathbf{A}, \mathbf{B}$ cũng đồng thời làm thay đổi độ lớn (Norm) và góc xoay (Angle) của từng cột trọng số theo một tỷ lệ phụ thuộc chặt chẽ vào nhau (Coupled Dynamics). Điều này hạn chế khả năng tinh chỉnh tinh tế của mô hình: nó không thể xoay góc biểu diễn mà vẫn giữ nguyên độ lớn chuẩn của nơ-ron.
  DoRA giải quyết vấn đề này bằng cách phân tách:
  \begin{equation*}
  \mathbf{W} = \mathbf{m} \odot \frac{\mathbf{V}}{\|\mathbf{V}\|_c}, \quad \text{với } \mathbf{V} = \mathbf{W}_0 + \frac{\alpha}{r}\mathbf{A}\mathbf{B}
  \end{equation*}
  Về mặt hình học, gradient cập nhật cho vector hướng $\mathbf{v}$ được nhân với ma trận chiếu trực giao $\mathbf{P}_{\mathbf{v}} = \mathbf{I} - \frac{\mathbf{v}\mathbf{v}^T}{\|\mathbf{v}\|_2^2}$. Vì $\mathbf{v}^T \mathbf{P}_{\mathbf{v}} = \mathbf{0}$, mọi gradient đẩy lên nhánh LoRA đều vuông góc tuyệt đối với vector hướng hiện tại. Nhờ đó, nhánh LoRA chỉ thuần túy trượt trên mặt cầu đơn vị để tìm góc biểu diễn tối ưu, trong khi độ lớn của nơ-ron được bàn giao hoàn toàn cho vector $\mathbf{m}$ điều chỉnh độc lập, tái lập chính xác động học của Full Fine-Tuning.
\end{itemize}

\subsection{Lời Giải Nhóm B: Bài Toán Tính Toán Giải Tích Bằng Tay}

\paragraph{Lời giải Bài B.1}:
\begin{itemize}
  \item[a)] Dung lượng VRAM lưu trọng số gốc Llama-3-70B dạng NF4 ($0.5\text{ byte/tham số}$):
  \begin{equation*}
  V_{\text{base}} = 70 \times 10^9 \times 0.5\text{ Bytes} = 35.0 \times 10^9\text{ Bytes} \approx 32.60\text{ GiB}
  \end{equation*}

  \item[b)] Dung lượng cho nhánh LoRA ($10^8$ tham số):
  \begin{itemize}
    \item Trọng số LoRA (BF16, 2 bytes): $10^8 \times 2\text{ B} = 0.2\text{ GB}$.
    \item Gradients LoRA (BF16, 2 bytes): $10^8 \times 2\text{ B} = 0.2\text{ GB}$.
    \item Trạng thái Optimizer AdamW (FP32, 8 bytes): $10^8 \times 8\text{ B} = 0.8\text{ GB}$.
  \end{itemize}
  Tổng dung lượng nhánh LoRA:
  \begin{equation*}
  V_{\text{LoRA}} = 0.2 + 0.2 + 0.8 = 1.2\text{ GB}
  \end{equation*}

  \item[c)] Tổng VRAM toàn cục:
  \begin{equation*}
  V_{\text{total}} = V_{\text{base}} + V_{\text{LoRA}} + V_{\text{activation}} \approx 32.60\text{ GiB} + 1.2\text{ GiB} + 4.0\text{ GiB} \approx 37.8\text{ GiB}
  \end{equation*}
  Vì $37.8\text{ GiB} < 48.0\text{ GiB}$ (tổng bộ nhớ của 2 card RTX 3090/4090 24GB), hệ thống \textbf{hoàn toàn có thể huấn luyện mượt mà} trên máy tính cá nhân cao cấp!
\end{itemize}

\paragraph{Lời giải Bài B.2}:
Sai số làm tròn $e = X - \hat{X}$ nằm trong đoạn $[-\Delta / 2, \Delta / 2]$ và tuân theo phân phối đều liên tục $e \sim \mathcal{U}[-\Delta/2, \Delta/2]$. Hàm mật độ xác suất:
\begin{equation*}
f(e) = \frac{1}{\frac{\Delta}{2} - \left(-\frac{\Delta}{2}\right)} = \frac{1}{\Delta}
\end{equation*}
Kỳ vọng sai số:
\begin{equation*}
E[e] = \int_{-\Delta/2}^{\Delta/2} e \cdot \frac{1}{\Delta} \, de = \left[ \frac{e^2}{2\Delta} \right]_{-\Delta/2}^{\Delta/2} = 0
\end{equation*}
Phương sai sai số:
\begin{equation*}
\text{Var}(e) = E[e^2] - (E[e])^2 = \int_{-\Delta/2}^{\Delta/2} e^2 \cdot \frac{1}{\Delta} \, de = \frac{1}{\Delta} \left[ \frac{e^3}{3} \right]_{-\Delta/2}^{\Delta/2} = \frac{1}{\Delta} \left( \frac{\Delta^3}{24} - \left(-\frac{\Delta^3}{24}\right) \right) = \frac{\Delta^2}{12}
\end{equation*}
(Điều phải chứng minh).

\paragraph{Lời giải Bài B.3}:
Đặt $u = \beta \log \frac{\pi_\theta(y_w \mid x)}{\pi_{\text{ref}}(y_w \mid x)} - \beta \log \frac{\pi_\theta(y_l \mid x)}{\pi_{\text{ref}}(y_l \mid x)} = \hat{r}_\theta(x, y_w) - \hat{r}_\theta(x, y_l)$.
Hàm mất mát là $\mathcal{L}_{\text{DPO}} = -\log \sigma(u)$. Áp dụng quy tắc đạo hàm hàm hợp:
\begin{equation*}
\nabla_\theta \mathcal{L}_{\text{DPO}} = -\frac{1}{\sigma(u)} \cdot \sigma'(u) \cdot \nabla_\theta u
\end{equation*}
Sử dụng tính chất vi phân của hàm sigmoid: $\sigma'(u) = \sigma(u)(1 - \sigma(u)) = \sigma(u)\sigma(-u)$. Khi đó:
\begin{equation*}
\frac{\sigma'(u)}{\sigma(u)} = 1 - \sigma(u) = \sigma(-u) = \sigma\left(\hat{r}_\theta(x, y_l) - \hat{r}_\theta(x, y_w)\right)
\end{equation*}
Mặt khác, vì mô hình tham chiếu $\pi_{\text{ref}}$ là hằng số đối với $\theta$:
\begin{equation*}
\nabla_\theta u = \beta \nabla_\theta \left[\log \pi_\theta(y_w \mid x) - \log \pi_\theta(y_l \mid x)\right]
\end{equation*}
Kết hợp lại, ta thu được:
\begin{equation*}
\nabla_\theta \mathcal{L}_{\text{DPO}}(\theta) = -\beta \cdot \sigma\left( \hat{r}_\theta(x, y_l) - \hat{r}_\theta(x, y_w) \right) \cdot \left[ \nabla_\theta \log \pi_\theta(y_w \mid x) - \nabla_\theta \log \pi_\theta(y_l \mid x) \right]
\end{equation*}
(Điều phải chứng minh).

\paragraph{Lời giải Bài B.4}:
\begin{enumerate}
  \item Tính ma trận Hessian $\mathbf{H}$ và $\mathbf{H}^{-1}$:
  \begin{equation*}
  \mathbf{X} \mathbf{X}^T = \begin{bmatrix} 1 & 2 & 1 \\ 0 & 1 & 2 \end{bmatrix} \begin{bmatrix} 1 & 0 \\ 2 & 1 \\ 1 & 2 \end{bmatrix} = \begin{bmatrix} 6 & 4 \\ 4 & 5 \end{bmatrix}
  \end{equation*}
  Ma trận Hessian:
  \begin{equation*}
  \mathbf{H} = 2 \mathbf{X} \mathbf{X}^T = \begin{bmatrix} 12 & 8 \\ 8 & 10 \end{bmatrix}
  \end{equation*}
  Định thức: $\det(\mathbf{H}) = 12 \times 10 - 8 \times 8 = 120 - 64 = 56$.
  Ma trận Hessian nghịch đảo:
  \begin{equation*}
  \mathbf{H}^{-1} = \frac{1}{56} \begin{bmatrix} 10 & -8 \\ -8 & 12 \end{bmatrix} = \begin{bmatrix} \frac{5}{28} & -\frac{1}{7} \\ -\frac{1}{7} & \frac{3}{14} \end{bmatrix} \approx \begin{bmatrix} 0.17857 & -0.14286 \\ -0.14286 & 0.21429 \end{bmatrix}
  \end{equation*}

  \item Tính vector bù sai số $\Delta \mathbf{w}$ khi lượng tử hóa $w_1 = 1.5 \to \hat{w}_1 = 2.0$:
  Sai số cục bộ: $w_1 - \hat{w}_1 = 1.5 - 2.0 = -0.5$.
  Phần tử đường chéo $[\mathbf{H}^{-1}]_{11} = \frac{5}{28}$.
  Áp dụng công thức cập nhật bù trừ OBS:
  \begin{equation*}
  \Delta \mathbf{w} = - \frac{w_1 - \hat{w}_1}{[\mathbf{H}^{-1}]_{11}} \cdot \mathbf{H}^{-1}_{:, 1} = - \frac{-0.5}{\frac{5}{28}} \begin{bmatrix} \frac{5}{28} \\ -\frac{1}{7} \end{bmatrix} = 2.8 \begin{bmatrix} \frac{5}{28} \\ -\frac{4}{28} \end{bmatrix} = \begin{bmatrix} 0.5 \\ -0.4 \end{bmatrix}
  \end{equation*}

  \item Giá trị mới của $w_2'$:
  \begin{equation*}
  w_2' = w_2 + \Delta w_2 = -2.0 + (-0.4) = \mathbf{-2.4}
  \end{equation*}
  Đầu ra ban đầu: $\mathbf{y} = [1.5, 1.0, -2.5]$.
  - Sai số khi dùng RTN ($\widehat{\mathbf{w}} = [2.0, -2.0]$):
  $\widehat{\mathbf{y}}_{\text{RTN}} = [2.0, 2.0, -2.0] \implies \|\mathbf{y} - \widehat{\mathbf{y}}_{\text{RTN}}\|_2^2 = (1.5-2.0)^2 + (1.0-2.0)^2 + (-2.5 - (-2.0))^2 = 0.25 + 1.0 + 0.25 = \mathbf{1.50}$.
  - Sai số khi dùng OBS ($\mathbf{w}_{\text{OBS}} = [2.0, -2.4]$):
  $\widehat{\mathbf{y}}_{\text{OBS}} = [2.0, 1.6, -2.8] \implies \|\mathbf{y} - \widehat{\mathbf{y}}_{\text{OBS}}\|_2^2 = (1.5-2.0)^2 + (1.0-1.6)^2 + (-2.5 - (-2.8))^2 = 0.25 + 0.36 + 0.09 = \mathbf{0.70}$.
  OBS giảm thiểu hơn $53.3\%$ sai số bình phương tái tạo kích hoạt!
\end{enumerate}

\paragraph{Lời giải Bài B.5}:
Hàm mục tiêu: $\min_{\Delta \mathbf{w}} \frac{1}{2} \Delta \mathbf{w}^T \mathbf{H} \Delta \mathbf{w}$ dưới ràng buộc $\mathbf{e}_q^T \Delta \mathbf{w} + (w_q - \hat{w}_q) = 0$.
Thiết lập hàm Lagrange:
\begin{equation*}
\mathcal{L}(\Delta \mathbf{w}, \lambda) = \frac{1}{2} \Delta \mathbf{w}^T \mathbf{H} \Delta \mathbf{w} + \lambda \left(\mathbf{e}_q^T \Delta \mathbf{w} + w_q - \hat{w}_q\right)
\end{equation*}
Lấy đạo hàm riêng theo vector $\Delta \mathbf{w}$:
\begin{equation*}
\nabla_{\Delta \mathbf{w}} \mathcal{L} = \mathbf{H} \Delta \mathbf{w} + \lambda \mathbf{e}_q = \mathbf{0} \implies \mathbf{H} \Delta \mathbf{w} = -\lambda \mathbf{e}_q \implies \Delta \mathbf{w} = -\lambda \mathbf{H}^{-1} \mathbf{e}_q
\end{equation*}
Nhân vô hướng cả hai vế với $\mathbf{e}_q^T$:
\begin{equation*}
\mathbf{e}_q^T \Delta \mathbf{w} = -\lambda \mathbf{e}_q^T \mathbf{H}^{-1} \mathbf{e}_q = -\lambda [\mathbf{H}^{-1}]_{qq}
\end{equation*}
Theo điều kiện ràng buộc: $\mathbf{e}_q^T \Delta \mathbf{w} = -(w_q - \hat{w}_q) = \hat{w}_q - w_q$. Suy ra:
\begin{equation*}
\hat{w}_q - w_q = -\lambda [\mathbf{H}^{-1}]_{qq} \implies \lambda = \frac{w_q - \hat{w}_q}{[\mathbf{H}^{-1}]_{qq}}
\end{equation*}
Thế $\lambda$ vào công thức $\Delta \mathbf{w}$:
\begin{equation*}
\Delta \mathbf{w} = - \frac{w_q - \hat{w}_q}{[\mathbf{H}^{-1}]_{qq}} \cdot \mathbf{H}^{-1} \mathbf{e}_q = - \frac{w_q - \hat{w}_q}{[\mathbf{H}^{-1}]_{qq}} \cdot \mathbf{H}^{-1}_{:, q}
\end{equation*}
Tính mức gia tăng sai số toàn cục tối thiểu:
\begin{align*}
\Delta \mathcal{E}_q &= \frac{1}{2} \Delta \mathbf{w}^T \mathbf{H} \Delta \mathbf{w} = \frac{1}{2} \left(-\lambda \mathbf{e}_q^T \mathbf{H}^{-1}\right) \mathbf{H} \left(-\lambda \mathbf{H}^{-1} \mathbf{e}_q\right) \\
&= \frac{1}{2} \lambda^2 \mathbf{e}_q^T \mathbf{H}^{-1} \mathbf{e}_q = \frac{1}{2} \left(\frac{w_q - \hat{w}_q}{[\mathbf{H}^{-1}]_{qq}}\right)^2 [\mathbf{H}^{-1}]_{qq} = \frac{1}{2} \frac{(w_q - \hat{w}_q)^2}{[\mathbf{H}^{-1}]_{qq}}
\end{align*}
(Điều phải chứng minh).

\paragraph{Lời giải Bài B.6}:
\begin{enumerate}
  \item Chứng minh tổng độ lợi thế nhóm luôn triệt tiêu về $0$:
  Theo định nghĩa, $\bar{r} = \frac{1}{G}\sum_{i=1}^G r_i \iff \sum_{i=1}^G r_i = G \bar{r}$.
  Lấy tổng của tất cả các độ lợi thế $\widetilde{A}_i$:
  \begin{equation*}
  \sum_{i=1}^G \widetilde{A}_i = \sum_{i=1}^G \frac{r_i - \bar{r}}{\sigma_r + \epsilon} = \frac{1}{\sigma_r + \epsilon} \sum_{i=1}^G (r_i - \bar{r}) = \frac{1}{\sigma_r + \epsilon} \left( \sum_{i=1}^G r_i - \sum_{i=1}^G \bar{r} \right) = \frac{1}{\sigma_r + \epsilon} \left( G \bar{r} - G \bar{r} \right) = 0
  \end{equation*}
  Do đó, kỳ vọng mẫu của độ lợi thế nhóm:
  \begin{equation*}
  \mathbb{E}_{\text{group}}[\widetilde{A}] = \frac{1}{G} \sum_{i=1}^G \widetilde{A}_i = \frac{0}{G} = 0
  \end{equation*}
  (Điều phải chứng minh).

  \item Khi phương sai điểm thưởng nhóm triệt tiêu ($\sigma_r \to 0$):
  Tất cả các câu trả lời sinh ra trong nhóm nhận điểm thưởng bằng hệt nhau (ví dụ: $r_i = 1.0$ cho mọi $i$ trên câu hỏi quá dễ, hoặc $r_i = 0.0$ cho mọi $i$ trên câu hỏi quá khó).
  Trong trường hợp này, dữ liệu không cung cấp bất kỳ tín hiệu tương đối nào để mô hình biết hành động nào tốt hơn hành động nào.
  Bằng cách gán $\widetilde{A}_i = 0$, gradient của hàm mục tiêu chính sách triệt tiêu:
  \begin{equation*}
  \nabla_\theta \mathcal{J}_{\text{GRPO}} = \mathbf{0}
  \end{equation*}
  Đây là một cơ chế tự bảo vệ tuyệt đẹp của học tăng cường: Nó ngăn mô hình không cập nhật mù quáng khi không có tín hiệu đối chứng, tiết kiệm năng lực tính toán và giữ ổn định tham số mạng!
\end{enumerate}

\paragraph{Lời giải Bài B.7}:
\begin{enumerate}
  \item Dung lượng thông tin Entropi theo Shannon của biến rời rạc 3 trạng thái đẳng xác suất:
  \begin{equation*}
  H(X) = -\sum_{i=1}^3 P(x_i) \log_2 P(x_i) = -3 \times \left(\frac{1}{3} \log_2 \frac{1}{3}\right) = \log_2(3) = \frac{\ln 3}{\ln 2} \approx \frac{1.098612}{0.693147} \approx \mathbf{1.5849625\text{ bits}}
  \end{equation*}
  (Điều phải chứng minh).

  \item Đóng gói 5 tham số bậc ba vào 1 byte:
  Tổng số trạng thái tổ hợp của 5 tham số bậc ba là $3^5 = 243$.
  Một byte ($8\text{ bits}$) có dải biểu diễn $2^8 = 256$ trạng thái. Vì $243 \le 256$, toàn bộ 5 tham số có thể được mã hóa duy nhất vào 1 byte theo hệ cơ số 3:
  \begin{equation*}
  B = w_0 + 3 w_1 + 3^2 w_2 + 3^3 w_3 + 3^4 w_4, \quad \text{với } w_k \in \{0, 1, 2\}
  \end{equation*}
  Đối với mô hình 70 tỷ tham số ($70 \times 10^9$ trọng số):
  \begin{equation*}
  \text{Dung lượng lưu trữ} = \frac{70 \times 10^9 \text{ tham số}}{5 \text{ tham số/byte}} = 14 \times 10^9\text{ Bytes}
  \end{equation*}
  Quy đổi sang GiB ($1\text{ GiB} = 1024^3\text{ Bytes} = 1{,}073{,}741{,}824\text{ Bytes}$):
  \begin{equation*}
  V_{\text{BitNet}} = \frac{14 \times 10^9}{1024^3} \approx \mathbf{13.0385\text{ GiB}}
  \end{equation*}
  So sánh với mô hình FP16 gốc ($2\text{ bytes/tham số}$):
  \begin{equation*}
  V_{\text{FP16}} = \frac{70 \times 10^9 \times 2}{1024^3} \approx \mathbf{130.385\text{ GiB}}
  \end{equation*}
  Mô hình BitNet b1.58 cắt giảm chính xác \textbf{$10$ lần ($90\%$)} dung lượng bộ nhớ, giúp một mô hình 70 tỷ tham số chạy mượt mà ngay trên RAM $16\text{ GB}$ của máy tính cá nhân hoặc thiết bị biên thông minh!
\end{enumerate}

\paragraph{Lời giải Bài B.8}:
\begin{enumerate}
  \item Dung lượng scale đơn (FP32, 32 bits) trên từng tham số với block size $B_1 = 64$:
  \begin{equation*}
  b_{\text{single}} = \frac{32\text{ bits}}{64} = 0.50\text{ bits/tham số}
  \end{equation*}

  \item Khi áp dụng Double Quantization với $B_2 = 256$, scale cấp 1 được lượng tử hóa sang FP8 ($8\text{ bits}$), và scale cấp 2 giữ FP32 ($32\text{ bits}$):
  \begin{align*}
  b_{\text{double}} &= \frac{8\text{ bits}}{B_1} + \frac{32\text{ bits}}{B_1 \times B_2} = \frac{8}{64} + \frac{32}{64 \times 256} \\
  &= 0.125 + \frac{32}{16384} \approx 0.125 + 0.001953 = 0.126953\text{ bits/tham số}
  \end{align*}
  Mức tiết kiệm bộ nhớ trên từng tham số:
  \begin{equation*}
  \Delta b = b_{\text{single}} - b_{\text{double}} = 0.50 - 0.126953 = \mathbf{0.373047\text{ bits/tham số}} \approx \mathbf{0.373\text{ bits/tham số}}
  \end{equation*}
  (Điều phải chứng minh).

  \item Tính dung lượng VRAM tiết kiệm được cho mô hình 70B và 405B:
  \begin{itemize}
    \item Với Llama-3-70B ($N_1 = 70 \times 10^9$ tham số):
    \begin{equation*}
    \Delta M_{\text{70B}} = \frac{70 \times 10^9 \times 0.373047}{8 \times 1024^3} = \frac{26.113 \times 10^9}{8{,}589{,}934{,}592} \approx \mathbf{3.04\text{ GiB}}
    \end{equation*}
    (Tương đương xấp xỉ $3.26\text{ GB}$ theo hệ thập phân).
    \item Với Llama-3.1-405B ($N_2 = 405 \times 10^9$ tham số):
    \begin{equation*}
    \Delta M_{\text{405B}} = \frac{405 \times 10^9 \times 0.373047}{8 \times 1024^3} = \frac{151.084 \times 10^9}{8{,}589{,}934{,}592} \approx \mathbf{17.59\text{ GiB}}
    \end{equation*}
    (Tương đương xấp xỉ $18.89\text{ GB}$ theo hệ thập phân).
    Mức tiết kiệm $17.59\text{ GiB}$ này tương đương với dung lượng của gần một card đồ họa RTX 4090 độc lập, cho phép mở rộng batch size hoặc chiều dài ngữ cảnh gấp nhiều lần!
  \end{itemize}
\end{enumerate}

\paragraph{Lời giải Bài B.9}:
\begin{enumerate}
  \item Đặt vector trọng số hiệu dụng $\mathbf{w} = m \frac{\mathbf{v}}{\|\mathbf{v}\|_2}$ với $m = \|\mathbf{w}\|_2 > 0$ và $\mathbf{v} \neq \mathbf{0}$.
  Ký hiệu $\|\mathbf{v}\|_2 = (\mathbf{v}^T \mathbf{v})^{1/2}$. Vi phân của chuẩn L2:
  \begin{equation*}
  d\|\mathbf{v}\|_2 = \frac{\mathbf{v}^T d\mathbf{v}}{\|\mathbf{v}\|_2}
  \end{equation*}
  Vi phân của vector $\mathbf{w}$ theo $\mathbf{v}$ (khi cố định $m$):
  \begin{equation*}
  d\mathbf{w} = m \left( \frac{d\mathbf{v}}{\|\mathbf{v}\|_2} - \frac{\mathbf{v} d\|\mathbf{v}\|_2}{\|\mathbf{v}\|_2^2} \right) = m \left( \frac{d\mathbf{v}}{\|\mathbf{v}\|_2} - \frac{\mathbf{v} (\mathbf{v}^T d\mathbf{v})}{\|\mathbf{v}\|_2^3} \right) = \frac{m}{\|\mathbf{v}\|_2} \left( \mathbf{I} - \frac{\mathbf{v}\mathbf{v}^T}{\|\mathbf{v}\|_2^2} \right) d\mathbf{v}
  \end{equation*}
  Theo quy tắc đạo hàm hàm hợp, cho gradient $\mathbf{g} = \nabla_{\mathbf{w}} \mathcal{L}$:
  \begin{equation*}
  d\mathcal{L} = \mathbf{g}^T d\mathbf{w} = \mathbf{g}^T \left[ \frac{m}{\|\mathbf{v}\|_2} \left( \mathbf{I} - \frac{\mathbf{v}\mathbf{v}^T}{\|\mathbf{v}\|_2^2} \right) d\mathbf{v} \right] = \left[ \frac{m}{\|\mathbf{v}\|_2} \left( \mathbf{I} - \frac{\mathbf{v}\mathbf{v}^T}{\|\mathbf{v}\|_2^2} \right) \mathbf{g} \right]^T d\mathbf{v}
  \end{equation*}
  Do đó:
  \begin{equation*}
  \nabla_{\mathbf{v}} \mathcal{L} = \frac{m}{\|\mathbf{v}\|_2} \left( \mathbf{I} - \frac{\mathbf{v}\mathbf{v}^T}{\|\mathbf{v}\|_2^2} \right) \mathbf{g}
  \end{equation*}
  (Điều phải chứng minh).

  \item Tính tích vô hướng $\mathbf{v}^T \nabla_{\mathbf{v}} \mathcal{L}$:
  \begin{align*}
  \mathbf{v}^T \nabla_{\mathbf{v}} \mathcal{L} &= \mathbf{v}^T \left[ \frac{m}{\|\mathbf{v}\|_2} \left( \mathbf{I} - \frac{\mathbf{v}\mathbf{v}^T}{\|\mathbf{v}\|_2^2} \right) \mathbf{g} \right] = \frac{m}{\|\mathbf{v}\|_2} \left( \mathbf{v}^T - \frac{(\mathbf{v}^T \mathbf{v}) \mathbf{v}^T}{\|\mathbf{v}\|_2^2} \right) \mathbf{g} \\
  &= \frac{m}{\|\mathbf{v}\|_2} \left( \mathbf{v}^T - \frac{\|\mathbf{v}\|_2^2 \mathbf{v}^T}{\|\mathbf{v}\|_2^2} \right) \mathbf{g} = \frac{m}{\|\mathbf{v}\|_2} \left( \mathbf{v}^T - \mathbf{v}^T \right) \mathbf{g} = \frac{m}{\|\mathbf{v}\|_2} \cdot \mathbf{0}^T \cdot \mathbf{g} = 0
  \end{align*}
  (Điều phải chứng minh).
  Hệ quả: Gradient cập nhật cho nhánh hướng $\mathbf{v}$ luôn trực giao hoàn toàn với vector $\mathbf{v}$ hiện tại, chứng minh DoRA phân rã triệt để giữa góc hướng và độ dài vector.
\end{enumerate}

\subsection{Lời Giải Nhóm C: Thực Hành Kỹ Nghệ \& Lập Trình}
Dưới đây là mã nguồn hoàn chỉnh cho các bài thực hành phòng lab:

\begin{lstlisting}[language=Python, caption={Cài đặt MultiLoRAManager và kiểm tra tính hợp lệ đa khách hàng.}]
from typing import Dict, Optional, Tuple
import numpy as np

class MultiLoRAManager:
    def __init__(self, d_in: int, d_out: int, r: int = 8, alpha: float = 16.0):
        self.d_in = d_in
        self.d_out = d_out
        self.r = r
        self.scaling = alpha / r
        # Mo hinh goc duy nhat (Shared Base Model)
        self.w0 = np.random.standard_normal((d_in, d_out)).astype(np.float32)
        # Tu dien luu adapter cua cac tenant: tenant_id -> (A, B)
        self.adapters: Dict[str, Tuple[np.ndarray, np.ndarray]] = {}

    def register_adapter(self, tenant_id: str, a_matrix: np.ndarray, b_matrix: np.ndarray):
        """Dang ky hoac nap adapter moi cho mot khach hang cu the."""
        assert a_matrix.shape == (self.d_in, self.r), "Kich thuoc ma tran A khong khop!"
        assert b_matrix.shape == (self.r, self.d_out), "Kich thuoc ma tran B khong khop!"
        self.adapters[tenant_id] = (a_matrix, b_matrix)

    def forward(self, x: np.ndarray, tenant_id: Optional[str] = None) -> np.ndarray:
        """Suy luan voi adapter cua tenant ma khong ton chi phi sao chep W0."""
        base_out = np.matmul(x, self.w0)
        if tenant_id is None or tenant_id not in self.adapters:
            return base_out

        a, b = self.adapters[tenant_id]
        lora_out = self.scaling * np.matmul(np.matmul(x, a), b)
        return base_out + lora_out
\end{lstlisting}

\begin{lstlisting}[language=Python, caption={Cài đặt OBSQuantizerEngine nén ma trận trọng số với bù sai số bậc hai.}]
class OBSQuantizerEngine:
    """Bo luong tu hoa Optimal Brain Surgeon 1 tang voi bu sai so bac hai."""
    def __init__(self, damping: float = 0.01):
        self.damping = damping

    def quantize_layer(
        self,
        w: np.ndarray,
        x: np.ndarray,
        bits: int = 4
    ) -> Tuple[np.ndarray, float, float]:
        """
        w: ma tran trong so (d_out, d_in)
        x: ma tran kich hoat (d_in, N)
        Tra ve: (w_quantized, mse_obs, mse_rtn)
        """
        d_out, d_in = w.shape
        w_obs = w.copy().astype(np.float64)
        w_rtn = w.copy().astype(np.float64)

        # 1. Tinh ma tran Hessian H = 2 * X @ X.T
        h = 2.0 * np.matmul(x, x.T)
        diag_mean = float(np.mean(np.diag(h)))
        if diag_mean == 0.0:
            diag_mean = 1.0

        # Dieu hoa Damping factor chong suy bien
        h += self.damping * diag_mean * np.eye(d_in)

        # 2. Tinh nghich dao Hessian H^-1
        h_inv = np.linalg.inv(h)

        # 3. Tinh Scale luong tu hoa INT cho tung hang
        q_max = (2 ** (bits - 1)) - 1
        scales = np.max(np.abs(w), axis=1) / q_max
        scales = np.maximum(scales, 1e-8)

        # 4. Luong tu hoa doc lap RTN de so sanh
        for r in range(d_out):
            s = scales[r]
            q_rtn = np.clip(np.round(w_rtn[r, :] / s), -q_max, q_max)
            w_rtn[r, :] = q_rtn * s

        # 5. Luong tu hoa OBS voi bu sai so lien tuc tren tung cot
        for col in range(d_in):
            h_inv_qq = h_inv[col, col]
            h_inv_col = h_inv[:, col]

            for r in range(d_out):
                w_orig = w_obs[r, col]
                s = scales[r]
                q_val = np.clip(np.round(w_orig / s), -q_max, q_max)
                w_hat = q_val * s
                err = w_orig - w_hat

                # Cap nhat bu sai so cho cac cot con lai
                delta_w = - (err / h_inv_qq) * h_inv_col
                w_obs[r, :] += delta_w
                w_obs[r, col] = w_hat

        # 6. Danh gia sai so tai tao kich hoa Y = W @ X
        y_true = np.matmul(w, x)
        y_obs = np.matmul(w_obs, x)
        y_rtn = np.matmul(w_rtn, x)

        mse_obs = float(np.mean((y_true - y_obs) ** 2))
        mse_rtn = float(np.mean((y_true - y_rtn) ** 2))

        return w_obs.astype(np.float32), mse_obs, mse_rtn
\end{lstlisting}

\paragraph{Lời giải Bài C.5}:
Xem chi tiết mã nguồn hoàn chỉnh của lớp \code{GRPOAdvantageEstimator} tại Mục~\ref{sec:python-production-modules}.
Quy trình hiện thực hóa gồm 3 thành phần chính:
\begin{enumerate}
  \item \textbf{Bộ kiểm tra quy tắc tất định (\code{rule_based_verifier})}: Dùng biểu thức chính quy Regex để nhận diện cấu trúc \verb|<think>...</think>| và trích xuất đáp số bên trong \verb|\boxed{...}|. Nếu đúng cả cấu trúc và đáp số, câu trả lời nhận tối đa điểm thưởng $r = 1.0 + \lambda_{\text{format}} \times 1.0$.
  \item \textbf{Bộ chuẩn hóa lợi thế nhóm (\code{compute_group_advantages})}: Tính trung bình mẫu $\bar{r}$ và độ lệch chuẩn mẫu $\sigma_r$ trên nhóm $G$ rollouts. Trả về $\widetilde{A}_i = \frac{r_i - \bar{r}}{\sigma_r + \epsilon}$. Nếu $\sigma_r < 10^{-8}$, trả về véc-tơ 0 để ngăn chặn cập nhật vô nghĩa.
  \item \textbf{Tổn thất Clipped Surrogate PPO (\code{compute_clipped_surrogate_loss})}: Tính tỷ số $\rho_t = \exp(\log \pi_\theta - \log \pi_{\text{old}})$, kẹp tỷ số trong $[1-\epsilon, 1+\epsilon]$ và lấy $\min$ giữa nhánh không kẹp và nhánh kẹp để tối ưu hóa chính sách an toàn.
\end{enumerate}

\paragraph{Lời giải Bài C.6}:
Xem chi tiết mã nguồn hoàn chỉnh của lớp \code{BitNetTernaryLinear} tại Mục~\ref{sec:python-production-modules}.
Phương thức \code{forward_pure_addition} hiện thực hóa cơ chế tính toán không dùng bộ nhân (Multiplier-less):
\begin{equation*}
y_{q, j} = \sum_{k: \widetilde{W}_{j,k} = +1} x_{q, k} - \sum_{k: \widetilde{W}_{j,k} = -1} x_{q, k}
\end{equation*}
Kết quả kiểm thử tự động trong phòng lab chứng minh rằng đầu ra của phương thức thuần cộng/trừ số nguyên trùng khớp chính xác $100\%$ (\code{assert_array_equal}) với phép nhân ma trận chuẩn $\mathbf{x}_q \mathbf{W}^T$, trong khi cắt giảm hơn $36\times$ năng lượng tiêu thụ trên phần cứng vi mạch!

\paragraph{Lời giải Bài C.7, C.8, C.9}:
Xem chi tiết mã nguồn hoàn chỉnh của các lớp \code{AWQScaleOptimizer}, \code{DoRALinearLayer}, và \code{QLoRADoubleQuantCalculator} tại Mục~\ref{sec:python-production-modules}.
\begin{itemize}
  \item \textbf{Bài C.7 (\code{AWQScaleOptimizer})}:
  Thuật toán thực hiện tìm kiếm lưới trên 20 điểm mẫu của siêu tham số $\alpha \in [0, 1]$. Trên dữ liệu hiệu chuẩn có kênh ngoại lai lớn ($X_0$ lớn gấp $50\times$), thuật toán tự động nhận diện giá trị tối ưu $\alpha^* \approx 0.40 - 0.60$, kéo giảm hơn $10\times$ sai số bình phương tái tạo so với lượng tử hóa thông thường không có AWQ. Phương thức \code{fold_scale_into_layernorm} chia trước véc-tơ trọng số LayerNorm cho $\mathbf{s}$, triệt tiêu hoàn toàn chi phí trễ thời gian chạy ($0\text{ ns}$).
  \item \textbf{Bài C.8 (\code{DoRALinearLayer})}:
  Khởi tạo $\mathbf{m} = \|\mathbf{W}_0\|_c$ và $\mathbf{B} = \mathbf{0}$, bảo đảm tính tương đương toán học chính xác $100\%$ với mô hình cơ sở tại thời điểm $t=0$. Khi tinh chỉnh, vector độ lớn $\mathbf{m}$ và ma trận hướng $\mathbf{V}$ được cập nhật độc lập, giúp nhánh LoRA tự do xoay góc mà không bị phụ thuộc độ lớn. Phương thức \code{merge_weights} hợp nhất vĩnh viễn $\mathbf{W}_{\text{merged}} = \mathbf{m} \odot \frac{\mathbf{W}_0 + \frac{\alpha}{r}\mathbf{A}\mathbf{B}}{\|\mathbf{W}_0 + \frac{\alpha}{r}\mathbf{A}\mathbf{B}\|_c}$ để phục vụ suy luận không trễ.
  \item \textbf{Bài C.9 (\code{QLoRADoubleQuantCalculator})}:
  Định cỡ bộ nhớ chi tiết cho mô hình Llama-3-70B: Double Quantization nén các hệ số co giãn cấp 1 (FP32) xuống FP8 với block size 256, tiết kiệm chính xác $0.373\text{ bits/tham số}$ ($3.04\text{ GiB}$). Toàn bộ chi phí VRAM cho mô hình tĩnh, adapter LoRA, đạo hàm và trạng thái bộ tối ưu AdamW được khống chế ở mức $34.7\text{ GiB}$, cho phép tinh chỉnh mượt mà trên 2 card đồ họa RTX 3090/4090 ($48\text{ GB}$).
\end{itemize}

\subsection{Lời Giải Nhóm D: Phân Tích Thất Bại \& Debugging Thực Tế}

\paragraph{Lời giải Sự cố 1 (Năng lực giải thuật suy thoái sau khi GPTQ)}:
\begin{enumerate}
  \item \textbf{Nguyên nhân gốc rễ (Root Cause)}:
  \begin{itemize}
    \item \textit{Lệch phân phối dữ liệu hiệu chỉnh (Calibration Mismatch)}: Tập dữ liệu calibration dùng cho GPTQ thường lấy từ văn bản tiếng Anh tổng hợp (C4, WikiText). Tập này thiếu vắng hoàn toàn mã nguồn (Code) và các ký hiệu toán học đặc thù, khiến ma trận Hessian $\mathbf{H}$ chỉ tối ưu hóa các kênh ngôn ngữ chung và làm méo mó các kênh kích hoạt biểu diễn logic thuật toán.
    \item \textit{Lượng tử hóa cả các tầng nhạy cảm}: Các ma trận chiếu đầu ra $\mathbf{W}_o$ của khối Attention và ma trận $\mathbf{W}_{\text{down}}$ của MLP có độ nhạy cảm sai số rất cao khi thực hiện suy luận chuỗi nhiều bước.
  \end{itemize}
  \item \textbf{3 Biện pháp khắc phục chuẩn mực}:
  \begin{itemize}
    \item Chuyển sang sử dụng AWQ với Group Size nhỏ ($G=64$ hoặc $G=32$) để bảo vệ các kênh ngoại lai quan trọng theo độ lớn thực tế.
    \item Trộn ít nhất $30\%$ dữ liệu mã nguồn Python/C++ và các bài toán GSM8K vào tập dữ liệu hiệu chỉnh.
    \item Áp dụng lượng tử hóa hỗn hợp (Mixed-Precision): Giữ nguyên $5\%$ các tầng nhạy cảm nhất ở định dạng FP16 hoặc INT8.
  \end{itemize}
\end{enumerate}

\paragraph{Lời giải Sự cố 2 (Hội chứng sụp đổ độ dài trong huấn luyện DPO)}:
\begin{enumerate}
  \item \textbf{Nguyên nhân gốc rễ}:
  Hiện tượng này được gọi là \textit{Hội chứng thiên vị độ dài (Length Exploitation / Verbosity Bias)}. Do hàm mất mát DPO tối ưu hóa tổng log-xác suất không được chuẩn hóa theo độ dài chuỗi:
  \begin{equation*}
  \log \pi_\theta(y \mid x) = \sum_{t=1}^{|y|} \log \pi_\theta(y_t \mid x, y_{<t})
  \end{equation*}
  Khi một câu trả lời $y_w$ có nhiều token hơn, nếu từng token có log-likelihood dương nhẹ so với mô hình tham chiếu, tổng tích lũy sẽ tăng tuyến tính theo số lượng token. Mô hình nhanh chóng phát hiện lỗ hổng này và sinh ra văn bản kéo dài vô tận để tối đa hóa điểm thưởng tiềm ẩn.
  \item \textbf{Giải pháp khắc phục}:
  \begin{itemize}
    \item Tăng hệ số phạt KL $\beta$ từ $0.1$ lên $0.3 - 0.5$ để hạn chế trôi dạt chính sách.
    \item Chuyển sang thuật toán SimPO với cơ chế chuẩn hóa độ dài bắt buộc $\frac{\beta}{|y|} \log \pi_\theta(y \mid x)$.
    \item Lọc và cân bằng độ dài của các cặp dữ liệu huấn luyện: loại bỏ các mẫu mà $|y_w| > 1.5 \times |y_l|$.
  \end{itemize}
\end{enumerate}

\paragraph{Lời giải Sự cố 3 (Sụp đổ số học do ma trận Hessian suy biến)}:
\begin{enumerate}
  \item \textbf{Nguyên nhân gốc rễ (Root Cause)}:
  Trong các mô hình Transformer sâu, hàm kích hoạt thường tạo ra hiện tượng nơ-ron chết (Dead Neurons, giá trị kích hoạt luôn bằng 0 qua toàn bộ token) hoặc hai nơ-ron có kích hoạt cộng tuyến (Collinear Activations, kênh $i$ bằng bội số của kênh $j$). Khi đó:
  \begin{equation*}
  \text{rank}(\mathbf{X} \mathbf{X}^T) < d_{\text{in}} \implies \det(\mathbf{H}) = 0
  \end{equation*}
  Ma trận Hessian không còn xác định dương, có trị riêng bằng 0 hoặc âm do sai số làm tròn số thực. Khi chạy thuật toán phân rã Cholesky $\mathbf{H}^{-1} = \mathbf{L} \mathbf{L}^T$, phép khai căn số âm $\sqrt{L_{ii}}$ sẽ thất bại ngay lập tức. Nếu dùng giả nghịch đảo không ổn định, các giá trị trên đường chéo $[\mathbf{H}^{-1}]_{qq} \to 0$, khiến phép chia trong công thức OBS $\frac{w_q - \hat{w}_q}{[\mathbf{H}^{-1}]_{qq}}$ bị nổ số mũ ra vô cực ($NaN / \pm\infty$).
  \item \textbf{Giải pháp kỹ nghệ chuẩn mực}:
  \begin{itemize}
    \item Bắt buộc thêm số hạng điều hòa Tikhonov (Damping Factor) $\lambda \mathbf{I}$ vào ma trận Hessian trước khi phân rã Cholesky:
    \begin{equation*}
    \widehat{\mathbf{H}} = 2 \mathbf{X} \mathbf{X}^T + \left( \alpha \cdot \frac{\text{Tr}(2 \mathbf{X} \mathbf{X}^T)}{d_{\text{in}}} \right) \mathbf{I}, \quad \text{với } \alpha = 0.01
    \end{equation*}
    \item Hệ số $\lambda > 0$ đảm bảo mọi trị riêng của $\widehat{\mathbf{H}}$ đều thỏa mãn $\lambda_i \ge \lambda > 0$, triệt tiêu hoàn toàn tính suy biến và bảo đảm thuật toán phân rã Cholesky luôn thành công $100\%$.
  \end{itemize}
\end{enumerate}
